\documentclass[11pt]{article}
\usepackage[T1]{fontenc}
\usepackage{microtype}
\usepackage{lmodern,amsmath,amssymb,amsthm,geometry,hyperref}
\geometry{margin=1in}
\hypersetup{hidelinks,pdftitle={Validated Dirichlet approximation and weighted arithmetic moments},pdfauthor={GRH research working note}}
\newtheorem{theorem}{Theorem}
\newtheorem{proposition}[theorem]{Proposition}
\title{Validated Dirichlet approximation and weighted arithmetic moments}
\author{GRH research working note}
\date{5 November 2026}
\begin{document}\maketitle
\begin{abstract}Adaptive outward dyadic arithmetic and theta-integral quadrature give a requested-error construction for three fixed completed Dirichlet functions. Fresh samples and executed rational readout training certify the supplied uniform $10^{-24}$ examples on the stated compact; no practical extreme-precision runtime bound is asserted. Separately, a fixed-twist sextic Hecke-family moment improves an original masked remainder estimate, while retaining a positive cutoff loss. We give the precise external analytic inputs and distinguish these results from zero-location and spectral determinant claims. A classical quadratic pivot improves one row sector; joint Gauss phases have an exact masked identity. An ordinary all-character conductor comparison improves the unsigned raw envelope to $H+DH^{1/6}+(DH)^{2/3}$, while an executed finite arithmetic Fourier readout certifies correlation and signed-observable errors without asserting a varying-family signed saving. A moving common-profile construction aligns the actual coefficients on nonprincipal sixth-power rows, and an exact joint example rules out variance-only centering; the fixed-profile signed target remains open. An exact conductor-12 character twist and gcd--M\"obius inversion give a separate bound for the signed coprime slice, retaining the unresolved $7/6$ power. A signed common-divisor zeta expansion and a complete-column exact laboratory retain every raw mask; weighted divisor certificates do not close the required asymptotic saving. No proof of GRH or complete Lean formalization is asserted.\end{abstract}
\section{A finite, validated approximation result}
Let $\chi_3=(0,1,-1)$, $\chi_4=(0,1,0,-1)$ and
$\chi_5=(0,1,-1,-1,1)$, extended periodically, with parities $a=1,1,0$.
We approximate the completed functions
\[
 \Lambda_q(s)=(q/\pi)^{(s+a)/2}\Gamma((s+a)/2)L(s,\chi_q).
\]
The primitive theta/Mellin identity is the explicit classical analytic input;
no hypothesis about zeros is used. These three examples do not assert a
uniform result for all conductors, general Hecke functions or Artin holomorphy.

\begin{theorem}[Scope of the supplied finite certificate]
For each $q=3,4,5$, the supplied exact-rational calculation constructs a
complex polynomial $P_q$ of degree $28$ satisfying
\[
 \sup_{|s-1/2|\le7/10}|P_q(s)-\Lambda_q(s)|<5\cdot10^{-15}.
\]
The computation includes enclosures of all $64$ true Fourier-circle nodes,
completed-function sample error, analytic truncation, Fourier aliasing,
the exact closed-form $180$-step gradient-descent iterate and coefficient
midpoint rounding. The Python implementation and informal analytic proof
have independent checks; this is not a kernel-checked Lean theorem.
\end{theorem}
\begin{proof}
Write $F(c+z)=\sum_{n\ge0}a_nz^n$, with $c=1/2$, $R=7/10$,
$\rho=6/5$, $r=12/5$. The independently recomputed rational theta majorants
are strictly below $M_q=4,10,6$, respectively, on $|z|\le r$.
Use true roots $\zeta_j=e^{2\pi i j/64}$ and rational labels $y_j$
with proved error $|y_j-F(c+\rho\zeta_j)|\le\eta$.
Use the real Euclidean gradient on complex coefficient coordinates for
\[
 \mathcal L(b)=\frac1{128}\sum_{j=0}^{63}
 \left|\sum_{n=0}^{28}b_n\zeta_j^n-y_j\right|^2.
\]
Discrete Fourier orthogonality gives the least-squares optimum
$b_n^*=64^{-1}\sum_jy_j\bar\zeta_j^n$, $0\le n\le28$,
and gradient $\nabla\mathcal L(b)=b-b^*$.
For step $1/4$ and zero initialization, the exact $180$-step iterate is
$b_n^{180}=(1-(3/4)^{180})b_n^*$. The code evaluates this closed form;
it does not claim an explicit180-update optimizer run.

Set $u=R/\rho=7/12$, $v=R/r=7/24$ and $w=(\rho/r)^{64}=2^{-64}$.
Let $\widehat b_n$ be an outward rectangle enclosing $b_n^*$ and
$\epsilon_n$ the sum of the coordinate halfwidths of the rectangle
for $b_n^{180}$. Choose each actual polynomial coefficient as that
rectangle's exact rational midpoint $c_n$. The normalized basis is explicit:
\[
 P_q(s)=\sum_{n=0}^{28}c_n\left(\frac{s-1/2}{6/5}\right)^n.
\]
Thus the exported coefficient with index $n$ multiplies the normalized
variable, rather than the raw power $(s-1/2)^n$. Its error on the compact disk is bounded by
\begin{align*}
 E_{\rm opt}&=(3/4)^{180}\sum_{n=0}^{28}|\widehat b_n|_1u^n,&
 E_{\rm arithmetic}&=\sum_{n=0}^{28}\epsilon_nu^n,\\
 E_{\rm oracle}&=\eta\frac{1-u^{29}}{1-u},&
 E_{\rm alias}&=M_q\frac{w}{1-w}\frac{1-v^{29}}{1-v},\\
 E_{\rm Taylor}&=M_q\frac{v^{29}}{1-v}.&&
\end{align*}
Here $|\widehat b_n|_1$ is the sum of the maximum absolute real and imaginary
coordinates, a safe upper bound for every enclosed complex modulus.
Cauchy's estimate $|a_n|\le M_qr^{-n}$ and the Fourier identity
$b_n^*=a_n\rho^n+\sum_{\ell\ge1}a_{n+64\ell}\rho^{n+64\ell}$
for exact labels prove the alias and Taylor bounds. Label perturbations
contribute at most $\eta$ per Fourier coefficient. The remaining two
terms follow directly from the exact iterate and coefficient rectangles.
Each term is rounded upward to a $180$-bit dyadic rational; the reported total
is the exact sum of those reported terms. The certified totals below are
strictly below the asserted threshold. Detailed sample enclosures are proved
in the extended edition and encoded with exact integer endpoints.
\end{proof}

\begin{center}\begin{tabular}{c|cc}Conductor & Maximum sample radius & Uniform polynomial bound\\\hline
3 & $3.026\times10^{-17}$ & $1.786\times10^{-15}$ \\
4 & $7.206\times10^{-17}$ & $4.455\times10^{-15}$ \\
5 & $1.404\times10^{-16}$ & $2.906\times10^{-15}$ \\
\end{tabular}\end{center}
The displayed values are rounded upward; exact rational records establish the strict inequalities.


\section{What this does and does not certify}
The final JSON records exact rational centers, interval endpoints, source
hashes and all five error terms. Float values are display summaries only.
A high-precision comparison is an independent diagnostic, not the reason
an interval contains its mathematical target. Arithmetic uses outward
integer operations, and invalid float/order inputs are rejected, including
through typed caches. Earlier pilots and failed validation probes are retained.

This is a fixed finite tolerance/configuration. Increasing accuracy needs
new truncation and precision choices; the fixed $N=11,T=128$ interface is
not an arbitrary-accuracy oracle. The broader approximation existence theorem
is separate. The certificate says nothing about the location of zeros:
$P_q$ is not required to be real-rooted or a self-adjoint spectral determinant.
The independent arithmetic determinant identity, positivity and original
signed short-row/covariance estimates remain open. No GRH proof or complete
new Lean4 formalization is claimed.

\section{Validated primitive operations}
Real intervals are closed dyadic intervals $[l2^{-224},h2^{-224}]$ with
integer endpoints; complex intervals are Cartesian products. Addition is
exact at this scale. For multiplication take the minimum/maximum of all four
integer endpoint products, divide by $2^{224}$ and round down/up. For a
reciprocal whose interval omits zero, reverse the endpoint order and round
$2^{448}/h$ down and $2^{448}/l$ up. Negative denominators and scalar division
use the same endpoint rules. Each operation encloses all represented real
values by monotonicity; induction proves inclusion for every expression
built from them. Rectangular complex multiplication uses these real rules.
The sum of coordinate halfwidths bounds the complex distance from its exact
rational center. Binary floating arguments are never silently converted.

For $0<x<1$, the alternating arctangent sum and the endpoint obtained by
adding the next signed term enclose $\arctan x$. Sixty terms at $1/5$ and
$1/239$, together with Machin's identity
$\pi=16\arctan(1/5)-4\arctan(1/239)$, give the interval for $\pi$.
The negative coefficient reverses the second interval. The identity follows
from tangent addition and the range of the two angles; no decimal value of
$\pi$ is a trusted input.

For rational $x>0$, range-reduce $x=2^kv$, $1\le v<2$. With
$z=(v-1)/(v+1)\le1/3$,
\[
 \log v=2\sum_{j=0}^{J-1}\frac{z^{2j+1}}{2j+1}+\mathcal R,
 \qquad 0\le\mathcal R\le
 \frac{2z^{2J+1}}{(2J+1)(1-z^2)}.
\]
The geometric tail bounds each remaining denominator below by $2J+1$.
Use $J=92$, the analogous expression for $\log2$, and sign-aware interval
multiplication by $k$. This covers negative $k$ as well.

For complex $z$, repeatedly halve until its interval $\ell^1$ magnitude
bound is $A\le1/2$. The degree $60$ Taylor polynomial for $e^z$ then has remainder
at most
\[
 \frac{A^{61}}{61!}\frac1{1-A/62}.
\]
Successive absolute terms have ratio at most $A/62$, so geometric comparison
proves this bound. Add its symmetric interval to each coordinate, then square
outward the required number of times. The identity $e^{2z}=(e^z)^2$ and
interval inclusion prove the final enclosure. Caches are type-sensitive:
otherwise a cached exact argument can accidentally bypass rejection of an
equal-valued float. No implementation change is allowed during the recorded run;
the producer verifies source bytes again before emitting its certificate.

\section{Recomputing the compact growth constants}
For the larger growth disk, $c=1/2$, $r=12/5$, take $k=2$ when $a=1$
and $k=1$ when $a=0$. Both real Mellin exponents on that disk are at most
$k$. Define
\[
 Q_k(x)=\sum_{j=0}^k\frac{k!}{(k-j)!x^{j+1}},\qquad
 U_{16}(x)=(1+x/16)^{-16},\qquad \delta=3/q.
\]
Integration by parts gives
$\int_1^\infty t^ke^{-xt}\,dt=e^{-x}Q_k(x)$.
The same Gaussian split and geometric-tail comparison used below, now
with cutoff $n_0=8$, imply the exact rational majorant
\[
 \widehat M_q=2\left[
 \sum_{n=1}^{7}n^aU_{16}(\delta n^2)Q_k(\delta n^2)
 +C_aQ_k(64\delta)\frac{U_{16}(32\delta)}{1-U_{16}(8\delta)}\right],
 \quad C_0=1,\quad C_1=(1+\delta^{-1})/2.
\]
Replacing the actual $d=\pi/q$ by the smaller positive $\delta$ increases
the integral majorant. Exact fraction comparison gives
$\widehat M_3<4$, $\widehat M_4<10$, $\widehat M_5<6$.
The producer recomputes these fractions, asserts the strict comparisons and
exports them; the constants in the readout proof are thus checked rather
than supplied by a floating growth estimate.

\section{The completed-function integral and its tails}
Let $d=\pi/q$ and $\theta_q(t)=\sum_{n\ge1}\chi_q(n)n^ae^{-dn^2t}$.
The standard primitive theta transformation, followed by splitting the Mellin
integral at1, gives
\[
 \Lambda_q(s)=\int_1^\infty\theta_q(t)t^{(s+a)/2-1}\,dt+
              \int_1^\infty\theta_q(t)t^{(1-s+a)/2-1}\,dt.
\]
The second root-number coefficient is1 for these three characters: their
finite Gauss sums are $i\sqrt3$, $2i$ and $\sqrt5$, respectively.
For the last, $2\cos(2\pi/5)-2\cos(4\pi/5)=\sqrt5$; the first two are direct
conjugate-root differences. Division by $i^a\sqrt q$ gives the claimed
normalization. Exponential decay makes the two integrals entire. This classical
theta identity, not an assumption of Artin holomorphy or desired positivity,
is the analytic input of the implementation.

The actual circle nodes are enclosed by computing $e^{2\pi i j/64}$ with
the validated primitives. All input boxes fit
$D=[-1,2]+i[-4/3,4/3]$. Thus both exponents
$p=(s+a)/2-1$ and $p'=(1-s+a)/2-1$ satisfy
$|\Re p|\le2$, $|\Im p|\le2/3$. The domain guard precedes all term discards.
Retain $1\le n\le11$ and integrate $1\le t\le128$.

Put $\delta=3/q<d$, $U(x)=(1+x/128)^{-128}$ and
$Q(x)=x^{-1}+2x^{-2}+2x^{-3}$. Since $(1+x/128)^{128}\le e^x$,
$e^{-x}\le U(x)<1$ for $x>0$. Integration by parts gives
$\int_1^\infty t^2e^{-xt}\,dt=e^{-x}Q(x)$.
With $C_0=1$, $C_1=(1+\delta^{-1})/2$ the omitted-$n$ tail of both integrals is
\[
 2C_aQ(\delta12^2)
 \frac{U(\delta12^2/2)}{1-U(12\delta)}.
\]
Indeed $n^ae^{-\delta n^2/2}\le C_a$, the function $Q$ decreases, and
$(12+j)^2\ge12^2+24j$. These inequalities yield a geometric bound after
splitting the exponent in halves. The retained-$n$ time tail is at most
\[
 \sum_{n=1}^{11}2n^a128^3U(128\delta n^2)Q(128\delta n^2).
\]
It follows by substituting $t=128u$ and majorizing $t^{\Re p}$ by $t^2$.
The two tails are disjoint: omitted $n$ use the full integral, whereas
retained $n$ contribute only $t>128$.

\section{High-order integration with a certified remainder}
Split each dyadic interval $[2^j,2^{j+1}]$, $j=0,\ldots,6$, into four equal
panels. For its midpoint $m$ and halfwidth $h$, $h/m\le1/9$.
For $f(z)=e^{-\lambda z}z^p$, $\lambda=\pi n^2/q$, the principal logarithm
is analytic on $|z-m|\le\rho=m/2$. On that circle,
\[
 |f(z)|\le M_\rho=3\max\{(3m/2)^2,(m/2)^{-2},1\}
 U(\delta n^2m/2).
\]
Here $\Re z\ge m/2$, $|\arg z|\le1$ and
$e^{|\Im p|}\le e^{2/3}\le3$. The maximum handles both signs of $\Re p$.
Cauchy's coefficient estimate and cancellation of odd powers on $[-h,h]$
therefore bound the error after degree $2K$ by
\[
 \frac{2hM_\rho(h/\rho)^{2K+2}}
 {(2K+3)(1-(h/\rho)^2)}.
\]
Indeed integrate each remaining even term, replace its denominator by the
first omitted denominator, then sum a geometric series. Here $K=12$.
Multiply by $2n^a$ for both Mellin terms and sum all $28$ panels and retained $n$.
Keeping this majorant even for later discarded terms only overcounts error.

To evaluate the polynomial stably, set $u=h/m$, $q_{-1}=0$, $q_0=1$ and use
\[
 q_{j+1}=\frac{u(p-\lambda m-j)q_j-(\lambda m)u^2q_{j-1}}{j+1}.
\]
Multiplying $f'=(-\lambda+p/t)f$ by $t=m+x$ and matching coefficients proves
this recurrence, with $q_j=f^{(j)}(m)h^j/(j!f(m))$. The integrated polynomial is
$2he^{-\lambda m}m^p\sum_{j=0}^{K}q_{2j}/(2j+1)$; all its arithmetic is outward.

Fixed absolute exponential-rounding errors can be amplified by high-order
coefficients in late panels. The code therefore first bounds a complete
$(n,\text{panel})$ pair, including both powers, by
\[
 4hn^a(m+h)^2U(\delta n^2(m-h)).
\]
For real $t\ge1$ its phase has modulus1 and $t^{\Re p}\le(m+h)^2$, proving
the bound. If it is below $2^{-90}$, skip that pair entirely and add its bound
to the final radius. This is controlled truncation, not numerical underflow.
The final sample radius is the interval-sum halfwidth plus the omitted-$n$,
time, Cauchy and discard bounds. All are recorded separately and checked as
nonnegative. No high-precision special-function routine is trusted here.

\section{Reproduction and remaining interfaces}
Run the preserved source on the existing Linux Python environment with no new
packages: \texttt{python3} on \path{examples/analytic_oracle/theta_oracle.py}, with argument \texttt{--output NEW}.
It refuses to overwrite an existing experiment directory. The final-v5
experiment contains $192$ samples, three rational degree $28$ polynomials, copied
source and SHA-256 hashes. Independent validation checks rational endpoint
arithmetic, transcendental enclosures, $1160$ exact coefficient identities, high-
precision completed-function comparisons and the exported component sums.

The fixed-parameter oracle solves a concrete evaluation/certification gap.
It does not identify a positive arithmetic operator, force real-rooted
polynomials, prove the signed original-family short-row bound or supply
complete Lean4 verification. Surface/order and neural architectures may use
these certificates as independent targets, but a small learned loss is not a
substitute for the remaining arithmetic and categorical theorems.

\section{Requested-error readouts from preserved certified labels}\label{ar:requested}
The accepted final-v5 sample certificate is preserved unchanged. For each
$q=3,4,5$, it supplies $m=64$ rational complex labels $y_j$, bounds
$|y_j-\Lambda_q(c+\rho\zeta_j)|\le\eta_j$, and outward rectangles enclosing
the true roots $\zeta_j=e^{2\pi i j/m}$, where $c=1/2$, $\rho=6/5$.
The classical primitive theta/Mellin identity and accepted growth bounds
$M_q=4,10,6$ on $|s-c|\le r=12/5$ remain the analytic inputs.
This extension accepts an exact rational error request $\varepsilon>0$ on
the fixed compact $|s-c|\le R=7/10$. It changes the finite readout, while
retaining the certified sample errors and the $64$-node alias budget.
It is not an arbitrary-accuracy theta-sampling interface.

\begin{theorem}[Executed rational readout with explicit refusal]\label{ar:returned}
For each supported conductor and rational request, the finite-data
selection below either returns a normalized rational polynomial $P$ and an
executed, exactly replayable gradient-descent trace with
\[
 \sup_{|s-c|\le R}|P(s)-\Lambda_q(s)|\le\varepsilon,
\]
or refuses the request with the exact minimum base budget over all
$64$ allowed degrees. Refusal concerns attainability by this particular
fixed-data majorant; it is not an impossibility theorem for approximation.
\end{theorem}
\begin{proof}
Let $B_n$ be an outward rectangle enclosing
$b_n^*=m^{-1}\sum_j y_j\bar\zeta_j^n$, $0\le n<64$, computed from the
preserved root rectangles. Put $t_n=\operatorname{mid}(B_n)$ and let
$e_n$ be the sum of its coordinate halfwidths. These are exact rationals,
and $|t_n-b_n^*|\le e_n$.
For a selected $d<64$ execute, using exact rational operations,
\[
 b_n^0=0,\qquad b_n^{k+1}=\tfrac34 b_n^k+\tfrac14 t_n,
 \qquad 0\le n\le d.
\]
This is genuine gradient descent with step $1/4$ for the explicitly
specified real Euclidean objective
$J(b)=\frac12\sum_{n=0}^d|b_n-t_n|^2$, whose gradient is $b-t$.
All iterates and exact objective values are retained. The output is literally
\[
 P_k(s)=\sum_{n=0}^d b_n^k\left(\frac{s-c}{\rho}\right)^n.
\]
Set $u=R/\rho=7/12$, $v=R/r=7/24$, $w=(\rho/r)^{64}=2^{-64}$,
$S_d(x)=\sum_{n=0}^d x^n$ and $\eta=\max_j\eta_j$. The five budgets are
\begin{align*}
 E_{\rm opt}(k)&=\sum_{n=0}^d|b_n^k-t_n|_1u^n,&
 E_{\rm oracle}&=\eta S_d(u),\\
 E_{\rm alias}&=M_q\frac{w}{1-w}S_d(v),&
 E_{\rm Taylor}&=M_q\frac{v^{d+1}}{1-v},\\
 E_{\rm arithmetic}&=\sum_{n=0}^d e_nu^n.&&
\end{align*}
Here $|x+iy|_1=|x|+|y|$. Cauchy's estimate and true-root Fourier
orthogonality give the alias and Taylor terms; certified label errors give
the oracle term. Triangle inequality applied to
$b_n^k-b_n^*=(b_n^k-t_n)+(t_n-b_n^*)$ gives the other two.
Thus their exact sum bounds the compact-uniform error.

Define $C_d=E_{\rm oracle}+E_{\rm alias}+E_{\rm Taylor}
+E_{\rm arithmetic}$. Select the smallest $d<64$ with $C_d<\varepsilon$;
if none exists, report $\min_{d<64}C_d$ and refuse. For a selected degree,
$E_{\rm opt}(k)=(3/4)^k\sum_n|t_n|_1u^n\longrightarrow0$.
Therefore executing updates until
$E_{\rm opt}(k)\le\varepsilon-C_d$ terminates after finitely many steps.
The serialized coefficients, budgets, targets and trace use exact rational
endpoints; displayed decimal approximations are not proof inputs.
\end{proof}

\section{Objective semantics, precision floors and replay}\label{ar:semantics}
The original true-circle sample-label least-squares objective is
\[
 \mathcal L_y(b)=\frac1{2m}\sum_{j=0}^{m-1}
 \left|\sum_{n=0}^d b_n\zeta_j^n-y_j\right|^2.
\]
With $d<m$, Fourier orthogonality gives $A^*A=mI$ and the real Euclidean
coordinate gradient $\nabla\mathcal L_y(b)=b-b^*$. The executed objective
$J$ instead has the exact rational midpoint target $t$.
The difference $\nabla\mathcal L_y-\nabla J=t-b^*$ is enclosed
coefficientwise by $e_n$.
Consequently this is an actual rational GD run for $J$ and a certified
perturbed-gradient implementation for $\mathcal L_y$. It is not the exact
optimizer trajectory of the ideal-root objective, nor a replay of neural or
quantum training. An interval enclosure of an ideal-root GD trajectory
would be a different explicitly labelled object.

Each actual update contracts $b^k-t$ by $3/4$, so the exact stored losses
satisfy $J(b^{k+1})=(9/16)J(b^k)$. A replay recomputes every rational
coordinate update, loss and weighted residual. The closed form
$b^k=(1-(3/4)^k)t$ is an independent check on the executed trace, not its
producer. Exact rational loop coordinates introduce no further midpoint or
floating export error. The arithmetic budget covers the difference between
the certified Fourier rectangles and their midpoint targets, including the
inherited root enclosure and outward interval operations. A subsequent
float conversion requires its own exact conversion-error bound.

The normalized polynomial convention is essential: a stored coefficient
with index $n$ multiplies $((s-1/2)/(6/5))^n$, not $(s-1/2)^n$.
The input certificate is pinned by SHA-256, including its source hashes,
all $192$ sample enclosures, root rectangles and analytic radius components.
Producing source is captured at entry, checked unchanged before emission,
and copied into a fresh output directory. An existing output is never
overwritten. An independent replay should also check the fixed input hash,
all dimensions, degree choice, the first stopping iterate, all five exact
terms and their literal sum against the requested rational tolerance.

The attainable boundary of this majorant is
\[
 C_{\min,q}=\min_{0\le d<64}C_{q,d}.
\]
Every request $\varepsilon>C_{\min,q}$ has strict budget slack for some
degree and the proved finite stop rule. At or below that boundary this
interface refuses. The degree minimizing $C_d$ need not be $63$:
Taylor error decreases with degree while the oracle and arithmetic sums
increase. In particular $C_d\ge\eta$ and
$C_d\ge M_q2^{-64}/(1-2^{-64})$. Increasing the number of optimization
steps does not lower these fixed-data budgets. Approximate boundary values
on the stated compact are $7.2929\cdot10^{-17}$,
$1.7370\cdot10^{-16}$ and $3.3741\cdot10^{-16}$ for $q=3,4,5$;
only the exact exported rational minima are used for acceptance.
These are certificate floors, not lower bounds on the actual errors.

A full adaptive theta oracle would separately need proof-preserving choices
of omitted-$n$ cutoff, time cutoff, quadrature order, whole-panel discard
budget, primitive series lengths, interval precision and serialized-bound
precision. The present $N=11$, $T=128$, $224$-bit arithmetic,
$60/92/60$ primitive series lengths, $2^{-90}$ discard threshold and
$180$-bit rounded budgets cannot support every positive request.
Increasing a nominal bit count alone leaves the fixed primitive remainders
and analytic tails in place. The preserved labels are not recomputed here.
The classical identity, finite arithmetic and informal proof are not new
Lean declarations. No zero-location, all-conductor, positive arithmetic
determinant, quantum-advantage or GRH conclusion follows.

The interactive replay is available at \href{https://amelie-iska.github.io/GRH-research/readout.html}{the certified readout laboratory}. All decimal graphics are illustrative views of the linked exact records.

\section{Requested-error theta samples for three fixed conductors}\label{at:sampling}
For the primitive characters $\chi_3=(0,1,-1)$, $\chi_4=(0,1,0,-1)$,
$\chi_5=(0,1,-1,-1,1)$ with parities $a=1,1,0$, retain the classical
primitive theta/Mellin identity and root number one already stated above.
Put $D=[-1,2]+i[-4/3,4/3]$ and $\delta=3/q$.
The following adaptive extension is implemented with independent interval
and sample checks; it is an informal analytic theorem, not a Lean result.
\begin{theorem}[Finite adaptive sample specification]\label{at:sample}
For every supported conductor, exact rational $s\in D$, and rational
$\eta>0$, the choices below terminate and give a rational complex midpoint
$y$ and exact rational radius $R\le\eta$ with
$|y-\Lambda_q(s)|\le R$.
The same conclusion holds for a specified true circle node when its
root and node enclosures are refined on every precision retry.
\end{theorem}
\begin{proof}
For $x\ge0$ define the exact rational bound $E(x)=2^{-\lfloor x\rfloor}$.
Since $e>2$, $e^{-x}\le E(x)$; this bound is nonincreasing and tends to zero.
Let $Q(x)=x^{-1}+2x^{-2}+2x^{-3}$ and
$C_0=1$, $C_1=(1+\delta^{-1})/2$.
For $N\ge1$, $n_0=N+1$, $J\ge1$, $T=2^J$, use
\begin{align*}
 B_n&=2C_aQ(\delta n_0^2)
       \frac{E(\delta n_0^2/2)}{1-E(\delta n_0)},\\
 B_t&=\sum_{n=1}^N2n^aT^3E(\delta n^2T)Q(\delta n^2T).
\end{align*}
Here $\delta n_0\ge6/5$ ensures a positive denominator; $E(x)$ equals
one for $0\le x<1$, so this safeguard is essential.
On the four panels of each $[2^j,2^{j+1}]$, write midpoint $m$,
halfwidth $h$, $\rho=m/2$, $v=h/\rho\le2/9$,
and $B(m)=\max\{(m+\rho)^2,(m-\rho)^{-2},1\}$.
Use the Cauchy remainder
\[
 B_Q=\sum_{\mathrm{panels}}\sum_{n=1}^N
 \frac{2n^a\,2h\,3B(m)E(\delta n^2m/2)v^{2K+2}}
 {(2K+3)(1-v^2)}.
\]
Choose in turn the first $N\ge1$, first $J\ge1$, and first $K\ge0$
for which their respective budgets are at most $\eta/8$.
Gaussian decay, then exponential time decay, then $v\le2/9$ prove
termination. These are the accepted tail and Cauchy arguments with each
true exponential separately replaced by its upper bound $E$.

There are $P=4JN$ potential pairs. A nonzero-character pair may be
omitted only if its entire two-power contribution has bound
\[
 b=4hn^a(m+h)^2E(\delta n^2(m-h))\le\eta/(8P).
\]
Add every omitted bound to $B_{\rm discard}$, giving
$B_{\rm discard}\le\eta/8$.
Evaluate all other integrated Taylor polynomials by the stated exact
coefficient recurrence using outward rectangles. Refine both interval
precision and all primitive series remainders until the final rectangle's
sum of coordinate halfwidths satisfies $R_{\rm arith}\le\eta/4$.
For this fixed finite expression all input and primitive intervals shrink;
continuity and outward arithmetic imply $R_{\rm arith}\to0$.
Precision doubling therefore terminates, and the sum of the five raw
components is at most $3\eta/4$.
Choose an integer output scale $B$ such that $5\,2^{-B}\le\eta/8$ and
round each component upward by
$\operatorname{up}_B(x)=2^{-B}\lceil2^Bx\rceil$.
The reported radius is their literal exact sum. With
$\mathbf b=(R_{\rm arith},B_n,B_t,B_Q,B_{\rm discard})$, it obeys
\[
 |y-\Lambda_q(s)|\le\sum_{i=1}^5\operatorname{up}_B(b_i)
 \le\frac{3\eta}{4}+5\,2^{-B}\le\frac{7\eta}{8}<\eta.
\]
The output retains exact rational midpoints, rounded components, $B$,
raw interval radius, parameters, input metadata and source hashes.
The raw analytic bounds are exactly reconstructible; they need not be
serialized as potentially enormous fractions. The output precision varies
with the request and introduces no fixed reporting floor.
\end{proof}

This fixed-conductor sample result uses no hypothesis on zeros.
It does not remove a downstream fixed-node Fourier alias floor, supply
an all-conductor theorem or a positive arithmetic determinant, or prove GRH.

\section{Adaptive sampling parameters and precision termination}\label{at:details}
The primitive theta/Mellin identity with root number one yields
\[
 \Lambda_q(s)=\sum_{n\ge1}\chi_q(n)n^a
 \int_1^\infty e^{-\pi n^2t/q}
 \bigl(t^{(s+a)/2-1}+t^{(1-s+a)/2-1}\bigr)\,dt.
\]
On $D=[-1,2]+i[-4/3,4/3]$, both exponents $p$ satisfy
$|\Re p|\le2$ and $|\Im p|\le2/3$.
For exact $x\ge0$, put $E(x)=2^{-\lfloor x\rfloor}$.
The elementary inequality $e>2$ gives
$e^{-x}\le e^{-\lfloor x\rfloor}\le E(x)$ and
$E(x)\le2e^{-x\log2}$.
Use this solely as a pointwise nonincreasing upper bound; it is not
multiplicative and has no asserted exponential integral identity.

Integration by parts gives
$\int_1^\infty t^2e^{-xt}\,dt=e^{-x}Q(x)$ for
$Q(x)=x^{-1}+2x^{-2}+2x^{-3}$.
For $a=0$ the bound $n^ae^{-\delta n^2/2}\le C_0=1$ is immediate.
For $a=1$, $n\le(1+n^2)/2$ and
$n^2e^{-\delta n^2/2}\le2/(e\delta)\le1/\delta$ give
$C_1=(1+1/\delta)/2$.
For $n=n_0+j$, the inequality
$n^2\ge n_0^2+2n_0j$ therefore bounds the omitted-$n$ sum by
\[
 2C_aQ(\delta n_0^2)
 \frac{e^{-\delta n_0^2/2}}{1-e^{-\delta n_0}}
 \le2C_aQ(\delta n_0^2)
 \frac{E(\delta n_0^2/2)}{1-E(\delta n_0)}=B_n.
\]
Take $n_0=N+1\ge2$; then $\delta n_0\ge6/5$ and the latter denominator
is positive. Substituting $t=Tu$ in each retained-$n$ time tail gives
\[
 B_t=\sum_{n=1}^N2n^aT^3E(\delta n^2T)Q(\delta n^2T).
\]
The two tails are disjoint because the omitted-$n$ tail uses the whole
integral and the time tail uses only retained $n$.

For $j=0,\ldots,J-1$ and $\ell=0,\ldots,3$ take
\[
 m=2^j\left(1+\frac{2\ell+1}{8}\right),\qquad h=\frac{2^j}{8},
 \qquad T=2^J.
\]
These panels cover $[1,T]$ and have $h/m\le1/9$.
On $|z-m|=\rho=m/2$, principal log is analytic,
$\Re z\ge m/2$, and $|\arg z|\le1$.
For $f(z)=e^{-\pi n^2z/q}z^p$, the angular multiplier is at most
$e^{2/3}\le3$, whence
\[
 |f(z)|\le3B(m)E(\delta n^2m/2),\qquad
 B(m)=\max\{(3m/2)^2,(m/2)^{-2},1\}.
\]
Cauchy's coefficient estimate, cancellation of odd terms, and geometric
summation yield the $B_Q$ in Theorem~\ref{at:sample}.
Including zero-character or later discarded terms in that sum is harmless.
Gaussian, time, and geometric decay prove termination of the sequential
first-success searches with $B_n,B_t,B_Q\le\eta/8$.
For fixed q and $L=\max\{1,\log(1/\eta)\}$ their asymptotic scales can be
chosen as $N=O(\sqrt L)$, $T=O(L)$, $J=O(\log L)$, and $K=O(L)$.
Polynomial prefactors in $N,T,J$ contribute only $O(\log L)$ to the
logarithm of the error estimates; this is a truncation-size statement,
not a rational-arithmetic bit-complexity bound.

For real $t\in[m-h,m+h]$, $t\ge1$ and $t^{\Re p}\le(m+h)^2$.
The bound for a whole pair, including both powers, is therefore
$4hn^a(m+h)^2E(\delta n^2(m-h))$.
There are $P=4JN$ potential pairs; summing only bounds at most
$\eta/(8P)$ gives $B_{\rm discard}\le\eta/8$.
No partial recurrence or isolated exponential is discarded.
For active pairs retain the coefficient recurrence
\[
 r_{-1}=0,\quad r_0=1,\quad
 r_{j+1}=\frac{(h/m)(p-\lambda m-j)r_j-
                   \lambda m(h/m)^2r_{j-1}}{j+1},
 \qquad\lambda=\pi n^2/q.
\]
Coefficient comparison in $t f'=(-\lambda t+p)f$ proves that the
integrated Taylor polynomial is
\[
 2he^{p\log m-\lambda m}\sum_{j=0}^K\frac{r_{2j}}{2j+1}.
\]
Outward arithmetic encloses its exact value, including the signed
character weight and both powers.

Use an isolated p-bit dyadic context, with all primitive remainders
refined as p increases. Put $\delta_p=2^{-p-16}$.
Each arctangent series separately stops when its next-term bound is at
most $\delta_p$, giving total weighted analytic tail width at most
$20\delta_p$ in Machin's formula, before outward dyadic rounding.
Each reduced logarithm series separately has tail at most $\delta_p$;
for $x=2^kv$, the combined bound is $(1+|k|)\delta_p$, before rounding.
For the complex exponential, after reduction to $A\le1/2$, choose the
Taylor order so the bound
\[
 \frac{A^{H+1}}{(H+1)!}\frac1{1-A/(H+2)}.
\]
is at most $\delta_p$. Choose H by finite exact search; its denominator
is positive. Add a symmetric exponential error enclosure to each coordinate
before outward squaring. The fixed guard16 is not itself an acceptance
bound: the final complete rectangle radius must satisfy
$R_{\rm arith}\le\eta/4$. All interval caches must include p and preserve
exact input-type guards. The successful p and retry precisions/radii are
recorded. Individual series lengths are determined by these, the exact
inputs and the pinned source and are reconstructible by replay; they are
not explicitly stored.

For fixed N,J,K and exact target point, the expression has finitely many
operations and only nonzero fixed integer denominators. Each primitive
input is bounded for sufficiently large p; exponential halving counts
are consequently bounded. Exact primitive tails and dyadic rounding
widths tend to zero. Induction through the finite continuous expression,
including finitely many outward squarings, gives $R_{\rm arith}\to0$.
Doubling p after a failed comparison terminates. The argument also covers
true roots and nodes regenerated at every p, since their widths tend to zero.
A fixed nondegenerate s rectangle has no corresponding shrinking-input
claim.

\paragraph{A growing-order primitive estimate.}
To justify convergence despite a growing exponential Taylor order, write
$\epsilon_p=2^{-p}$, $\delta_p=2^{-p-16}$, and let $d(X)$ be the sum of a
rectangle's coordinate widths and $m(X)$ its $\ell^1$ magnitude bound.
The rounded reduced input $W$ satisfies $a=m(W)\le1/2+2\epsilon_p<3/5$
and $d(W)\le2^{-k}d(Z)+4\epsilon_p$; the remainder is at most $2^{-H}$,
so the first accepted order satisfies $H\le p+18$.
For the computed term boxes, outward multiplication and integer division give
$m_j\le a m_{j-1}+6\epsilon_p$ and
$d_j\le a d_{j-1}+m_{j-1}d(W)+12\epsilon_p$.
Thus $m_j\le a^j+15\epsilon_p$ and
$\sum_{j=0}^H d_j\le(25/4)d(W)+75H\epsilon_p$; adding the remainder
increases width by at most $4(\delta_p+\epsilon_p)$.
The pre-squaring width is therefore $O(d(Z)+p2^{-p})$ and its magnitude
is bounded (indeed below $3$ for $p\ge32$).
Each outward square has width at most $2m(X)d(X)+8\epsilon_p$;
the bounded number of squarings consequently preserves convergence to zero.
Finally, exact rational partial sums and endpoint rounding give Machin width
at most $20\delta_p+40\epsilon_p$, and logarithm width at most
$(1+|k|)(\delta_p+2\epsilon_p)$ for fixed $x=2^kv$.
These estimates justify the primitive-width step without assuming a fixed
Taylor order or asserting a global arithmetic complexity bound.


For serialization, initialize
\[
 B=\max\{32,\operatorname{bitlen}(\operatorname{den}\eta)
             -\operatorname{bitlen}(\operatorname{num}\eta)+8\}
\]
and increase B until $5\,2^{-B}\le\eta/8$.
Write $\widehat x=2^{-B}\lceil2^Bx\rceil$ for each of the five raw
components. Their exact reported sum is at most
$3\eta/4+5\,2^{-B}\le7\eta/8$, and every component remains an upper
bound. The output records B, the raw interval radius and complete rectangle,
N/J/K, discard threshold/pair counts, exact target metadata and source hashes.
These reconstruct all raw analytic bounds and pair sets without serializing
huge exact tail denominators. For the practical requests below B is87
and114, respectively. Exact endpoints remain rational; the mathematical
specification assumes an integer representation sufficient for the requested
size. Dynamic rounding removes the old fixed reporting floor. For rational s, validate its exact coordinates in closed D before
conversion: an outward dyadic enclosure of $\Im s=4/3$ extends beyond D,
but the analytic bounds concern the true point. Circle nodes satisfy the
strict exact bounds $-7/10\le\Re s\le17/10$, $|\Im s|\le6/5$.

The exact parameter-design checker gives
\[
\begin{array}{c|c|rrrr|r}
\eta&q&N&J&T&K&\text{active pairs}\\\hline
10^{-24}&3&12&7&128&19&71\\
10^{-24}&4&14&8&256&20&65\\
10^{-24}&5&16&8&256&20&111\\
10^{-32}&3&14&7&128&25&84\\
10^{-32}&4&16&8&256&26&72\\
10^{-32}&5&18&8&256&26&129
\end{array}
\]
with exact rational budgets and no special-function oracle.
The arithmetic precision is selected from the actual interval radius,
not inferred from this table. This extension retains the existing
classical theta identity; it adds no zero-location, positive determinant,
all-conductor, Lean or GRH claim. The following end-to-end adaptive readout supplies its own node-count,
alias, degree, Fourier and optimization choices.

\section{An end-to-end adaptive analytic readout}\label{at:family}
The new driver changes the sample truncations, interval precision, circle
size and degree; it does not reuse the frozen labels to assert a smaller
error. Its scope is still the three stated primitive quadratic characters
and the fixed compact $|s-1/2|\le7/10$.

\begin{theorem}[Requested-error construction in the stated finite family]\label{at:requested}
Given a positive rational $\varepsilon$ and $q\in\{3,4,5\}$, the specified
exact-arithmetic algorithm terminates in the mathematical unbounded-resource
model and returns a normalized rational polynomial $P$ with
\[
 \sup_{|s-1/2|\le7/10}|P(s)-\Lambda_q(s)|\le\varepsilon.
\]
Its sample enclosures, rational optimizer updates and all error budgets are
explicit. This is an approximation theorem, with the stated classical theta
identity as an input; it supplies no zero-location or positive determinant
conclusion and is not a Lean theorem or a practical runtime bound.
\end{theorem}
\begin{proof}
Put $u=7/12$, $v=7/24$. Choose $d\ge0$ so
$10v^{d+1}/(1-v)\le\varepsilon/8$, and choose $m>d$ so
\[
 10\frac{2^{-m}}{1-2^{-m}}\frac1{1-v}\le\varepsilon/8.
\]
The driver rounds $m$ upward to a multiple of eight, which only decreases
this bound. Set the requested sample error
$\eta=\varepsilon(1-u)/8$. The adaptive sampler gives all $m$ true circle
nodes and labels with error at most $\eta$. Form their outward Fourier
rectangles through degree $d$, and refine the nodes and finite evaluation
arithmetic if their weighted midpoint radii exceed $\varepsilon/8$.
This refinement terminates: for these fixed finite $d,m$, node uncertainty
and every finite-expression arithmetic width tend to zero with precision.
Merely lifting fixed boxes would not suffice, so the driver regenerates the
nodes and samples on this retry.

Choose rational midpoint targets $t_n$ and execute the previously specified
quadratic-objective updates until
$\sum_{n=0}^d|b_n^k-t_n|_1u^n\le\varepsilon/4$.
Their geometric contraction proves finite termination. The oracle,
aliasing, Taylor and Fourier-midpoint terms are each at most
$\varepsilon/8$, and the optimization term is at most $\varepsilon/4$.
Hence the raw five-term sum is at most $3\varepsilon/4$.
Choose a serialization precision $B$ with $5\,2^{-B}\le\varepsilon/8$,
round each term upward, then add the five reported rationals exactly.
The printed bound is at most $7\varepsilon/8\le\varepsilon$.
The polynomial uses $((s-1/2)/(6/5))^n$, and the preceding Fourier/Cauchy
proof gives its asserted compact-uniform bound.
\end{proof}

The trace is losslessly encoded by the fixed target vector and a rational
scale extracted from each \emph{executed} coefficient state, together with
its objective, weighted residual and full-state hash. Every coefficient
can be recovered exactly; a separate replay expands the states and checks
every update. The producer does not generate its trace from a closed form.
The saved objective remains the rational Fourier-midpoint quadratic, with
its perturbation from ideal-root least squares explicitly budgeted.

\section{Termination, compression and the boundary of the guarantee}\label{at:family-details}
The degree and circle-size choices use only rational comparisons. Since
$0<v<1$, the Taylor test succeeds at finite $d$; since $2^{-m}\to0$, the
alias test succeeds at finite $m>d$. The common growth bound10 is safe for
all three families; the certificate also recomputes the previously proved
exact conductor-specific majorants below $4,10,6$.

At fixed $d,m,\eta$, the analytic truncation parameters are fixed while
precision is doubled. The computed sample centers stay bounded: their
underlying finite expression is fixed and the true completed-function
values are bounded on the sample circle. Its arithmetic widths tend to
zero, as do the regenerated true-root enclosures. Fourier computation is
a finite composition of additions, products and division by the positive
integer $m$. Its midpoint-radius sum therefore tends to zero. This proves
the driver's additional precision retry, without pretending that changing
a box's nominal precision shrinks its uncertainty.

The accepted sampler error may remain above its actual arithmetic width
because its analytic tails and output bounds are fixed for this $\eta$.
That causes no problem: the oracle allocation uses the total certified
sample radius, whereas the Fourier-midpoint allocation uses only finite
coefficient-rectangle width. No analytic sample error is silently moved
into the latter term or removed by numerical convergence.

For the stored training trace, the exact loop gives $b^{k+1}=3b^k/4+t/4$
from $b^0=0$. Inductively every state is a rational scalar multiple of the
same target vector. The producer first executes the actual vector updates,
then extracts this scalar from a nonzero target coordinate, checks all
coordinates, and records a canonical full-vector SHA-256 hash. If the
entire target vector is zero, its zero state needs only scale0. The compressed
record is thus an exact representation of the executed states, not a
floating-point sketch. Hashes are integrity aids; exact vector reconstruction
and the independently checked recurrence establish the mathematical claim.

Coefficient export uses exact rational numerators and denominators. A
chunked decimal integer encoder/parser handles endpoints beyond Python's
usual digit-conversion guard without changing global interpreter settings.
No fixed $180$-bit serialization floor remains. The decimal graphics are
separate approximate projections, not polynomial coefficients to which the
exact certificate may be applied without a conversion-error bound.

All producing source bytes are captured, copied into fresh output paths and
checked unchanged. Historical fixed producers are pinned and untouched.
For the new modules the run also checks the import-time source snapshots,
so a stale imported module cannot be attributed to later source bytes.
After output creation, a failed evaluation, final provenance check or
serialization emits uncertified failure metadata. Invalid requests and stale
import-time sources are refused before creating output artifacts. Old
outputs are never overwritten. Extreme requests can be
computationally impractical: the termination proof gives no machine-time or
memory guarantee, and finite test runs establish only their reported cases.

Nothing requires the output polynomials to be real-rooted, to equal an
arithmetic Fredholm determinant, or to have a self-adjoint realization.
Consequently arbitrary compact approximation of these known functions does
not imply GRH. Uniformity in changing conductor, general Hecke/Dedekind/
elliptic/Artin scopes, the independent determinant comparison, and actual
analytic/category Lean objects remain separate obligations.

% Accepted prose AFM proof; explicit external inputs, not a Lean theorem.
\section{An arithmetic-family moment with every original row retained}
\label{mm:core}
Let $K=\mathbb Q(\sqrt{-3})$, $\mathcal O=\mathbb Z[\omega]$, and write
$N$ for the ideal or element norm. Fix a finite-order Hecke character
$\nu$ and a finite set $S$ of prime ideals containing those above $6$ and
those dividing its conductor. For $0\ne u\in\mathcal O$, let $\kappa_u$
be the character of $K(u^{1/6})/K$ defined using arithmetic Frobenius, and
put $\psi_u=\operatorname{prim}(\nu\kappa_u)$. Here
$\operatorname{prim}$ means the underlying primitive Hecke character:
it is formed before introducing any divisor variable or evaluating
conductor-dependent zero extensions. The rows below are all actual
Eisenstein elements, including units and prime powers supported in $S$.

The external inputs are abelian Kummer reciprocity and character
uniqueness, the primitive Hecke functional equation and its usual
vertical growth, Stirling's formula, scalar Dedekind-zeta convexity, and
Blomer--Goldmakher--Louvel's large sieve for their specified sixth-order
family. We use the pinned
\href{https://arxiv.org/pdf/1112.1650v1}{arXiv:1112.1650v1},
Theorem~1.3, Lemma~2.1 and Appendix~II; its Section~5 supplies the
all-ideal decomposition whose fixed-twist proof is redone here.
No statement about all primitive characters of comparable conductor,
Artin holomorphy, Lindel\"of bounds or zero locations is assumed.

\begin{theorem}[The original-family numerator moment]
\label{mm:afm}
For every $\varepsilon>0$, uniformly for $H\ge1$ and $t\in\mathbb R$,
\begin{equation}\label{mm:central}
 \sum_{0<Nu\le H}|L(1/2+it,\psi_u)|^2
 \ll_{\nu,S,\varepsilon}H^{1+\varepsilon}(1+|t|)^{1+\varepsilon}.
\end{equation}
For every fixed $0<\delta<1/2$, the same estimate holds with
$1/2+it$ replaced by $1/2+\delta+it$, with the constant also depending
on $\delta$. If the central bound has constant $C_0$ at an exponent
$0<\varepsilon_0\le1/4$, an admissible off-central constant is
\begin{equation}\label{mm:poisson-constant}
 200 C_0\left(\frac{3/2+\delta}{1/2-\delta}\right)^2.
\end{equation}
\end{theorem}
\begin{proof}
For the canonical BGL character $\vartheta_{\mathfrak a}$,
Appendix~II gives a radicand $\alpha_{\mathfrak a}$ with
$(\alpha_{\mathfrak a})_S=\mathfrak a\mathfrak g^{-6}$.
Since $\mathcal O_S$ is principal, comparison with any literal $u$
whose outside-$S$ ideal is $\mathfrak a$ gives
$\kappa_u=\eta\vartheta_{\mathfrak a}$ for one of finitely many fixed
$S$-unit Kummer twists $\eta$. This is an equality of primitive
characters, established outside a finite set and extended by uniqueness.
For each outside-$S$ ideal $\mathfrak a$ and $S$-supported integral
ideal $\mathfrak b$, at most six elements generate
$\mathfrak a\mathfrak b$. The resulting $S$-supported multiplicity sum
converges at exponent $1+\varepsilon$.

For each fixed $\xi=\nu\eta$, split the canonical parameter as
$\mathfrak a=\mathfrak a_0\mathfrak a_1\mathfrak a_2^6$, where
$\mathfrak a_0$ is squarefree, $\mathfrak a_1$ is squarefull and
sixth-power-free, and $(\mathfrak a_0,\mathfrak a_1)=1$.
The primitive conductor is at most a fixed multiple of
$N(\mathfrak a_0\mathfrak a_1)$.
The common-gamma approximate functional equation, with
$X=H/(N\mathfrak a_2)^6$ and $T=1+|t|$, has common maximal length
$B=X^{1/2}T(HT)^\eta$ for any sufficiently small fixed $\eta>0$.
Mellin inversion removes its row-dependent conductor weight.
Split a column as $\mathfrak b_S\mathfrak b_1\mathfrak b_2$, with
$\mathfrak b_S$ supported in $S$, $\mathfrak b_1$ squarefree and
$\mathfrak b_2$ squarefull, the last two coprime and prime to $S$.
The original primitive values at $\mathfrak b_S$ remain bounded row
multipliers. The $S$-supported and squarefull column weights are summable
at $1/2+\eta$. Multiplying squarefree column coefficients by the fixed
$\xi$ preserves their squared norm, whatever its finite order.
BGL's large sieve therefore bounds the energy for fixed
$\mathfrak a_1,\mathfrak a_2$ by
\[
 (HT)^{O(\eta)}\{M+B+(MB)^{2/3}\},\qquad
 M=X/N\mathfrak a_1.
\]
Squarefull counting and summation give
$X+B\sqrt X+(XB)^{2/3}\ll (HT)^{O(\eta)}XT$.
The $\mathfrak a_2$ sum converges at exponent $6$.
Principal rows are separated before the pole-free functional equation
and bounded by scalar zeta convexity. This proves~\eqref{mm:central}.

For the off-central statement divide each $L(s,\psi)$ by
$(s+1)^{5/4}$ and, on principal rows, also multiply by $(s-1)/(s+1)$.
These are bounded analytic functions on $\operatorname{Re}s>1/2$,
continuous on the boundary. Poisson representation and Jensen's
inequality transfer the central squared moment. The convolution of
$P_\delta(y)=\delta/[\pi(\delta^2+y^2)]$ with
$(1+|t|)^{-3/2+\varepsilon_0}$ is at most eleven times that weight
when $0<\varepsilon_0\le1/4$. Restoring the denominator and the
principal pole-removal factor gives~\eqref{mm:poisson-constant}.
The extended edition gives the detailed derivation with all primitive local values.
\end{proof}

\paragraph{Separate fixed-unitary extension.}
The moment theorem also holds for any one fixed unitary Hecke character
$\nu$. For precision choose
$\nu(\mathfrak a)=\nu_0(\mathfrak a)(N\mathfrak a)^{-i\tau}$, with
angular parameter $m\in\mathbb Z$ and completed gamma factor
\[
 \gamma_\nu(s)=2(2\pi)^{-(s+|m|/2+i\tau)}
                    \Gamma(s+|m|/2+i\tau).
\]
All finite Kummer twists have this same infinity type. Use
$T_\nu=1+|t+\tau|+|m|$, comparable to $1+|t|$ with
$\nu$-dependent constants. Pure norm rows have
$L(s,\psi)=\zeta_K(s+i\tau)$ and are separated before the pole-free
functional equation; their pole at $1-i\tau$ is removed by
$(s-1+i\tau)/(s+1+i\tau)$. No uniformity in $m$, $\tau$ or the
conductor of $\nu$ is asserted. This extension of the moment is separate
from the following finite-order type-I residue formula.

\subsection{The exact masked remainder and the remaining loss}
Retain the original zero-extended symbol $\chi_d(u)=(u/d)_6$ on good
squarefree ideals. Let $\mu$ be the ideal M\"obius function and set
\[
 \beta_U(d)=\sum_{ab=d\atop Na,Nb\le U}\mu(a)\mu(b),\quad
 T(u,d)=S\cup\operatorname{supp}(u)\cup\operatorname{supp}(d).
\]
Let $\psi_{u,2}=\operatorname{prim}(\psi_u^2)$, formed independently
of squaring any zero sentinel. With $\mathfrak f_u$ and
$\mathfrak f_{u,2}$ the primitive conductors, the full-mask series is
\begin{align}
 F_{u,d}(s)
 &=\frac{L(s,\psi_u)}{L(2s,\psi_{u,2})}C_{T(u,d)}(s),\label{mm:masked}\\
 C_T(s)
 &=\frac{\prod_{p\in T,\ p\nmid\mathfrak f_u}
              (1-\psi_u(p)(Np)^{-s})}
         {\prod_{p\in T,\ p\nmid\mathfrak f_{u,2}}
              (1-\psi_{u,2}(p)(Np)^{-2s})}.\notag
\end{align}
Initially it equals
$\sum_{c\ \mathrm{squarefree},\ (c,T)=1}\psi_u(c)(Nc)^{-s}$
for $\operatorname{Re}s>1$. Every prime of $T$ remains deleted,
including those forgotten by either primitive conductor.

For finite-order $\nu$, $D\ge1$, $U\ge1$, $U^2\le D$, and
$W\in C_c^\infty((0,\infty))$ supported in $[1,2]$, put
\[
 \widehat W(s)=\int_0^\infty W(x)x^{s-1}\,dx,\qquad
 \mathcal M_4(W)=\sum_{j=0}^4\sup_x|(x\partial_x)^jW(x)|.
\]
On $\sigma=1/2+\delta$, the exact residue-subtracted type-I remainder is
\begin{equation}\label{mm:remainder}
 E_u=\frac{D^\delta}{2\pi}\int_{\mathbb R}
 \widehat W(\sigma+it)D^{it}
 \sum_{d\ \mathrm{squarefree}\atop(d,S)=1}
 \frac{\beta_U(d)\nu(d)\chi_d(u)}{(Nd)^{\sigma+it}}
 F_{u,d}(\sigma+it)\,dt.
\end{equation}
This definition retains the actual $\beta_U$, all phases, $d$-masks,
conductors and the independently primitive square.

\begin{theorem}[An improved masked residual estimate]
\label{mm:residual}
Under the preceding finite-order hypotheses, for every $\varepsilon>0$,
\begin{equation}\label{mm:energy}
 \sum_{0<Nu\le H}|E_u|^2
 \ll_{\nu,S,\delta,\varepsilon}
 \mathcal M_4(W)^2D^{2\delta}U^{2-4\delta}H(HD)^\varepsilon.
\end{equation}
\end{theorem}
\begin{proof}
The same deleted set in both factors of~\eqref{mm:masked} gives
$|F_{u,d}(\sigma+it)|\ll (Nu\,Nd)^\eta|L(\sigma+it,\psi_u)|$:
the reciprocal incomplete Euler product at $2\sigma>1$ is bounded
by $\zeta_K(2\sigma)$, and its remaining finite numerator is bounded
by a small power of its radical. Apply Minkowski to~\eqref{mm:remainder}
in the exact row space and use Theorem~\ref{mm:afm}. The actual coefficient
sum is at most
\[
 \sum_d\frac{|\beta_U(d)|}{(Nd)^{\sigma-\eta}}
 \le\left(\sum_{Na\le U}(Na)^{-\sigma+\eta}\right)^2
 \ll U^{1-2\delta+2\eta}.
\]
Four Mellin integrations by parts bound the remaining height integral
by a constant times $\mathcal M_4(W)$. Square the result and absorb
the small losses using $U^2\le D$.
\end{proof}

At $H=D^{3/4}$ and $U=D^a$ with fixed $0<a\le1/2$, the upper bound's excess
over $H$ is $2\delta+(2-4\delta)a$, up to small losses; as
$\delta\downarrow0$ it still leaves $2a$. This is a limitation of the
upper bound, not an energy lower bound. The extracted principal residue,
the original short term at smaller scales and the signed type-II minus
type-I covariance are not controlled by~\eqref{mm:energy}. The next
arithmetic task is to estimate the complete $d$-weighted moment before
its absolute coefficient sum, or establish sufficient signed covariance.
No positive arithmetic determinant, zero-location theorem or GRH proof
follows. The analytic argument uses the external classical inputs above;
finite diagnostics and selected existing Lean local identities do not
formalize the Hecke functions, approximate functional equation, Poisson
transfer or these moment bounds.

% Full derivation of the scoped arithmetic-family moment and masked remainder.
\section{Family identification, the approximate functional equation and masks}
\label{mm:details}
We prove the analytic statements of Section~\ref{mm:core} from the
explicit classical inputs stated there. Write $I(S)$ for the nonzero
integral ideals prime to $S$, and $\mathcal O_S$ for the ring of
$S$-integers. Constants may depend on the fixed field, $S$, $\nu$ and
any declared positive small exponent. No such constant is asserted
uniform in the conductor or infinity type of a varying $\nu$.

\subsection{A finite comparison with the literal Kummer rows}
Use Appendix~II of BGL, with $n=6$. Its fixed local modulus and
representatives may be chosen once, since $S$ contains the primes above
$6$ and $\mathcal O_S$ has class number one. For each
$\mathfrak a\in I(S)$ the construction writes
\[
 \mathfrak a=(x)E\mathfrak g^6,\qquad
 \alpha_{\mathfrak a}=xm_E,\qquad
 (m_E)_S=E,
\]
so $(\alpha_{\mathfrak a})_S=\mathfrak a\mathfrak g^{-6}$.
Choose $y\in K^\times$ generating $\mathfrak g$ as an
$\mathcal O_S$-ideal. If $(u)_S=\mathfrak a$, then
\[
 e=\alpha_{\mathfrak a}y^6/u\in\mathcal O_S^\times,\qquad
 [\alpha_{\mathfrak a}]=[eu]\quad\hbox{in }K^\times/K^{\times6}.
\]
There is a precision here: Appendix~II allows choosing $\mathfrak g$
relative to a tested denominator. Fix one admissible global decomposition
first. Its radicand gives the canonical character
$\vartheta_{\mathfrak a}$ on all primes outside the finite set
$S\cup\operatorname{supp}(\mathfrak a\mathfrak g)$ where that
choice applies. Uniqueness of Hecke characters, equivalently Chebotarev
in the finite abelian quotient, then identifies the underlying primitive
characters everywhere. With matching arithmetic-Frobenius conventions,
\begin{equation}\label{mm:unit-comparison}
 \kappa_u=\kappa_{e^{-1}}\vartheta_{\mathfrak a}.
\end{equation}
If the residue-symbol convention is conjugated, conjugate the canonical
family and its coefficient sequence; the same large sieve follows.

The group $\mathcal O_S^\times/(\mathcal O_S^\times)^6$ is finite.
Thus~\eqref{mm:unit-comparison} uses one of a finite set of fixed
Kummer twists $\eta$, all unramified outside $S$ and of bounded
$S$-conductor. This is a comparison of primitive characters, not a
pointwise multiplication rule for independently zero-extended arrays at
bad primes. In particular, a conductor drop must not be replaced by a
product of artificial zeros.

Since $\mathcal O$ is principal with six units, for fixed
$\mathfrak a\in I(S)$ and an $S$-supported integral ideal
$\mathfrak b$ there are at most six $u$ satisfying
$(u)=\mathfrak a\mathfrak b$. For any nonnegative row quantity, the
literal moment is consequently at most
\begin{equation}\label{mm:literal-sum}
 6\sum_{\mathfrak b\ S\text{-supported}\atop N\mathfrak b\le H}
   \sum_{\eta}\sum_{\mathfrak a\in I(S)\atop N\mathfrak a\le H/N\mathfrak b}
 \left|L\bigl(s,\operatorname{prim}(\nu\eta
                       \vartheta_{\mathfrak a})\bigr)\right|^2.
\end{equation}
An ideal bound $X^{1+\varepsilon}T^{1+\varepsilon}$ therefore transfers
to all literal rows, since
\[
 \sum_{\mathfrak b\ S\text{-supported}}(N\mathfrak b)^{-1-\varepsilon}
 =\prod_{p\in S}(1-(Np)^{-1-\varepsilon})^{-1}<\infty.
\]
No quotient of the original unit rows or deletion of $S$-prime powers
has occurred.

\subsection{Fixed twists, exact tame conductors and the squarefree sieve}
Fix $\xi=\nu\eta$ and write
$\rho_{\mathfrak a}=\operatorname{prim}
(\xi\vartheta_{\mathfrak a})$. Its fixed conductor is supported in
$S$ before the Kummer twist. On columns prime to $S$, $|\xi(\mathfrak b)|=1$.
BGL Theorem~1.3 thus gives, for any fixed $\mathfrak a_1$,
\begin{align}\label{mm:ls}
 &\sum_{\mathfrak a_0\in I(S)\ \mathrm{squarefree}\atop N\mathfrak a_0\le M}
 \left|\sum_{\mathfrak b_1\in I(S)\ \mathrm{squarefree}\atop N\mathfrak b_1\le B}
 \lambda_{\mathfrak b_1}\xi(\mathfrak b_1)
 \vartheta_{\mathfrak a_1}(\mathfrak b_1)
 \vartheta_{\mathfrak a_0}(\mathfrak b_1)\right|^2\notag\\
 &\hspace{1cm}\ll (MB)^\varepsilon
 \{M+B+(MB)^{2/3}\}\sum_{\mathfrak b_1}|\lambda_{\mathfrak b_1}|^2.
\end{align}
The first fixed multiplier preserves the coefficient norm; the second
has modulus at most one. Restricting the positive outer row sum is
harmless. A coprimality condition with another fixed ideal belongs to
$\lambda$. In particular, $\nu$ need not have order dividing six and
is not asserted to belong to BGL's finite ray-class character set.

There is a unique valuation decomposition
\[
 \mathfrak a=\mathfrak a_0\mathfrak a_1\mathfrak a_2^6,
 \quad \mathfrak a_0\ \hbox{squarefree},\quad
 \mathfrak a_1\ \hbox{squarefull and sixth-power-free},\quad
 (\mathfrak a_0,\mathfrak a_1)=1.
\]
Indeed, for $v_p(\mathfrak a)=6q+r$, $0\le r\le5$, put the residual
prime in $\mathfrak a_0$ if $r=1$, and its exponent $r$ in
$\mathfrak a_1$ if $2\le r\le5$. BGL Lemma~2.1 gives
$\vartheta_{\mathfrak a\mathfrak b^6}=\vartheta_{\mathfrak a}$ and
$\vartheta_{\mathfrak a_0\mathfrak a_1}
=\vartheta_{\mathfrak a_0}\vartheta_{\mathfrak a_1}$ on ideals
prime to $S$ in this sixth-power-free case. Those are the identities
needed here; no arbitrary multiplicativity assertion is substituted.

At a prime outside $S$, tame Kummer inertia at residual exponent $r$
has order $6/\gcd(6,r)$. It is nontrivial for $r=1,\ldots,5$ and
has conductor exponent one. The fixed $\xi$ is unramified there, so
cannot cancel that inertia. At primes in $S$ the local sixth-power
classes are finite, hence only finitely many primitive conductor types
occur after multiplication by the fixed $\xi$. Consequently
\begin{equation}\label{mm:conductor}
 Q_{\rho_{\mathfrak a}}=|D_K|N\mathfrak f_{\rho_{\mathfrak a}}
 \le C_{\xi,S}N(\mathfrak a_0\mathfrak a_1),
\end{equation}
independently of $\mathfrak a_2$. The equality of the conductor's
outside-$S$ part with $\operatorname{rad}(\mathfrak a_0\mathfrak a_1)$
does not identify that conductor with the original larger deleted set.

\subsection{Removing the row-dependent functional-equation weights}
First take finite-order $\nu$. All characters then have trivial infinity
type at the complex place. The primitive completion may be normalized as
\[
 \Lambda(s,\rho)=Q_\rho^{s/2}\gamma_K(s)L(s,\rho),\qquad
 \gamma_K(s)=2(2\pi)^{-s}\Gamma(s),\qquad
 \Lambda(s,\rho)=\epsilon_\rho\Lambda(1-s,\bar\rho),
 \quad |\epsilon_\rho|=1.
\]
For nonprincipal rows the completed function is entire. Apply this
functional equation to the Mellin integral with even auxiliary factor
$G(z)=e^{z^2}$ and denominator $z$, shifting across $z=0$. At
$s=1/2+it$ this gives two weighted Dirichlet series, truncated below to
polynomials; the dual prefactor has modulus one. Principal rows are excluded from
this pole-free step and treated separately below.

For fixed $\mathfrak a_2$, let $X=H/(N\mathfrak a_2)^6\ge1$ and
$T=1+|t|$. An explicit normalized kernel is
\begin{equation}\label{mm:afe-kernel}
 V_t(x)=\frac1{2\pi i}\int_{(\eta)}
 e^{z^2}\frac{\Gamma(s+z)}{\Gamma(s)}T^{-z}x^{-z}\,\frac{dz}{z},
 \qquad x=\frac{2\pi N\mathfrak b}{\sqrt{Q_\rho}T},
\end{equation}
with $\eta>0$ small. Its factors combine to
$(\sqrt{Q_\rho}/(2\pi N\mathfrak b))^z$ in the actual completed-L
kernel; there is no extra factor of $T$. The denominator $z$ supplies
the residue at zero.
On $z=\eta+iy$, Stirling and $|t+y|\le T(1+|y|)$ give
\[
 \left|\frac{\Gamma(1/2+\eta+i(t+y))}
                 {\Gamma(1/2+it)T^{\eta+iy}}\right|
 \ll_\eta (1+|y|)^{C_\eta}e^{\pi|y|/2}.
\]
The factor $e^{-y^2}$ makes this an integrable majorant uniformly in
$t$. The analogous bound holds on every fixed positive shifted line
and for the dual gamma kernel. Moving far enough to the right gives
arbitrary power decay beyond a common total-ideal cutoff
\begin{equation}\label{mm:common-cutoff}
 B'=X^{1/2}T(HT)^\eta,
\end{equation}
up to fixed $\xi,S,K$ constants. For any requested $A>0$, a sufficiently
large positive kernel shift, depending also on $\eta$, makes the
truncation error $O((HT)^{-A})$ uniformly in all rows in this block.
Ideal counting $O_K(Y)$ bounds the absolute tails. Such errors are
negligible after the finite row sums by taking $A$ larger if needed.

Insert~\eqref{mm:afe-kernel} on $\operatorname{Re}z=\eta$.
The factor $(\sqrt{Q_\rho}T)^z$ is independent of the column and is a
row multiplier of modulus at most $(CXT^2)^{\eta/2}$.
Its imaginary part is a phase. The remaining column coefficient is
$(N\mathfrak b)^{-1/2-\eta-it-iy}$, apart from fixed constants.
Minkowski or Cauchy with the integrable gamma majorant therefore removes
the row-dependent weight at cost $(HT)^{O(\eta)}$ in the squared norm.
No independence of the root numbers is assumed or needed.

For each column use the unique decomposition
\[
 \mathfrak b=\mathfrak b_S\mathfrak b_1\mathfrak b_2,
 \quad \mathfrak b_S\ S\hbox{-supported},\quad
 \mathfrak b_1\ \hbox{squarefree},\quad
 \mathfrak b_2\ \hbox{squarefull},\quad
 (\mathfrak b_1,\mathfrak b_2)=1,
\]
where $\mathfrak b_1,\mathfrak b_2\in I(S)$; a prime of outside-$S$
valuation one goes in $\mathfrak b_1$, and a prime of valuation at least
two goes with its full exponent in $\mathfrak b_2$. For fixed
$\mathfrak b_S,\mathfrak b_2$ retain
$\rho_{\mathfrak a}(\mathfrak b_S)$ as the original primitive row
value. On the good factors the exact canonical identity applies.
The remaining factors involving $\mathfrak b_2$, including
$\vartheta_{\mathfrak a_0}(\mathfrak b_2)$, are bounded row
multipliers. Their moduli are at most one; only then are they discarded
in an inequality. In particular the value at $\mathfrak b_S$ is not
replaced by the product of zero-extended primitive arrays.

For the $\mathfrak b_1$ polynomial the coefficient sequence is
independent of $\mathfrak a_0$; it includes the indicator
$(\mathfrak b_1,\mathfrak b_2)=1$ and the common cutoff
$N\mathfrak b_1\le B'/N(\mathfrak b_S\mathfrak b_2)$.
Padding it by zeros up to $B'$ is legitimate. Its squared norm is at most
$\zeta_K(1+2\eta)$. Minkowski over $\mathfrak b_S,\mathfrak b_2$
is legitimate because
\[
 \sum_{\mathfrak b_S,\mathfrak b_2}
       N(\mathfrak b_S\mathfrak b_2)^{-1/2-\eta}<\infty.
\]
The finite Euler product handles $\mathfrak b_S$; the squarefull series
converges above $1/2$. Apply~\eqref{mm:ls} to each such polynomial and
to the conjugate dual family. For fixed $\mathfrak a_1,\mathfrak a_2$
the resulting energy is
\begin{equation}\label{mm:block-energy}
 (HT)^{O(\eta)}\{M+B'+(MB')^{2/3}\},\qquad
 M=X/N\mathfrak a_1.
\end{equation}
The additional condition $(\mathfrak a_0,\mathfrak a_1)=1$ simply
restricts the positive row sum. The large-sieve small exponent can be
chosen sufficiently small along with $\eta$ and absorbed in the
indicated small loss.

\subsection{All-ideal summation and principal rows}
Every squarefull ideal $\mathfrak c$ has a unique expression
$\mathfrak d^2\mathfrak e^3$ with $\mathfrak e$ squarefree.
This follows primewise by using exponent zero or one in $\mathfrak e$
according to the parity of an exponent at least two. Therefore
\[
 \#\{\mathfrak c\ \mathrm{squarefull}:N\mathfrak c\le X\}
 \ll_K X^{1/2},\qquad
 \sum_{\mathfrak c\ \mathrm{squarefull}}(N\mathfrak c)^{-s}<\infty
 \quad(s>1/2).
\]
For the count, sum
$O((X/(N\mathfrak e)^3)^{1/2})$ and use
$\zeta_K(3/2)<\infty$. Summing~\eqref{mm:block-energy} over
$\mathfrak a_1$ gives
\[
 (HT)^{O(\eta)}
 \{X+B'\sqrt X+(XB')^{2/3}\}
 \ll (HT)^{O(\eta)}XT,
\]
since the squarefull series converges at $1$ and $2/3$,
$B'=\sqrt X\,T(HT)^\eta$, and $T^{2/3}\le T$ for $T\ge1$.
The $\mathfrak a_2$ sum is bounded by
$H\sum_{\mathfrak a_2}(N\mathfrak a_2)^{-6}=H\zeta_K(6)$.
Choosing all internal losses small for the requested final
$\varepsilon$ proves the fixed-twist all-ideal central moment.

There are $O_K(H)$ ideal rows in total. On any row where the primitive
character is principal its $L$-function is $\zeta_K$, and ordinary
scalar convexity gives
$|\zeta_K(1/2+it)|^2\ll_\varepsilon T^{1+\varepsilon}$.
This treats those rows without a pole-free approximate functional
equation. Finally~\eqref{mm:literal-sum} proves~\eqref{mm:central}
for every actual nonzero Eisenstein element.

\subsection{Poisson transfer with an explicit constant}
For a nonprincipal row define
\[
 f_\psi(s)=\frac{L(s,\psi)}{(s+1)^{5/4}},
\]
using the holomorphic logarithm on $\operatorname{Re}(s+1)>0$.
On a principal row additionally multiply by $(s-1)/(s+1)$.
Analytic continuation and ordinary degree-two convexity, used separately
at each fixed conductor, show that $f_\psi$ is bounded analytic on the
right half-plane $\operatorname{Re}s>1/2$, with continuous boundary
and decay at infinity. The principal pole is removed before applying
Poisson representation to its real and imaginary parts. Jensen gives
\[
 |f_\psi(1/2+\delta+it)|^2
 \le\int_{\mathbb R}P_\delta(t-v)
             |f_\psi(1/2+iv)|^2\,dv.
\]
The exact row family is finite for fixed $H$, so its sum interchanges
with the integral without a limiting-family issue. The boundary
pole-removal factor has modulus at most one.

Take $0<\varepsilon_0\le1/4$ and
$\alpha=3/2-\varepsilon_0\in[5/4,3/2)$. Since
$|3/2+iv|\ge(1+|v|)/\sqrt2$, the central squared moment of the
$f_\psi$ is bounded by
$3C_0H^{1+\varepsilon_0}w_\alpha(v)$, where
$w_\alpha(v)=(1+|v|)^{-\alpha}$.
For $|t|\ge2$, split $|v-t|\le|t|/2$ from its complement. The
near contribution to $P_\delta*w_\alpha$ is at most
$2^\alpha w_\alpha(t)<3w_\alpha(t)$. In the far region use
$P_\delta(t-v)\le4\delta/(\pi t^2)$ and
$\int w_\alpha=2/(\alpha-1)$, giving at most
\[
 \frac{2\delta3^\alpha}{\pi(\alpha-1)}w_\alpha(t)
 \le8w_\alpha(t).
\]
Here $(1+|t|)^\alpha/t^2\le3^\alpha/4$ for $|t|\ge2$.
For $|t|\le2$, the convolution is at most one and hence at most
$3^\alpha w_\alpha(t)<6w_\alpha(t)$. Thus
$P_\delta*w_\alpha\le11w_\alpha$ uniformly for $0<\delta<1/2$.
Restoring the denominator costs less than
$6(1+|t|)^{5/2}$, and the inverse principal factor is bounded by
$(3/2+\delta)/(1/2-\delta)$. The product $11\cdot3\cdot6=198<200$
proves~\eqref{mm:poisson-constant}. The constant is uniform as
$\delta\downarrow0$, but not as the contour approaches the principal
pole line $\operatorname{Re}s=1$. No zero-free hypothesis appears.

\subsection{The separately scoped unitary infinity type}
Let $\nu(\mathfrak a)=\nu_0(\mathfrak a)(N\mathfrak a)^{-i\tau}$
be fixed and unitary, with angular parameter $m$ and gamma factor
$\gamma_\nu$ stated in Section~\ref{mm:core}. The external primitive
functional equation is explicitly
\[
 Q_\rho^{s/2}\gamma_\nu(s)L(s,\rho)
 =\epsilon_\rho Q_\rho^{(1-s)/2}
                 \gamma_{\bar\nu}(1-s)L(1-s,\bar\rho),
 \qquad |\epsilon_\rho|=1.
\]
A constant norm-twist phase in another completion convention is absorbed
in $\epsilon_\rho$. All finite Kummer twists share this infinity type.
Replace $T$ in the kernel by $T_\nu=1+|t+\tau|+|m|$.
The numerator gamma ratio now has real base $1/2+|m|/2>0$ and
vertical parameter $t+\tau$; its conjugate dual has parameter
$-(t+\tau)$. Stirling gives the same integrable Gaussian-majorized
kernel on each fixed shifted line. Constants can depend on $m,\tau$;
$T_\nu$ is comparable to $T$ for this one fixed $\nu$.
The finite-prime conductor analysis, modulus-one coefficient twist,
all-ideal decomposition and all original primitive row values remain
unchanged.

A pure norm row has $L(s,\rho)=\zeta_K(s+i\tau)$; exclude it from
the pole-free approximate functional equation and use scalar convexity
at $t+\tau$. In the Poisson step cancel its actual pole by
\[
 h_\tau(s)=\frac{s-1+i\tau}{s+1+i\tau}.
\]
This factor is analytic in the relevant right half-plane. On its
boundary $|h_\tau|\le1$, while inside
$|h_\tau^{-1}|\le(3/2+\delta)/(1/2-\delta)$.
Hence the same transfer factor $200$ applies to the now
$\nu$-dependent central constant. If the opposite norm convention is
chosen, replace $\tau$ by $-\tau$ in the gamma factor and pole alike.
This proves the separate fixed-unitary moment assertion; it does not
claim uniformity in a family of infinity types or replace the
finite-order principal-residue expression below.

\subsection{Keeping the full type-I mask during transfer}
For $\sigma=1/2+\delta$, define $L^T$ by removing every prime of
$T(u,d)$, not just conductor primes. The identity~\eqref{mm:masked}
is equivalently
\[
 F_{u,d}(s)=\frac{L(s,\psi_u)}{L^{T(u,d)}(2s,\psi_{u,2})}
 \prod_{p\in T(u,d),\ p\nmid\mathfrak f_u}
                         (1-\psi_u(p)(Np)^{-s}).
\]
Absolute convergence at $2\sigma>1$ bounds the reciprocal incomplete
product by $\zeta_K(2\sigma)$, without requiring any zero-free theorem
in the critical strip. For every fixed $\eta>0$,
\[
 \prod_{p\in T}(1+(Np)^{-\sigma})
 \ll_{\sigma,\eta}(N\operatorname{rad}T)^\eta,
 \qquad N\operatorname{rad}T\le C_S Nu\,Nd.
\]
Indeed $1+q^{-\sigma}\le q^\eta$ beyond a fixed threshold, and only
finitely many prime ideals lie below it. Therefore
\begin{equation}\label{mm:mask-bound}
 |F_{u,d}(\sigma+it)|
 \ll_{\nu,S,\delta,\eta}(Nu\,Nd)^\eta
                              |L(\sigma+it,\psi_u)|.
\end{equation}
This is a bounded inequality after the exact mask identity, not an
identity that removes original zeros. For example, with principal $\nu$
and $u=\pi^3$ or $\pi^6$ at $N\pi=7$, the primitive local ratio,
restoration and full masked factor at $s=1$ are respectively
\[
 \begin{array}{c|ccc}
  v_\pi(u)&\hbox{primitive ratio}&\hbox{restoration}&\hbox{masked factor}\\\hline
 3&48/49&49/48&1\\
 6&8/7&7/8&1
 \end{array}
\]
The squared primitive character in the first row is unramified at
$\pi$ although the original primitive character is ramified there.
Squaring a zero sentinel would give the wrong denominator.

For finite-order $\nu$, the normalized type-I piece is the retained
$d$ sum of the smoothed squarefree series at $Y=D/Nd$, divided by
$\sqrt D$. Its possible pole is only at $s=1$, on principal rows.
Contour shifting with the full mask gives $I_u=P_u+E_u$, with $E_u$
exactly~\eqref{mm:remainder} and
\begin{align*}
 P_u&=\mathbf1_{\psi_u=1}\sqrt D\,\widehat W(1)
 \sum_{d\ \mathrm{squarefree}\atop(d,S)=1}
 \frac{\beta_U(d)\nu(d)\chi_d(u)}{Nd}\rho_{T(u,d)},\\
 \rho_T&=\frac{\operatorname*{res}_{s=1}\zeta_K(s)}{\zeta_K(2)}
                   \prod_{p\in T}(1+(Np)^{-1})^{-1}.
\end{align*}
This does not impose positivity on $P_u$. Since $\beta_U$ is supported
on $Nd\le U^2\le D$, every inner scale $Y$ is at least one.
Four integrations by parts give
$|\widehat W(\sigma+it)|\ll_\sigma\mathcal M_4(W)(1+|t|)^{-4}$,
which justifies the shifted integrals and their horizontal limits with
the classical Hecke growth estimates.

Apply Minkowski to~\eqref{mm:remainder} in $\ell^2$ of the actual
$u$ rows, then use~\eqref{mm:mask-bound} and the off-central moment.
For sufficiently small internal losses this bounds $\|E\|_2$ by
\[
 C D^\delta H^{1/2+O(\eta)}
 \sum_d\frac{|\beta_U(d)|}{(Nd)^{\sigma-\eta}}
 \int_{\mathbb R}|\widehat W(\sigma+it)|
                     (1+|t|)^{(1+O(\eta))/2}\,dt.
\]
All original $\nu(d)\chi_d(u)$ have modulus at most one and remain
inside the exact expression before this inequality. The actual truncated
factorization gives, without replacing $\beta_U$ by arbitrary coefficients,
\[
 \sum_d\frac{|\beta_U(d)|}{(Nd)^{\sigma-\eta}}
 \le\left(\sum_{Na\le U}(Na)^{-\sigma+\eta}\right)^2
 \ll_{\delta,\eta}U^{1-2\delta+2\eta}.
\]
Choose $\eta<\delta$ if necessary and use the fixed-field ideal count.
The height integral is bounded by $C\mathcal M_4(W)$. Squaring and
absorbing the small losses in $(HD)^\varepsilon$, using $U^2\le D$,
proves~\eqref{mm:energy}. Thus the formerly conditional numerator
moment input in this residual transfer is discharged, while its
positive-power cutoff loss is retained explicitly.

\subsection{What remains open and what has been verified}
The new remainder envelope improves the old scalar-convexity $H$ power
$3/2-\delta+\varepsilon$ to $1+\varepsilon$. At fixed $U$ one may
choose $\delta$ for a requested small-loss exponent; for $U=D^a$ with
fixed $a>0$, the remaining excess tends to $2a$ as $\delta\downarrow0$.
Neither statement controls $P_u$, the other signed piece, nor their
covariance. In the original Vaughan identity the short $2a_U$ term
must be restored at smaller scales where its support can meet the
annulus. Uniform $\mathcal M_4$ bounds do not justify interchanging a
profile supremum with the row moment or assuming $U^2\le D'$ after a
scale $D'$ is reduced.

The independent prose audit is recorded in
\path{planning/analytic-oracle/arithmetic-review/adversarial-afm-review-20261104.md}.
The source-hashed finite checker tests valuation/support decompositions,
squarefull coordinates, height exponents and rational Poisson envelopes.
Its finite success does not prove reciprocity, analytic continuation,
the large sieve, the approximate functional equation or the moments.
No new Lean declarations accompany these analytic theorems. Existing
Lean local Euler identities are a separately limited verification layer.
The independent positive arithmetic determinant comparison, original
signed low-gcd short-row estimate, all required contour uniformities and
GRH remain open; Dirichlet, Hecke, Dedekind and general Artin scopes are
not identified with one another.

\paragraph{Pinned primary sources and external inputs.}
The principal analytic family input is
\href{https://arxiv.org/pdf/1112.1650v1}{Blomer--Goldmakher--Louvel,
arXiv:1112.1650v1}, Theorem~1.3, Lemma~2.1, Appendix~II and Section~5.
Our fixed-$\nu$ proof changes column coefficients and redoes the
functional-equation reduction; it does not claim the published
corollary directly states every fixed twist.
For comparison,
\href{https://arxiv.org/pdf/2201.01885v2}{Gao--Zhao,
arXiv:2201.01885v2}, Lemma~2.9 and Section~3.2, uses the sextic
large-sieve/moment interface. Its GRH-conditional central-density theorem
is not invoked. Abelian Kummer identification uses the classical
reciprocity theorem, with the project's pinned Milne \emph{Class Field
Theory}, version~4.03 (6 August 2020), Theorem~0.8 and Chapter~I,
Section~1, a reference:
\href{https://www.jmilne.org/math/CourseNotes/CFT.pdf}{Milne's notes}.
The primitive degree-one Hecke continuation and completed functional
equation are explicit classical external inputs; the project's pinned
\href{https://u.cs.biu.ac.il/~reznikov/courses/kudla-1.pdf}{Kudla notes
on Tate's thesis}, Theorem~4.3 and Corollary~4.4, supply a reference.
Their sketch and incomplete displayed epsilon formulas are not used as
independently certified constants. Standard Stirling, Mellin inversion,
Poisson representation and scalar convexity complete the external
analytic dependencies. None is replaced by a finite diagnostic or a
GRH-equivalent assumption.

\section{Recombining the actual weighted cutoff}
\label{rc:core}
Keep the finite-order Eisenstein family and original masks of
Section~\ref{mm:core}. For $U\ge1$ and a squarefree ideal $m$, put
\[
 A_U(m)=\sum_{d\mid m}\beta_U(d)
 =\sum_{ab\mid m\atop(a,b)=1,\ Na,Nb\le U}\mu(a)\mu(b).
\]
Every summation variable here is an integral ideal. The last formula
uses squarefreeness of $m$; an unconstrained square of a M\"obius
polynomial would also include forbidden overlapping supports.
\begin{theorem}[Exact recombination and an absolute-half-plane rate]
\label{rc:recombination}
Write $T_0=S\cup\operatorname{supp}(u)$ and
$J_{u,U}(s)=\sum_{d\ \mathrm{squarefree},\ (d,S)=1}
\beta_U(d)\nu(d)\chi_d(u)(Nd)^{-s}F_{u,d}(s)$.
For $\operatorname{Re}s>1$,
\begin{align*}
 J_{u,U}(s)&=\sum_{m\ \mathrm{squarefree},\ (m,T_0)=1}
                   A_U(m)\psi_u(m)(Nm)^{-s},\\
 |A_U(m)|&\le3^{\omega(m)},\qquad
 A_U(m)=\mu(m)\quad(Nm\le U).
\end{align*}
For every $\kappa>0$ and $0<\epsilon<\kappa$, uniformly for
$\operatorname{Re}s\ge1+\kappa$, all original rows and all $U\ge1$,
\[
 \left|J_{u,U}(s)-\frac1{L^{T_0}(s,\psi_u)}\right|
 \le U^{-\kappa+\epsilon}
 \{\zeta_K(1+\epsilon)^3+\zeta_K(1+\epsilon)\}.
\]
Here $L^{T_0}$ removes every prime of $T_0$, including conductor-drop
primes; the bound is in the absolute-convergence half-plane only.
\end{theorem}
\begin{proof}
On a surviving term $m=dc$, the two squarefree ideals $c,d$ are
coprime, and both avoid $T_0$. The original product
$\nu(d)\chi_d(u)\psi_u(c)$ is then $\psi_u(m)$. Grouping by $m$
gives $A_U$. Each prime of $m$ has three possible positions in
$(a,b,m/(ab))$, so the absolute sum is at most $3^{\omega(m)}$.
When $Nm\le U$ all divisors fit both cutoffs; the primewise sum is
$1-1-1=-1$, giving $\mu(m)$.

The same bound justifies absolute convergence and regrouping for
$\operatorname{Re}s>1$, because on squarefree ideals
$3^{\omega(m)}=\tau_3(m)$ and
$\sum_m\tau_3(m)(Nm)^{-v}=\zeta_K(v)^3$ for $v>1$,
where this latter sum runs over all nonzero integral ideals.
Subtract the absolutely convergent M\"obius series for $1/L^{T_0}$.
Only $Nm>U$ remain. Rankin's inequality with exponent
$\eta=\kappa-\epsilon$ bounds their absolute sum by
\[
 U^{-\eta}\sum_m\{\tau_3(m)+1\}
                  (Nm)^{-1-\epsilon},
\]
which proves the stated constant independently of deleted primes,
row multiplicities or character phases.
\end{proof}

For $W$ supported in $[1,2]$, the same exact regrouping gives the
original normalized type-I polynomial
\[
 I_u=D^{-1/2}\sum_{m\ \mathrm{squarefree},\ (m,S)=1}
 A_U(m)\nu(m)\chi_m(u)W(Nm/D).
\]
This formula retains all original zero sentinels. It permits a direct
family large sieve before any absolute $d$ sum; its quantitative moment
is developed separately below. The displayed absolute-half-plane rate
alone is not that critical-strip moment.

\section{Euler packaging, coprimality and the analytic boundary}
\label{rc:details}
For $p\notin T_0$, set $z_p=\psi_u(p)(Np)^{-s}$ and
$t_p=z_p/(1+z_p)$. On $\operatorname{Re}s>1$, deleting an additional
surviving prime divides the squarefree Euler product by $1+z_p$.
Consequently the exact retained expression is
\[
 J_{u,U}(s)=F_{u,1}(s)B_{u,U}(s),\qquad
 B_{u,U}(s)=\sum_{d\ \mathrm{squarefree},\ (d,T_0)=1}
                 \beta_U(d)\prod_{p\mid d}t_p.
\]
Original zeros enforce $(d,u)=1$ even when a sixth power drops from
its primitive conductor. This factorization cannot restore a forbidden
$d$ row by replacing a zero sentinel with an unramified primitive value.
The equality continues wherever both sides are holomorphic.

Without the cutoffs, $\beta_\infty(d)=\mu(d)2^{\omega(d)}$ on
squarefree ideals. Absolute convergence justifies
\[
 B_{u,\infty}=\prod_{p\notin T_0}(1-2t_p)
 =\prod_{p\notin T_0}\frac{1-z_p}{1+z_p},\qquad
 F_{u,1}B_{u,\infty}=\prod_{p\notin T_0}(1-z_p)
                   =\frac1{L^{T_0}(s,\psi_u)}.
\]
Theorem~\ref{rc:recombination} gives a quantitative rate rather than an
unsupported exchange of this limit with a critical-strip contour.

\paragraph{A local counterexample to the naive square.}
In a one-prime projection at a good Eisenstein prime of norm $7$, take
the principal local phase and $s=1$. Then $z=1/7$, $t=1/8$ and,
for $U\ge7$, the exact polynomial is $B_U=1-2t=3/4$.
The unconstrained square $(1-t)^2=49/64$ differs by $1/64$:
it admits the forbidden pair $a=b=p$. After multiplying by
$F=1+z=8/7$ it gives $7/8$ instead of the correct $6/7$.
This is an actual local Euler calculation, not an infinite-series or
mean-square counterexample. It identifies the exact coprimality condition
that a proposed mollifier-square simplification must preserve.

\paragraph{The exact intermediate coefficient band.}
For a squarefree $m$ with $U<Nm\le U^2$, one has $A_U(m)=-\mu(m)$.
Indeed the unrestricted disjoint-assignment sum is $\mu(m)$. If $Na>U$,
every $b\mid m/a$ automatically has $Nb\le Nm/Na<U$, and
$\sum_{b\mid m/a}\mu(b)$ vanishes unless $a=m$. Hence the excluded
$a$ assignments contribute exactly $\mu(m)$; the excluded $b$
assignments contribute another $\mu(m)$. They cannot overlap because
$NaNb>U^2\ge Nm$ contradicts $ab\mid m$. Subtraction proves the claim.
Together with the small-norm band, this gives $A_U=\mu$ for $Nm\le U$
and $A_U=-\mu$ for $U<Nm\le U^2$. The original smoothed columns have
$Nm\ge D\ge U^2$, so the new band does not justify replacing their
coefficients by $-\mu(m)$: generally they lie beyond it.

\begin{theorem}[A stronger cutoff convergence would require zero-freeness]
\label{rc:boundary}
Fix a nonprincipal primitive $\psi_u$. If these actual $J_{u,U}$
converge locally uniformly on $\operatorname{Re}s>1/2$ as
$U\longrightarrow\infty$, then $L(s,\psi_u)$ has no zero there.
\end{theorem}
\begin{proof}
Classical abelian Hecke continuation makes the numerator of
$F_{u,1}=L^{T_0}(s,\psi_u)/L^{T_0}(2s,\psi_{u,2})$ holomorphic.
Its denominator is nonzero on $\operatorname{Re}s>1/2$ by the
absolute Euler product at $2s$. Each $B_{u,U}$ is a finite sum with
nonzero denominators $1+z_p$ on $\operatorname{Re}s>0$.
Thus every $J_{u,U}$ is holomorphic on the asserted half-plane.
A locally uniform limit $G$ is holomorphic there. On
$\operatorname{Re}s>1$, the already proved convergence gives
$L^{T_0}(s,\psi_u)G(s)=1$. The identity theorem extends this equality
to the connected half-plane, precluding a zero of $L^{T_0}$.
The finitely deleted factors $1-\psi_u(p)(Np)^{-s}$ do not vanish
when $\operatorname{Re}s>0$, so the same conclusion holds for the
primitive $L$-function.
\end{proof}
This is an implication, not a proof that the stronger convergence holds,
and it does not assert a converse. In particular, a family-averaged energy
bound is not silently promoted to locally uniform convergence for each
individual character. Principal rows have a pole at $1$ and were excluded
from this proposition. General Artin holomorphy is not an input.

\paragraph{Reproducible exact diagnostics.}
The new source \path{examples/arithmetic_cutoff/exact_recombination.py}
uses actual split/inert prime ideals and raw sextic symbols in
$\mathbb Z[\omega]$, with exact rational arithmetic. It checks the direct
masked sum, coefficient regrouping, Euler packaging, both finite full-cutoff
limits and original zeros, including third/sixth-power rows. Its overlapping-
support negative controls are not tests of an infinite large sieve or of
zero location. The analytic convergence argument above uses ordinary ideal
Euler products and the explicitly stated Hecke continuation theorem; no
new Lean declaration or complete GRH proof is claimed.

% Original masked cutoff family; no new Lean or GRH theorem.
\section{A cutoff-uniform envelope for the original masked remainder}
\label{wc:core}
Fix $K=\mathbb Q(\omega)$, $\mathcal O=\mathbb Z[\omega]$, a finite-order
Hecke character $\nu$, and a finite set $S$ containing the primes above $6$
and the conductor of $\nu$. All unspecified ideal sums below are over
nonzero integral ideals. Let $\chi_m(u)=(u/m)_6$ be the original symbol,
zero when $(m,u)>1$, and put
\[
 \beta_U(d)=\sum_{ab=d\atop Na,Nb\le U}\mu(a)\mu(b),\qquad
 A_U(m)=\sum_{d\mid m}\beta_U(d).
\]
The columns $d,c,m$ are squarefree and prime to $S$ where indicated.
Retain the primitive characters $\psi_u=\operatorname{prim}(\nu\kappa_u)$,
$\psi_{u,2}=\operatorname{prim}(\psi_u^2)$, and the full deleted set
$T(u,d)=S\cup\operatorname{supp}(u)\cup\operatorname{supp}(d)$.
For $H,D,U\ge1$, $U^2\le D$, and $W\in C_c^\infty((0,\infty))$
supported in $[1,2]$, the original normalized type-I piece is
\begin{equation}\label{wc:polynomial}
 I_u=D^{-1/2}\sum_{m\ \mathrm{squarefree}\atop(m,S)=1}
     A_U(m)\nu(m)\chi_m(u)W(Nm/D).
\end{equation}
Its exact principal residue is
\begin{align}\label{wc:principal}
 P_u&=\mathbf1_{\psi_u=1}\sqrt D\,\widehat W(1)
 \sum_{d\ \mathrm{squarefree}\atop(d,S)=1}
 \frac{\beta_U(d)\nu(d)\chi_d(u)}{Nd}\rho_{T(u,d)},\notag\\
 \rho_T&=\frac{\operatorname*{res}_{s=1}\zeta_K(s)}{\zeta_K(2)}
               \prod_{p\in T}(1+(Np)^{-1})^{-1}.
\end{align}
Define $E_u=I_u-P_u$; this is the same residue-subtracted remainder as the
fully masked Mellin contour, with no cutoff limit or profile supremum.

\begin{theorem}[Uniformity in the actual truncated cutoff]
\label{wc:uniform}
For every $\varepsilon>0$,
\begin{equation}\label{wc:bound}
 \sum_{0<Nu\le H}|E_u|^2
 \ll_{\nu,S,\varepsilon}\|W\|_\infty^2(HD)^\varepsilon
       \{H+DH^{1/3}+(HD)^{2/3}\}.
\end{equation}
The sum is over every actual nonzero Eisenstein element, including units,
$S$-supported powers and sixth-power rows. The constant is independent of
$U$ under $U^2\le D$. This is an application of the classical sextic
large sieve with exact original coefficients, not a zero-location theorem.
\end{theorem}
\begin{proof}
On an absolutely convergent line, expand the restored squarefree series
$F_{u,d}$. Its condition $(c,T(u,d))=1$ includes $(c,d)=1$, and the
original symbols multiply. Grouping $m=cd$ gives~\eqref{wc:polynomial}
by Mellin inversion. Each prime of $m$ may belong to the first truncated
M\"obius divisor, to the second, or to neither. Consequently
$|A_U(m)|\le3^{\omega(m)}\ll_\eta(Nm)^\eta$, uniformly in $U$.
The length-$D$ coefficient sequence in~\eqref{wc:polynomial} therefore
has squared norm $\ll_\eta\|W\|_\infty^2D^\eta$.

The all-ideal row norm in Section~3 of
\href{https://arxiv.org/pdf/1112.1650v1}{Blomer--Goldmakher--Louvel,
arXiv:1112.1650v1} gives
\[
 B_2^1(X,D)\ll_\eta(XD)^\eta
                 \{X+DX^{1/3}+(XD)^{2/3}\}.
\]
For completeness, split $a=a_1a_2^2\cdots a_5^5a_6^6$ dyadically;
$a_1,\ldots,a_5$ are squarefree and pairwise coprime. If their dyadic
sizes are $A_j$, put $R=\prod_{j=1}^6A_j$ and $M=\max(A_1,A_2)$.
Pivot on the first or second variable, using the squarefree powered
sieve proved simultaneously in that paper's Section~3.2. The other
variables are bounded column multipliers. The block energy is at most
\[
 (XD)^\eta\{R+DR/M+D^{2/3}R/M^{1/3}\}\sum_m|\lambda_m|^2.
\]
The inequalities $R\le X$, $R^3\le XM^3$ and $R^3\le X^2M$
follow from $\prod_jA_j^j\le X$ and the pivot choice. Summing the
logarithmically many blocks proves the displayed bound.

Finite $S$-unit and ray-class partitions convert the original literal
symbol into the canonical column value $\vartheta_m(a)$ with fixed
modulus-one coefficient multipliers. On noncoprime pairs both values
are zero: squarefree $m$ has outside-$S$ conductor exactly $m$.
This orientation preserves sixth-power zeros. At most six elements
generate each outside-$S$ ideal times an $S$-supported ideal; the
$S$-supported sums converge at powers $1$, $1/3$ and $2/3$.
Thus~\eqref{wc:bound} holds first with $I_u$ in place of $E_u$.

A principal $\psi_u$ requires every outside-$S$ valuation of $u$ to
be divisible by $6$. There are $O_{\nu,S}(H^{1/6})$ such literal
rows. Equation~\eqref{wc:principal}, $|\rho_T|\ll_K1$, and the actual
coefficient sum
\[
 \sum_d\frac{|\beta_U(d)|}{Nd}
 \le\left(\sum_{Na\le U}\frac1{Na}\right)^2\ll\log^2(2U)
\]
give $\sum|P_u|^2\ll DH^{1/6}\|W\|_\infty^2\log^4(2U)$.
This is absorbed by the right side of~\eqref{wc:bound}, using
$U^2\le D$ and a smaller internal exponent. Finally
$|E_u|^2\le2|I_u|^2+2|P_u|^2$. The extended proof specifies the
canonical-family transfer and the powered-sieve input precisely.
\end{proof}

Combining this with the previously proved fixed-twist AFM transfer yields
\begin{equation}\label{wc:minimum}
 \sum_{0<Nu\le H}|E_u|^2\ll_{\nu,S,\delta,\varepsilon}
 \mathcal M_4(W)^2(HD)^\varepsilon
 \min\{D^{2\delta}U^{2-4\delta}H,\ H+DH^{1/3}+(HD)^{2/3}\}.
\end{equation}
At $H=D^{3/4}$ and $U=D^a$, $0\le a\le1/2$, the excess over $H$
is at most $\min\{2\delta+(2-4\delta)a,1/2\}$, up to arbitrarily
small losses. The new envelope improves the old one for $a>1/4$ as
$\delta\downarrow0$, but still leaves an exponent loss of $1/2$.
It does not settle the original signed short-row bound, type-II covariance,
positive arithmetic determinant comparison or GRH. The principal-residue
bound used here is sufficient for this envelope, not the stronger target
needed to reassemble the full signed argument.

\section{All-ideal reciprocity and the surviving block obstruction}
\label{wc:details}
The proof uses no new zero-free input. Its external analytic ingredient
is the pinned BGL large sieve and the powered estimates established in
its proof; elementary ideal counting and the already restored Mellin
identity supply the remaining interfaces. All constants may depend on
fixed $K$, $S$, $\nu$ and declared small exponents, never on $U$.

\subsection{Exact recombination retains the primitive corrections}
On $\operatorname{Re}s>1$, start from the full defining series, with
its original zero masks:
\[
 F_{u,d}(s)=\sum_{c\ \mathrm{squarefree}\atop(c,T(u,d))=1}
                 \nu(c)\chi_c(u)(Nc)^{-s}.
\]
Its primitive ratio involves $L(s,\psi_u)$, the independently primitive
$L(2s,\psi_{u,2})$ and every deleted Euler factor. Expanding this
\emph{restored} series gives the exact equality
\begin{align}\label{wc:recombine}
 &\sum_{d\ \mathrm{squarefree}\atop(d,S)=1}
 \beta_U(d)\nu(d)\chi_d(u)(Nd)^{-s}F_{u,d}(s)\notag\\
 &\qquad=\sum_{m\ \mathrm{squarefree}\atop(m,S)=1}
               A_U(m)\nu(m)\chi_m(u)(Nm)^{-s}.
\end{align}
Indeed $(c,d)=1$ makes $m=cd$ squarefree, with $d\mid m$. If $m$
meets $u$, both original symbols retain their zero. Only on the
absolutely convergent line is this infinite series used. Mellin inversion
there gives the finite smoothed polynomial~\eqref{wc:polynomial}.
Nothing evaluates~\eqref{wc:recombine} as a Dirichlet series on
$\operatorname{Re}s=1/2+\delta$ or interchanges a cutoff limit with
continuation across unknown zeros.

For squarefree $m$, a factorization $ab=d\mid m$ assigns each prime
to $a$, to $b$, or to neither. The actual inequalities $Na,Nb\le U$
select a subset of these assignments. Thus
\[
 |A_U(m)|\le3^{\omega(m)}\ll_\eta(Nm)^\eta.
\]
The last divisor bound follows by choosing a fixed threshold so
$3\le(Np)^\eta$ beyond it, then absorbing the finitely many smaller
prime ideals in a constant. When $Nm\le U$ all assignments occur,
so $A_U(m)=\prod_{p\mid m}(1-1-1)=\mu(m)$; no such simplification
is made for the length-$D$ columns in general. With
$\lambda_m=D^{-1/2}A_U(m)\nu(m)W(Nm/D)$, ideal counting and a
smaller internal $\eta$ give
\begin{equation}\label{wc:l2}
 \sum_m|\lambda_m|^2\ll_\eta\|W\|_\infty^2D^\eta,
 \qquad \operatorname{supp}\lambda\subseteq\{D\le Nm\le2D\}.
\end{equation}

\subsection{The powered squarefree input and all-ideal pivot}
Let $\vartheta_a$ be BGL's canonical primitive sixth-order family.
Their Section~3 defines $B_1^j$ using squarefree row/column ideals and
kernel $\vartheta_a(m)^j$, while $B_2^j$ uses \emph{all} ideal rows
and kernel $\vartheta_m(a)^j$. Theorem~1.3 states the squarefree bound
for $j=1$; its Section~3.2 proves the estimates simultaneously for
$1\le j\le5$, because the recursion in Lemma~3.4 uses $j\ell$.
In particular the two inputs needed here are
\begin{equation}\label{wc:powered-input}
 B_1^j(M,D)\ll_\eta(MD)^\eta
                 \{M+D+(MD)^{2/3}\},\qquad j=1,2.
\end{equation}
To see the final step in the published recursion explicitly, its
$E_{4/3}$ estimate is $M^{4/3}+D+(MD)^{2/3}$.
If $M\le D$, its first term is at most $(MD)^{2/3}$; if $M>D$,
use reciprocity/duality (3.2) and the estimate with $M,D$ reversed.
The finite set of powers is treated together; no inference from a
fixed conductor count to an arbitrary primitive-character sieve occurs.

Write each ideal row uniquely as
\[
 a=a_1a_2^2a_3^3a_4^4a_5^5a_6^6,
\]
where $a_1,\ldots,a_5$ are pairwise-coprime squarefree ideals and
$a_6$ is arbitrary. For valuation $6q+r$, $0\le r\le5$, place the
prime in $a_r$ if $r>0$ and its power $q$ in $a_6$.
Dyadically let $A_j\le Na_j<2A_j$, with $A_j\ge1$, and put
$R=\prod_{j=1}^6A_j$. A nonempty block satisfies
$\prod_jA_j^j\le X$ when $Na\le X$.
Choose $\ell=1$ if $A_1\ge A_2$, and $\ell=2$ otherwise; put
$M=A_\ell$. Keep every other ideal fixed. Evaluation of the single
\emph{column} character gives
\[
 \vartheta_m(a)=\vartheta_m(a_\ell)^\ell
                  \vartheta_m(a/a_\ell^\ell).
\]
The second factor is a bounded column multiplier, including its original
zero if $m$ meets any fixed factor. It may depend on all the fixed
ideals but not on the pivot row. Pairwise-coprimality and the exact row
norm restriction merely select a subset of the positive pivot-row sum.
Apply~\eqref{wc:powered-input} to the first factor. There are
$O_K(R/M)$ choices of the other ideals by ideal counting. The block
energy is bounded by
\begin{equation}\label{wc:block}
 (XD)^\eta\{R+DR/M+D^{2/3}R/M^{1/3}\}\sum_m|\lambda_m|^2.
\end{equation}

All three monomial bounds used to sum these blocks are elementary.
First $R\le\prod_jA_j^j\le X$.
If $M=A_1$, then
\[
 (R/M)^3=A_2^3\cdots A_6^3
 \le A_1A_2^2A_3^3A_4^4A_5^5A_6^6\le X,
\]
since $A_2\le A_1$ and $j\ge3$ for the remaining factors.
If $M=A_2$, use $A_1\le A_2$ in the same comparison.
Finally,
\[
 \frac{R^3}{M(\prod_jA_j^j)^2}
 =\frac{A_1}{MA_2A_3^3A_4^5A_5^7A_6^9}\le1.
\]
Thus $R/M\le X^{1/3}$ and $R/M^{1/3}\le X^{2/3}$.
Insert these in~\eqref{wc:block} and sum the logarithmically many
blocks. This proves the all-ideal estimate used in Theorem~\ref{wc:uniform}.
It is precisely the mechanism of BGL Lemma~3.4 and (3.6)--(3.9),
with the final powered bound inserted. The simpler fixed-squarefull-row
argument would leave the weaker $D\sqrt X$ term.

\subsection{Literal rows, reciprocal orientation and principal count}
Let $a=(u)_S$ denote the outside-$S$ ideal. The finite $S$-unit
comparison gives $\kappa_u=\eta_u\vartheta_a$ as primitive
characters, with $\eta_u$ in a finite fixed set.
For $(a,m)=1$, BGL reciprocity therefore gives
\begin{equation}\label{wc:literal-reciprocity}
 \chi_m(u)=\eta_u(m)C([a],[m])\vartheta_m(a).
\end{equation}
The notation is deliberate: the left side is the project's original
$(u/m)_6$, not BGL's canonical character indexed by $m$.
When $(a,m)>1$, the left side is zero. The right side is also zero,
because squarefree $m$ has outside-$S$ primitive conductor exactly $m$.
Equation~\eqref{wc:literal-reciprocity} hence extends to every pair.
Using $\vartheta_a(m)$ on noncoprime pairs instead would be false:
primitive $\vartheta_{b^6}$ has restored primes of $b$, whereas the
original symbol must still vanish there.

Partition rows by $\eta_u$ and $[a]$, both finite. On each part the
factor $\eta_u(m)C([a],[m])$ is a fixed modulus-one column multiplier.
A positive row sum over that part is bounded by the full $B_2^1$ row
sum. For each $S$-supported ideal $b$ and outside-$S$ ideal $a$, at
most six actual elements generate $ab$. Thus the required literal energy
is bounded by the all-ideal norm summed over $b$, with $X=H/Nb$.
The geometric products
\[
 \sum_{b\ S\text{-supported}}(Nb)^{-q}
   =\prod_{p\in S}(1-(Np)^{-q})^{-1}<\infty,
 \qquad q=1,\ 1/3,\ 2/3,
\]
transfer the three terms without losing any unit or $S$-power row.
Together with~\eqref{wc:l2}, this proves the claimed $I$ energy.

If $\psi_u$ is principal, $\nu$ is unramified outside $S$, and tame
Kummer inertia at $v_p(u)$ has order $6/\gcd(6,v_p(u))$.
It must be trivial, so each such valuation is divisible by $6$.
Writing $(u)=ba_6^6$ with $b$ supported in $S$, ideal counting gives
\[
 \#\{0<Nu\le H:\psi_u=1\}
 \ll_K H^{1/6}\prod_{p\in S}(1-(Np)^{-1/6})^{-1}.
\]
This is only a necessary condition, sufficient for an upper count.
The exact residue formula~\eqref{wc:principal} and
$|\widehat W(1)|\le\|W\|_\infty$ then give the principal energy
bound used above. No sign or positivity of that residue is assumed.

\subsection{What the envelope still fails to prove}
The $DH^{1/3}$ term is still larger than $H(HD)^\varepsilon$ at
$H=D^{3/4}$. The monomial comparison itself is saturated by
$A_1=A_2=X^{1/3}$ and $A_3=\cdots=A_6=1$.
This proves sharpness of that \emph{particular comparison}, not an
energy lower bound for actual arithmetic coefficients. In particular,
no conclusion that the signed $A_U$ cannot cancel is justified.
The next load-bearing input must exploit those actual truncated M\"obius
coefficients inside the mixed moment, or establish the required signed
type-II minus type-I covariance. The smaller-scale short term, principal
contribution, full arithmetic determinant identification and GRH remain
separate obligations.

Exact finite checks use the split Eisenstein primes generated by
$3+\omega$, $4+\omega$, $5+2\omega$, $6+\omega$, of norms
$7,13,19,31$. Their actual residue symbols are computed in
$\mathbb Q(\zeta_6)$, retaining sixth-power zeros. Coefficient
recombination and dyadic inequalities pass; they check finite algebra
and bookkeeping, not the external analytic sieve or the contour theorem.
No new Lean declarations accompany this argument.

\section{High cutoffs, actual bilinear factors and a signed covariance test}\label{sb:core}
Retain the original Eisenstein family, fixed finite-order $\nu$, fixed $S$,
common $W\in C_c^\infty((0,\infty))$ supported in $[1,2]$, and every raw
shared-prime zero. Put
\[
 C_U(m)=\sum_{\substack{abc=m\\Na>U,\,Nb>U}}\mu(a)\mu(b),\qquad
 r_U(m)=\#\{(a,b):ab=m,\ Na,Nb>U\}.
\]
All factors are ideals, squarefree and pairwise coprime; on the active
annulus $A_U=C_U-\mu$. A coefficient identity is different from a row-energy bound.

\begin{proposition}[The high-cutoff window]\label{sb:window}
If $D>1$ and $2D/7\le U^2\le D$, then for every surviving squarefree
column $D\le Nm\le2D$,
\[
 C_U(m)=\mu(m)r_U(m),\qquad
 A_U(m)=\mu(m)(r_U(m)-1),\qquad |A_U(m)|\ge1.       
\]
For each fixed nonzero $W$, the normalized squared coefficient norm satisfies
\begin{equation}
 D^{-1}\sum_{\substack{m\ {\rm sf}\\(m,S)=1}}
 |A_U(m)W(Nm/D)|^2\gg_{W,S}1                     \label{sb:coefficient-lower}
\end{equation}
for large $D$, uniformly in this window. This is not a lower bound for row energy.
\end{proposition}
\begin{proof}
The least norm of a nonunit ideal outside $S$ is at least $7$. A long-long
term with $c\ne1$ would have $Nm>7U^2\ge2D$. Thus $ab=m$, including both
endpoints. The fixed-point-free swap $(a,b)\mapsto(b,a)$ makes $r_U$ even,
so $A_U$ is an odd nonzero integer. The last claim follows by counting
squarefree ideals on a fixed interval where $|W|$ is bounded below; the
constant and error are proved in the extended proof.
\end{proof}

The exact remaining polynomial is genuinely bilinear, with $(a,b)=1$,
$U<Na,Nb<2D/U$, $D\le Na\,Nb\le2D$, and
$x_{ab}(u)=x_a(u)x_b(u)$, including zeros. Both factors see the same row
$u=\varepsilon\alpha\beta^2$. Cauchy followed by the accepted powered BGL
bound gives only
\begin{equation}
 Q_{B_X}(C_U,C_U)\ll_{\nu,S,\epsilon}
 (DX)^\epsilon\|W\|_\infty^2
 \{\sqrt D\,X^2+DX+D^{5/6}X^{5/3}\}.              \label{sb:factor-bound}
\end{equation}
At $X\asymp D^{1/4}$ the exponents are $1,5/4,5/4$; the earlier regrouped
bound has $1/2,5/4,13/12$ and is stronger. No $BT_\sigma$ saving is obtained.
Multiplying two marginal operator norms does not supply a coupled fourth moment.

\paragraph{A finite test in the original family.}
There is no automatic nonnegative signed covariance, even when
$C_U(m)\mu(m)\ge0$. Take $\nu=1$, $S$ the primes above $2,3$, $D=400$,
$U=17$, and all $72$ literal rows of $B_7$. Choose the common nonnegative
smooth profile constructed in the extended proof, with values $10$ at $439/400$, $1$ at
$475/400$, and zero at every other integer norm divided by $400$.
There are exactly four active squarefree ideal columns, of norms
$439,439,475,475$. In that order,
\[
 A=(1,1,1,1),\quad C=(0,0,2,2),\quad
 \mu=(-1,-1,1,1),\qquad
 \Re G_{B_7}=
 \begin{pmatrix}72&0&0&18\\0&72&18&0\\0&18&72&0\\18&0&0&72\end{pmatrix}.
\]
Exact raw-symbol enumeration, independently reproducible from the ideal
and residue-field definitions, gives
\begin{equation}
 Q(C,\mu)=-\frac{27}{25},\qquad
 Q(C,C)=\frac{36}{25},\quad
 Q(\mu,\mu)=\frac{864}{25},\quad
 Q(A,A)=\frac{954}{25}.                          \label{sb:covariance-example}
\end{equation}
The real Gram matrix is positive definite, with eigenvalues $54,90$;
its positivity does not force a cross covariance to be positive.
This fixed finite example refutes that shortcut, not the asymptotic
$BT_\sigma$ target or the original signed mean-square estimate.

\paragraph{The two-cutoff contour remains coupled.}
For $T_0=S\cup\operatorname{supp}(u)$, $\psi=\psi_u$ and independently
primitive $\psi_2=\operatorname{prim}(\psi^2)$, put $q_p=\psi(p)(Np)^{-s}$,
$\alpha_p=(Np)^{-z}$, $\beta_p=(Np)^{-w}$ outside $T_0$. The actual
three-placement Euler factor is $1+q_p-q_p\alpha_p-q_p\beta_p$. Its exact
factorization is
\begin{equation}
 \mathcal F(s,z,w)=
 \frac{L^{T_0}(s,\psi)}{L^{T_0}(s+z,\psi)L^{T_0}(s+w,\psi)}
 \prod_{p\notin T_0}\left(1-
 \frac{q_p^2(1-\alpha_p)(1-\beta_p)}
 {(1-q_p\alpha_p)(1-q_p\beta_p)}\right).           \label{sb:perron-factor}
\end{equation}
Exact double Perron inversion is valid on positive auxiliary lines
$c,d>1-\Re s$ for $1/2<\Re s<1$, retaining the original principal pole
exclusion. These safe lines reproduce the previous $U^{2-4\delta}$
energy loss at $\Re s=1/2+\delta$, with the same common-profile derivative
and fixed-margin dependencies. The correction product alone does not
justify moving those lines through reciprocal-$L$ poles. In fact
$\mathcal F(s,z,0)=1/L^{T_0}(s+z,\psi)$ as a meromorphic identity.
No zero-free region, reciprocal moment, analytic saving, new Lean
formalization or GRH proof is supplied by this construction.

\section{Proof details for the high-window and two-cutoff tests}\label{sb:details}
\subsection{The coefficient lower bound and its precise limitation}
The element lattice $\mathbb Z[\omega]\subset\mathbb C$ has covolume
$\sqrt3/2$. Counting disks with boundary error $O(\sqrt x)$ and dividing
by the six units, using that this ring is a PID, gives
\[
 N_K(x)=\kappa x+O(\sqrt x),\qquad \kappa=\frac{\pi}{3\sqrt3}.
\]
Finite inclusion-exclusion at $S$ gives
$N_{K,S}(y)=\kappa y\prod_{p\in S}(1-(Np)^{-1})+O_S(\sqrt y)$.
Insert $1_{\rm sf}(m)=\sum_{d^2\mid m}\mu(d)$. Absolute convergence at
$2$, the tail $\sum_{Nd>\sqrt x}(Nd)^{-2}=O(x^{-1/2})$, and
$\sum_{Nd\le\sqrt x}(Nd)^{-1}=O(\log(2x))$ yield
\[
 N_{{\rm sf},S}(x)=c_Sx+O_S(\sqrt x\log(2x)),\qquad
 c_S=\frac{\kappa}{\zeta_K(2)}\prod_{p\in S}(1+(Np)^{-1})^{-1}>0.
\]
These tail estimates follow by partial summation from the displayed lattice
count. For a fixed nonzero $W$, choose $[a,b]\subset(1,2)$ with
$|W|\ge\eta>0$. Proposition~\ref{sb:window} gives, for large $D$,
\[
 D^{-1}\sum|A_U(m)W(Nm/D)|^2\ge\tfrac12\eta^2c_S(b-a).
\]
There is no uniform lower constant over all profiles with one specified
supremum norm. Nor can a coefficient lower bound be combined with an
operator upper bound to produce a row-energy lower bound: no lower singular
value has been established on the actual vector.

\subsection{What legal factor Cauchy actually proves}
For dyadic $Na\asymp A,Nb\asymp B$, retain the exact $(a,b)=1$ mask,
hyperbolic support and $W(Na\,Nb/D)$ in
\[
 T_a(u)=\sum_b\mu(b)\nu(b)W(Na\,Nb/D)x_b(u).
\]
Cauchy gives
\[
 |\sum_a\mu(a)\nu(a)x_a(u)T_a(u)|^2\ll A\sum_a|T_a(u)|^2.
\]
For fixed $a$, fixed second row factor and fixed unit/ray classes, the
$b$ coefficients are independent of the pivot row and have squared norm
$\ll B\|W\|_\infty^2$. The prior classical power-one/two BGL input has
$\Delta(X,B)=X+B+(XB)^{2/3}$. Summing the $O(A)$ fixed columns and
$O(X)$ other row factors gives
\[
 \mathcal E_{A,B}\ll_{\nu,S,\epsilon}(DX)^\epsilon\|W\|_\infty^2
 \frac{A^2B}{D}X\Delta(X,B)
 \ll (DX)^\epsilon\|W\|_\infty^2
 \{AX^2+DX+D^{2/3}A^{1/3}X^{5/3}\}.
\]
In the high window $A,B\asymp\sqrt D$, and only boundedly many dyadic
pairs occur. This proves~\eqref{sb:factor-bound}; exchanging factors
or row pivots leaves its $DX$ term. A product of marginal second-moment
bounds supplies no estimate for the two transforms on the same $(r,s)$.
Any useful saving must control that coupled fourth moment or the signed
actual-coefficient covariance, with all masks retained.

\subsection{A smooth profile and the exact covariance computation}
Let
\[
 \phi(t)=\begin{cases}\exp(1-1/(1-t^2)),&|t|<1,\\0,&|t|\ge1,\end{cases}
 \qquad W_x(y)=x\phi(4(400y-439))+\phi(4(400y-475)).
\]
For $x\ge0$ this is nonnegative and smooth; its disjoint supports lie
strictly inside $(1,2)$. For every integer $n$, $W_x(n/400)$ equals
$x$ at $n=439$, $1$ at $n=475$ and zero otherwise. Thus the computation
uses exact values of a genuine common $C_c^\infty$ profile, rather than
claiming a finite rational weight is smooth. At $x=10$, $\|W_x\|_\infty=10$.

The primary ideal generators, in the order used above, are
\[
 (-23,-18),\quad(-5,18),\quad(10,-15),\quad(25,15),
\]
where $(a,b)$ denotes $a+b\omega$. The first two have norm $439$.
The last two factor as $(-2,3)(-5,0)$ and $(-5,-3)(-5,0)$, with norms
$19\cdot25$. Since $19,25>17$, the long-factor count is $2$ there and
zero on the prime columns. All squarefree ideals at both selected norms
are included. The complete $B_7$ uses the four good primes of norms
$7,7,13,13$ for distinct ordered row factors and all six units. It has
$72$ distinct rows, maximum norm $2197$; each is nonprincipal by its
outside-$S$ Kummer valuation $1$ or $2$ modulo $6$, whose nontrivial
tame character has conductor exponent $1$.

For the raw Gram $G_{ij}=\sum_{u\in B_7}x_{m_i}(u)
\overline{x_{m_j}(u)}$, exact finite-field exponentiation and lattice
enumeration give the displayed real matrix. The weighted real vectors are
$\ell=(0,0,2,2)$, $v=(-x,-x,1,1)$, and $a=(x,x,1,1)$.
Therefore
\[
 Q(C_U,\mu)=D^{-1}\ell^t(\Re G)v=\frac{36-9x}{50}.
\]
Its diagonal contribution is $36/50>0$ and its off-diagonal contribution
is $-9x/50$. For $x=10$ this is $-27/25$ and gives
\eqref{sb:covariance-example}. The identity
$(\Re G-72I)^2=18^2I$ proves the real eigenvalues $54,90$;
these are not being asserted as the eigenvalues of the full complex Gram.
The exact source is \path{examples/arithmetic_cutoff/signed_covariance.py}.
A second test retains raw zeros on the $246$ rows $0<Nu\le89$ whose
outside-$S$ part is nonunit. Its covariance at $x=10$ is $-69/50$.
This subset is guaranteed nonprincipal; it is not called the entire
nonprincipal family. Here $89=\lfloor400^{3/4}\rfloor$.
These are finite counterexamples to automatic covariance positivity,
not asymptotic counterexamples to $BT_\sigma$ or GRH.

\subsection{Double Perron without a divergent critical-strip series}
For $T_0=S\cup\operatorname{supp}(u)$ retain
$F_0(s)=L^{T_0}(s,\psi)/L^{T_0}(2s,\psi_2)$, with the independently
primitive square character. Set $t_p=q_p/(1+q_p)$ outside $T_0$.
The finite actual cutoff is
\[
 J_U(s,u)=F_0(s)B_U(s),\qquad
 B_U(s)=\sum_{\substack{a,b\ {\rm sf},\ (a,b)=1\\
                  (ab,T_0)=1,\ Na,Nb\le U}}
 \mu(a)\mu(b)\prod_{p\mid ab}t_p.
\]
For $\sigma=\Re s>1/2$, the series
\[
 \mathcal G(s,z,w)=\sum_{\substack{a,b\ {\rm sf},\ (a,b)=1\\(ab,T_0)=1}}
 \mu(a)\mu(b)\prod_{p\mid ab}t_p\,(Na)^{-z}(Nb)^{-w}
 =\prod_{p\notin T_0}(1-t_p(Np)^{-z}-t_p(Np)^{-w})
\]
converges absolutely on $\sigma+\Re z>1$, $\sigma+\Re w>1$, uniformly
in all vertical parts at fixed positive margins. This follows from
$|t_p|\le(Np)^{-\sigma}/(1-(Np)^{-\sigma})$. There
$\mathcal F=F_0\mathcal G$, without evaluating an infinite $c$ series
in the critical strip. Initial absolute convergence also proves
\eqref{sb:perron-factor} by the local identity
\[
 (1-q)(1+q-q\alpha-q\beta)
 =(1-q\alpha)(1-q\beta)-q^2(1-\alpha)(1-\beta).
\]
The correction product is holomorphic on
\[
 \sigma+\Re z>0,\quad\sigma+\Re w>0,\quad
 2\sigma+\min(\Re z,0)+\min(\Re w,0)>1.
\]
On margins $\sigma+\Re z,\sigma+\Re w\ge\eta>0$ and last exponent
$\ge1+\epsilon$, its local deviation from $1$ is at most
$4(1-2^{-\eta})^{-2}(Np)^{-1-\epsilon}$. This proves normal convergence
and a phase/deletion-uniform vertical upper bound. Nonvanishing of the
correction is not asserted.

For $1/2<\sigma<1$, positive lines $c,d>1-\sigma$, and
$V=\lfloor U\rfloor+1/2$, define
\[
 \mathcal I_{T_1,T_2}=\frac1{(2\pi i)^2}
 \int_{c-iT_1}^{c+iT_1}\int_{d-iT_2}^{d+iT_2}
 \mathcal F(s,z,w)\frac{V^{z+w}}{zw}\,dw\,dz.
\]
The half-integer $V$ includes the actual equality cutoff, since
$Na\le U$ is exactly $Na<V$. Put
$\Delta_V=\inf_{n\in\mathbb Z_{\ge1}}|\log(V/n)|\ge1/(2V+1)$, and let
$\mathcal A=\sum_{a,b}|\mu(a)\mu(b)\prod_{p\mid ab}t_p|
(Na)^{-c}(Nb)^{-d}$ with the same disjoint support.
For $c,T>0$ and $v>0$, $v\ne1$, define
\[
 K_{c,T}(v)=\frac1{2\pi i}\int_{c-iT}^{c+iT}\frac{v^z}{z}\,dz.
\]
Finite-height termwise integration of the absolutely convergent
$\mathcal G$ and the one-variable tail estimate
$|K_{c,T}(v)-1_{v>1}|\le2v^c/(\pi T|\log v|)$ give
\begin{equation}
 |\mathcal I_{T_1,T_2}-J_U|
 \le |F_0|\mathcal A V^{c+d}
 \left\{\frac2{\pi T_1\Delta_V}+\frac2{\pi T_2\Delta_V}
       +\frac4{\pi^2T_1T_2\Delta_V^2}\right\}.      \label{sb:perron-error}
\end{equation}
Thus $\mathcal I_{T_1,T_2}\to J_U$ as both heights tend independently to
infinity, away from the original principal pole. This is an oscillatory
limit, not an absolutely integrable infinite vertical integral.

Taking $T_1=T_2=V^2$ gives the safe upper bound
$|J_U|\ll (Nu)^\eta |L(s,\psi)|V^{c+d}\log^2(2V)$, with constants
retaining fixed margins, $\nu,S$, and $\sigma$. Choosing $c,d$ close to
$1-\sigma$ reproduces $U^{2-2\sigma+\eta}$ in amplitude and
$U^{2-4\delta+\eta}$ in energy on $\sigma=1/2+\delta$.
Insertion into the original common-profile Mellin argument retains
$M_4(W)^2$ and all previously required height/conductor uniformities.
No derivative-free improvement follows from this Perron estimate.

The correction converges even on some negative auxiliary lines, but
moving there can cross $z=\rho-s$ or $w=\rho-s$, from zeros of the
primitive $L$-function. The exact meromorphic slice
$\mathcal F(s,z,0)=1/L^{T_0}(s+z,\psi)$ forbids assuming universal
cancellation of these poles. At a numerator zero the slice must be
understood meromorphically, not by an unjustified pointwise substitution.
Sequential shifts have single boundary integrals and crossed-zero
residues as well as the kernel double residue; away from $s$ zeros or
poles that double residue is $1/L^{T_0}(s,\psi)$.
A principal $L$ pole supplies a reciprocal zero, not a reciprocal pole.
Uniform control of actual reciprocal integrals/residues or the signed
covariance remains a new input. Numerator AFM alone supplies neither.

\section{A quantitative covariance budget without a new exponent}\label{cbq:core}
Retain the original Eisenstein sextic symbols, fixed finite-order $\nu$,
fixed $S$, one possibly complex common $W\in C_c^\infty$ supported in
$[1,2]$, all squarefree columns prime to $S$, and every raw shared-prime
zero. In the accepted window $2D/7\le U^2\le D$, put $r_m=r_U(m)$,
$\lambda_m=\mu(m)\nu(m)W(Nm/D)$ and $v_m=\overline{\lambda_m}$.
Thus $C_U=\mu r$ and $A_U=\mu(r-1)$ on the active annulus.
For a finite literal row block $B$, let
\[
 G_{mn}=\sum_{u\in B}x_m(u)\overline{x_n(u)},\qquad
 R=\operatorname{diag}(r_m),\qquad
 M_j=\sum_mr_m^j|W(Nm/D)|^2\quad(0\le j\le2).
\]
The primitive $\psi_u,\psi_{u,2}$ and full deleted sets retain their
previous meanings; no conductor-drop value replaces a raw zero.

\begin{proposition}[An actual-vector lower budget]\label{cbq:budget}
If $0\le G\le\Lambda_B I$, then
\begin{equation}\label{cbq:vector}
 \Re Q_B(C_U,\mu)\ge
 -\frac{\Lambda_B}{2D}\{\sqrt{M_0M_2}-M_1\}.
\end{equation}
In this high window, uniformly in $U$,
\begin{equation}\label{cbq:moments}
 \sum_{D\le Nm\le2D}r_m\ll_K D,\qquad
 \sum_{D\le Nm\le2D}r_m^2\ll_K D\log(2D).
\end{equation}
Consequently
\[
 \Re Q_B(C_U,\mu)\ge
 -C_K\|W\|_\infty^2\Lambda_B\sqrt{\log(2D)}.
\]
\end{proposition}
The vector inequality is sharp among positive contractions with the
specified $v,Rv$; no such sharpness is asserted for the arithmetic Gram.
The moment bound counts actual balanced ideal factorizations, rather than
assuming independence of the two row factors.

On $B_X$, inserting the previously accepted powered-BGL norm gives only
\[
 \Re Q_{B_X}(C_U,\mu)\ge
 -C_{\nu,S,\epsilon}\|W\|_\infty^2(DX)^\epsilon
 \{X^2+DX+D^{2/3}X^{5/3}\}.
\]
At $X\asymp D^{1/4}$ the exponents remain $1/2,5/4,13/12$.
At $H=D^{3/4}$ this budget still has the old $D^{1/2}$ excess over $H$;
it supplies no $BT_\sigma$ saving. A negative-part budget is weaker than
$2\Re Q(C_U,\mu)\ge Q(C_U,C_U)+Q(\mu,\mu)-R_\sigma$.

For $S_R=\sqrt R$, the exact identity
\[
 \tfrac12(RG+GR)=S_RGS_R+\tfrac12[S_R,[S_R,G]]
\]
localizes the commutator to unequal factor-count levels, with entry
$(\sqrt{r_m}-\sqrt{r_n})^2G_{mn}$. If $B$ retains all six literal units,
the actual unit sum further gives $G_{mn}=0$ unless $Nm\equiv Nn\pmod{36}$.
These are exact deletions, not a power saving. Even zero different-level
pairings leave the nonnegative within-level energies
$\sum_k(k-1)^2Q_B(\mu\mathbf1_{r=k},\mu\mathbf1_{r=k})$.
The unequal-product coupled fourth moment and its mixed M\"obius term
remain unestimated. No new analytic exponent, energy lower bound,
Lean formalization or GRH conclusion follows.

\section{Proof and arithmetic boundary of the covariance budget}\label{cbq:details}
\subsection{The positive-contraction calculation}
Since $G$ is the exact raw Gram, it is positive semidefinite. If
$\Lambda_B>0$, put $T=2G/\Lambda_B-I$, a self-adjoint contraction. Then
\[
 \Re(v^*RGv)=\frac{\Lambda_B}{2}
 \{M_1+\Re(v^*RTv)\}
 \ge\frac{\Lambda_B}{2}\{M_1-\|Rv\|_2\|v\|_2\},
\]
which is~\eqref{cbq:vector}; $\Lambda_B=0$ is immediate.
The Hermitian rank-two matrix
$\{(Rv)v^*+v(Rv)^*\}/2$ has possible nonzero eigenvalues
$(M_1\pm\sqrt{M_0M_2})/2$. Its negative spectral projection realizes
the lower bound among all $0\le G\le\Lambda_B I$, a linear-algebra
sharpness statement only. If $a\le r_m\le b$ on the weighted support,
$a+b>0$, the pointwise inequality $r_m^2\le(a+b)r_m-ab$ also gives
\[
 \Re Q_B(C_U,\mu)\ge
 -\frac{\Lambda_B M_0(b-a)^2}{8D(a+b)}.
\]
For $a=b=0$ the covariance is zero. The constant follows by maximizing
$\sqrt{(a+b)t-ab}-t$ on $[a,b]$.

\subsection{The balanced-divisor second moment}
Set $T_0=\sqrt D$. The high-window factor restrictions imply
\[
 \sqrt{2/7}\,T_0<Na,Nb<\sqrt{14}\,T_0.
\]
The elementary ideal count $\#\{a:Na\le Y\}\ll_K Y$ proves the
first assertion of~\eqref{cbq:moments}. For the second, two ordered
factorizations $ab=cd=m$ have the unique form
\[
 a=gh,\quad c=gk,\quad b=k\ell,\quad d=h\ell,
 \qquad g=(a,c).
\]
Squarefreeness makes $g,h,k,\ell$ pairwise coprime and squarefree.
To upper-bound this positive count, enlarge to arbitrary ideals.
For $M=\max(Nh,Nk)$ the individual bounds are
\[
 Ng\le\frac{\sqrt{14}\,T_0}{M},\qquad
 N\ell\le\frac{\sqrt{14}\,T_0}{M},\qquad
 1\le M\le\sqrt{14}\,T_0.
\]
On $2^j\le M<2^{j+1}$ there are $O_K(2^{2j})$ choices of $(h,k)$
and $O_K((T_0/2^j)^2)$ choices of $(g,\ell)$, with bounded terminal
ranges absorbed in the constant. Each dyadic range contributes $O_K(D)$;
there are $O(\log(2D))$ ranges. This proves the second moment with no
extra logarithm from counting only $NgN\ell$. Restricting to primes
outside fixed $S$ decreases the count. Hence
\[
 M_0,M_1\ll_K\|W\|_\infty^2D,\qquad
 M_2\ll_K\|W\|_\infty^2D\log(2D).
\]
These constants are independent of the cutoff in the stated window.
Only the common supremum norm is used, including for complex $W$.

\subsection{What the commutator actually controls}
Write
\[
 t_{mn}=D^{-1}\Re\{\lambda_m\overline{\lambda_n}G_{mn}\},\qquad
 \delta_{mn}=(\sqrt{r_m}-\sqrt{r_n})^2.
\]
The exact commutator identity proves
\[
 \Re Q_B(C_U,\mu)\ge
 Q_B(\sqrt r\,\mu,\sqrt r\,\mu)
 -\tfrac12\sum_{m,n}\delta_{mn}|t_{mn}|.
\]
The displayed loss retains the actual character and profile phases before
taking absolute values; it is at most the phase-blind loss obtained from
$|t_{mn}|\le D^{-1}|\lambda_m\lambda_nG_{mn}|$.
On the accepted four-column smooth $B_7$ fixture at amplitude $x\ge0$,
the positive term is $18/25$, this coefficient-adapted loss is $9x/50$,
and the lower bound is attained. Purely imaginary entries of the full
complex Gram make the phase-blind loss larger. This calibrates the new
budget and does not give an asymptotic obstruction.

For the unit deletion let $\zeta_6=1+\omega$. Every good prime norm $q$
is $1\pmod6$, and direct residue-field exponentiation gives
$x_p(\zeta_6)=\zeta_6^{(q-1)/6}$. Multiplying over $p\mid m$ and reducing
$Nm=\prod_{p\mid m}(1+6e_p)$ modulo $36$ gives
\[
 x_m(\zeta_6)=\zeta_6^{(Nm-1)/6\bmod6}.
\]
Multiplication of all literal rows by $\zeta_6$ preserves their raw zeros
and proves the norm-congruence deletion in the core. The vector budget
may be applied separately on these six classes. No D-dependent number
of classes or conductor saving is produced.

The moment proof above counts the coefficient diagonal $ab=cd$.
It has not estimated the unequal-product character quadruples in the
coupled row fourth moment. Replacing the inherited $DX$ norm by
$D^{1-\sigma}X$ merely because equal-level and different-class entries
vanish is the unsupported step. Small negative commutator loss alone
also leaves the within-level energies in the core. Principal residues
are unchanged; this covariance is used on the nonprincipal blocks where
the original residue is zero, not as a new centering prescription.

\subsection{The complete near-saturating face}
For a six-factor row block let $R_0=\prod_{j=1}^6A_j$,
$M=\max(A_1,A_2)$ and $\prod_j A_j^j\le H$.
The accepted exact-mask block norm is
\[
 \Lambda_A\ll_{\nu,S,\epsilon}(HD)^\epsilon
 \{R_0+DR_0/M+D^{2/3}R_0/M^{1/3}\}.
\]
The column moment and covariance proofs apply to every such block.
Write $x_j=\log A_j$, $h=\log H$, $c_0=\sum_jjx_j$ and
$r_0=\sum_jx_j$. Then, for $x_1\ge x_2$,
\[
 h-3(r_0-\max(x_1,x_2))
 =h-c_0+(x_1-x_2)+x_4+2x_5+3x_6;
\]
for $x_2\ge x_1$ replace $x_1-x_2$ by $2(x_2-x_1)$.
All terms are nonnegative. Thus the middle monomial is saturated on
\[
 A_1=A_2,\qquad A_1A_3=H^{1/3},\qquad A_4=A_5=A_6=1,
\]
with full capacity. The balanced $A_3=1$ case is only an endpoint.
The $a_3$ direction remains nonprincipal by its outside-$S$ valuation
$3\pmod6$. If $R_0/M\ge H^{1/3}D^{-\eta}$, the total displayed slack
is at most $3\eta\log D$. Therefore a fixed power gain must be uniform
on this full neighboring region. The proved budget still costs $DR_0/M$;
it gives no global improvement over the accepted envelope.

The independent standard-library diagnostic checks nine ideal moment
cases, eight raw Gram cases with $27\,000$ phase values and $5\,184$ raw
zeros, exact radical commutator/vector budgets, and $227$ exponent cases.
Finite tables use common norm profiles realizable by smooth disjoint bumps;
they are not tests of an infinite analytic estimate. Full derivation,
source bindings and preserved checks are under
\path{planning/covariance-budget-20261105/attempt/} and
\path{outputs/checks/2026-11-05-covariance-budget/attempt/}.

\section{The fourth-moment conductor and its full raw zero set}
\label{upq:core}
Retain the original Eisenstein sextic symbols $x_m(u)=(u/m)_6$,
fixed finite-order $\nu$, fixed $S$, and one possibly complex common
smooth $W$ supported in $[1,2]$. Put
$\lambda_m=\mu(m)\nu(m)W(Nm/D)$ and $r_m=r_U(m)$.
Columns remain squarefree and prime to $S$. In $2D/7\le U^2\le D$, the actual polynomial is
\[
 \mathcal C_U(u)=
 \sum_{\substack{(a,b)=1,\ Na,Nb>U\\D\le Na\,Nb\le2D}}
 \mu(a)\mu(b)\nu(ab)W(Nab/D)x_a(u)x_b(u).
\]
Its squared row energy has four factor indices $a,b,c,d$, with
$(a,b)=(c,d)=1$. These internal masks are essential. At each prime,
\[
 e_p=\mathbf1_{p\mid a}+\mathbf1_{p\mid b}
       -\mathbf1_{p\mid c}-\mathbf1_{p\mid d}
 \quad\hbox{belongs to }\{0,\pm1\}.
\]
Thus, for the two squarefree products $m=ab,n=cd$, the primitive
oscillatory conductor outside fixed $S$ is $q=mn/(m,n)^2$.
The full raw zero-preserving period is $L=mn/(m,n)$:
\[
 x_m(u)\overline{x_n(u)}
 =\mathbf1_{(u,(m,n))=1}\theta_{m/(m,n),\,n/(m,n)}(u).
\]
No common-prime zero is restored by taking the primitive character.

A generic fourth moment $|\sum_m\lambda_mx_m(u)|^4$ has four individually
squarefree columns $m_1,\ldots,m_4$ with no such internal coprimality.
Its signed local exponent
$e_p=\sum_{i=1}^2\mathbf1_{p\mid m_i}-\sum_{i=3}^4\mathbf1_{p\mid m_i}$
can also be $\pm2$, giving a nontrivial cubic component. For $L=\operatorname{rad}(m_1m_2m_3m_4)$, the exact row-power
interface is
\[
 Q_j=\prod_{\substack{p\mid L\\je_p\not\equiv0\ (6)}}p,\qquad
 Z_j=L/Q_j,\qquad
 \prod_{i=1}^2x_{m_i}(v^j)
       \overline{\prod_{i=3}^4x_{m_i}(v^j)}
 =\theta_{j,e}(v)\mathbf1_{(v,Z_j)=1}.
\]
Here $\theta_{j,e}$ is primitive on $Q_j$, with its own zero extension.
At $j=3$ an exponent $\pm2$ becomes primitive-principal but its raw
zero remains. For the original two-squarefree-product expansion,
$Q_j=q$ for all $1\le j\le5$, including the quadratic $j=3$ direction.
For the accepted row blocks $Nr_j\asymp A_j$, put
$R_0=\prod_jA_j$, $P=\max(A_1,\ldots,A_5)$ and retain
$\prod_jA_j^j\le H$. The full $A_3$ saturation face and all row-factor
masks remain.

The small-conductor pair count recovers an already controlled high-gcd
sector; it is not a new analytic advance. On the remaining original
sector $Nq>P^2$, favorable smooth completion
costs $\min(A_j,\sqrt{Nq})=A_j$ for $1\le j\le5$.
The literal sharp cutoff is not replaced by this smooth trial.
The signed large-conductor fourth-moment and mixed M\"obius sums remain
open; no improved full exponent, $BT_\sigma$, new Lean theorem,
zero-location result or GRH conclusion follows.

\section{Scope of the raw fourth-moment reduction}
\label{upq:details}
Write $c_{a,b}=\mu(a)\mu(b)\nu(ab)W(Nab/D)$ on the exact factor support
in the core. Before absolute Cauchy, the normalized energy is
\[
 \frac1D\sum_{\substack{(a,b),(c,d)\ {\rm on\ that\ support}}}
 c_{a,b}\overline{c_{c,d}}
 \sum_{u\in B}x_a(u)x_b(u)\overline{x_c(u)x_d(u)}.
\]
Squarefreeness of $ab$ and $cd$ forces each positive and negative
prime multiplicity to be at most one, proving $e_p\in\{0,\pm1\}$.
For the mixed term against the actual $\mu$ polynomial the exponent is
$\mathbf1_{p\mid a}+\mathbf1_{p\mid b}-\mathbf1_{p\mid n}$,
with the same conclusion and the same full union mask.
The common possibly complex $W$, $\nu$, hyperbolic support and all raw
zeros stay in these identities.

In contrast, unrestricted four-column memberships $(1,1,0,0)$ give
$e_p=2$, a cubic component; $(1,1,1,0)$ give $e_p=1$ although both
product supports contain $p$. Neither the parity symmetric difference
of four supports nor the symmetric difference of two nonsquarefree
product supports gives the general conductor.
For the interface in the core, its local character order is
\[
 \begin{array}{c|cccc}
 j&1,5&2,4&3&6\\ \hline
 e_p=\pm1&6&3&2&1\\
 e_p=\pm2&3&3&1&1
 \end{array}
\]
Nontrivial entries have tame conductor exponent one; trivial entries
remain deleted at $p\mid v$. Multiplying the interface over all six
row factors, with the fixed unit/$S$ phase, preserves every original
zero, including those at sixth-power factors. All inter-factor,
squarefree, ray and capacity restrictions remain external restrictions
on the same row sum; this identity asserts no independence.

At a residue field of size $Q$, let $\chi$ have order six and $\rho$
be nontrivial additive. The zero-extended complete sum
$G_k(t)=\sum_{z\ne0}\chi(z)^k\rho(tz)$ satisfies
\[
 \begin{array}{c|cc}
 &t=0&t\ne0\\ \hline
 k\not\equiv0\ (6)&0&|G_k(t)|=\sqrt Q\\
 k\equiv0\ (6)&Q-1&-1.
 \end{array}
\]
For $t\ne0$, substitution gives
$G_k(t)=\overline{\chi(t)^k}G_k(1)$; for nontrivial $k$,
expanding the norm square gives
$(Q-1)-\sum_{a\ne0,1}\chi(a)^k=Q$.
The proof applies at split and inert primes. By CRT the full period-$L$
transform vanishes if any local frequency on $Q_j$ is zero; otherwise
its magnitude is
\[
 (NQ_j)^{1/2}
 \prod_{p\mid Z_j}
 \left|-1+Np\,\mathbf1_{\text{local frequency at }p\text{ is }0}\right|.
\]
Canceled common primes thus supply Ramanujan factors.

Exact deletion, with all remaining restrictions kept in $F$, gives
\[
 \sum_r\theta(r)F(r)\mathbf1_{(r,Z)=1}
 =\sum_{d\mid Z}\mu(d)\theta(d)\sum_z\theta(z)F(dz).
\]
For a squarefree pivot let $E$ be the radical of its fixed excluded
primes, including $Z$ and the relevant other row factors, with
conductor primes omitted since $\theta$ already vanishes there.
The additional squarefree identity is
\[
 \mu(r)^2\mathbf1_{(r,E)=1}
 =\sum_{\substack{h^2\mid r\\(h,E)=1}}\mu(h)
       \sum_{\substack{d\mid E\\d\mid r}}\mu(d).
\]
Since $(d,h)=1$, its lifted row sum has terms
$\mu(d)\mu(h)\theta(d)\theta(h)^2\sum_z\theta(z)F(dh^2z)$.
The scaled $F$ retains the exact window, primary/ray restrictions and
other masks. Divisor triangles cost at most $(HD)^\epsilon$;
they do not restore a forbidden nonzero.

The previous fixed-radial-profile high-gcd proof
(\path{research/short-row-high-gcd.md}, independently reviewed in
\path{planning/reviews/book/short-row-formal.md})
already bounds its off-diagonal tail by
$O_{\Phi,\epsilon}(D^{2+\epsilon}/G_0^2)$ for $Nh\ge G_0$.
For annular $m=hr,n=ht$, $Nq\le Y$ implies $Nh\ge D/\sqrt Y$,
so the old smooth-scope bound is $O_{\Phi,\epsilon}(D^\epsilon Y)$.
The present elementary count $O(D\sqrt Y)$, factor multiplicities
$r_U(m)r_U(n)\ll_\epsilon D^\epsilon$, and $\#B_A=O(R_0)$ give only
$O(\|W\|_\infty^2D^\epsilon R_0\sqrt Y)$ on a literal sharp block.
This is a weaker conductor-form recovery in its own bookkeeping scope,
not a new high-gcd theorem; the smooth estimate is not silently
transferred to a selected sharp off-diagonal piece.
The full $A_3$ face and neighborhood slack are also inherited.

Finally, primitive smooth Poisson completion has bound
$O_V(\min(A_j,\sqrt{Nq}))$ on the original support.
For $Nq>P^2$ and $1\le j\le5$, every scaled length
$A_j/(Nd\,Nh^2)$ is below $\sqrt{Nq}$, so even this favorable trial
takes its trivial length term. Its constant uses fixed smooth
seminorms and need not stay bounded when a sharp boundary is approached.
In the admissible subsector $(m,n)=1$, $Nq=Nm\,Nn\asymp D^2$.
A complete zero mean does not control the nonzero frequencies of the
short original masked row window. The unresolved target is their
joint signed sum with weights $r_U(m)r_U(n)\lambda_m\overline{\lambda_n}$
and $r_U(m)\lambda_m\overline{\lambda_n}$, before per-conductor absolute
estimates. No new global analytic conclusion is claimed.

% Independently reviewed proposal; classical GL input, no new sieve/global exponent.
\section{A quadratic pivot on the cube factor}\label{qp:core}
Retain the original zero-extended sextic symbols
$\chi_m(u)=(u/m)_6$ over $K=\mathbb Q(\omega)$, the fixed $S$ containing
all primes above $6$ and comparison moduli, and the squarefree good columns
$D\le Nm\le2D$. Write an outside-$S$ row ideal uniquely as
$r_1r_2^2r_3^3r_4^4r_5^5r_6^6$: only $r_1,\ldots,r_5$ are pairwise-coprime
squarefree ideals; $r_6$ is arbitrary and may overlap them. Use primary
generators. Fix one literal unit and one $S$-supported row factor, retaining
their original values. Let $A_j\le Nr_j<2A_j$, $A_j\ge1$,
$R_0=\prod_jA_j$ and $\prod_jA_j^j\le H$; the actual capacity restriction
on rows is retained. Coefficients $b_m\in\mathbb C$ are arbitrary and fixed
independently of every variable row.

\begin{proposition}[The cube-factor block norm]\label{qp:bound}
For every $\varepsilon>0$, the raw Gram on this actual block obeys
\[
 \sum_{u\ \mathrm{in\ the\ block}}
       \left|\sum_m b_m\chi_m(u)\right|^2
 \ll_{\nu,S,\varepsilon}(HD)^\varepsilon
       \left\{R_0+\frac{DR_0}{A_3}\right\}\sum_m|b_m|^2.
\]
The constant is uniform in the fixed unit/$S$-factor, frozen good factors,
$H,D$ and the coefficient sequence. There is no rational-norm-squarefree
restriction. Consequently, with $M=\max(A_1,A_2)$, the old powered-BGL
bound and this bound give
\begin{equation}\label{qp:min}
 \Lambda_A\ll_{\nu,S,\varepsilon}(HD)^\varepsilon
 \left\{R_0+\min\left(
  \frac{DR_0}{M}+\frac{D^{2/3}R_0}{M^{1/3}},
  \frac{DR_0}{A_3}\right)\right\}.
\end{equation}
\end{proposition}
The external analytic input is Goldmakher--Louvel,
\href{https://arxiv.org/pdf/1112.1642v2}{arXiv:1112.1642v2},
Corollary~1.2 with $n=6$, already an explicit classical cube-power sieve.
The refinement is this application, not a new large-sieve theorem.
For the original normalized coefficients
$b_m=D^{-1/2}A_U(m)\nu(m)W(Nm/D)$, the accepted divisor bound gives
$\sum_m|b_m|^2\ll_\eta\|W\|_\infty^2D^\eta$, uniformly in
$U\ge1$ with $U^2\le D$ (and the common $W$ may be complex).
Choosing smaller internal exponents yields the corresponding uniform
block energy envelope. General cutoff or factor-count weights remain
inside the arbitrary coefficient norm; that norm is not silently replaced
by one. Summing $S$-supported row factors uses the reduced capacity
$H/N\mathfrak b$ and the established convergent geometric sums.

On the complete old saturated face $A_1=A_2=A$, $A_3=B$,
$A_4=A_5=A_6=1$, $AB=H^{1/3}$, the new bound is
$A^2B+DA^2$ in place of the old middle monomial $DAB$.
For $H=D^{3/4}$ put $A=D^a$, $B=D^{1/4-a}$,
$0\le a\le1/4$. When $a<1/8$, its leading exponent is $1+2a<5/4$;
the gain requires $B/A$ to grow by a power for a power improvement.
For $a\ge1/8$, the combined bound still has leading exponent $5/4$.
Thus the $A_3$-dominant part improves, while no new global exponent,
$BT_\sigma$ estimate, signed covariance theorem, Lean theorem or GRH
conclusion follows.

\paragraph{Inspecting the surviving sector.}
The \href{https://amelie-iska.github.io/GRH-research/quadratic-pivot.html}{quadratic-pivot inspector} plots the complete face, the five monomial powers and the remaining target gap. These are analytic bound exponents, not measured operator norms or trained-network predictions. The refined norm can be inserted in the earlier covariance budget; it does not remove its within-level energies.

\paragraph{A necessary near-face condition.}\label{qp:near}
Put $\mathcal M_A=DR_0/\max(M,A_3)$.
The combined envelope is at most
$(HD)^\varepsilon\{R_0+\mathcal M_A+D^{2/3}R_0/M^{1/3}\}$.
If $\mathcal M_A\ge DH^{1/3}D^{-\eta}$, then the earlier
near-face slack is at most $3\eta\log D$ and, additionally,
$A_3\le M D^\eta$. Thus a uniform remaining leading-term estimate
may focus on this smaller neighborhood; it must still cover the entire
$B\le A$ part of the face. This is a necessary condition on explicit
upper-bound monomials, not a lower bound for actual arithmetic energy.

% Independently reviewed proof proposal; all raw collision zeros retained.
\section{Proof and scope of the quadratic pivot}\label{qp:details}
Freeze all factors except $r_3$, and write the actual literal row
$u=Fr_3^3$. Multiplicativity of the original numerator gives the exact
zero-preserving identity
\[
 \chi_m(u)=\chi_m(F)\chi_m(r_3)^3.
\]
Put $c_m=b_m\chi_m(F)$, a coefficient sequence independent of $r_3$ with
$\sum|c_m|^2\le\sum|b_m|^2$. In particular the arbitrary sixth-power
factor is the coefficient mask
\[
 \chi_m(r_6^6)=\mathbf1_{(m,r_6)=1},
\]
not the constant one. No condition $(r_3,r_6)=1$ is introduced.
The actual capacity and coprimality of $r_3$ with
$r_1,r_2,r_4,r_5$ only restrict this positive outer row sum; after freezing,
it may be enlarged to all squarefree good pivot ideals in the dyad.
The variable collision mask $(m,r_3)$ remains in the original symbol.

Apply the previously proved literal comparison
\eqref{wc:literal-reciprocity} to the primary generator of $r_3$ itself:
\[
 \chi_m(r_3)^3=
 \eta_{r_3}(m)^3 C([r_3],[m])^3\vartheta_m(r_3)^3.
\]
The twist $\eta_{r_3}$ ranges over a fixed finite set of $S$-unit Kummer
characters, and $[r_3]$ over the fixed finite reciprocity group.
Partition pivot rows by $(\eta_{r_3},[r_3])$. On each part the first two
factors are one fixed modulus-one column multiplier. Thus its coefficient
sequence is independent of the pivot row and has squared norm at most
$\sum|b_m|^2$. This finite partition, rather than an assertion that arbitrary
row-dependent phases may be absorbed, resolves the fixed-$S$ phase issue.

Goldmakher--Louvel Definition~1 and Remark~3 identify the cubic powers
of a sixth-order Hecke family as a quadratic Hecke family; their
Corollary~1.2 gives directly
\[
 \sum_{a\ \mathrm{sf,good}\atop Na\le2A_3}
 \left|\sum_{m\ \mathrm{sf,good}\atop Nm\le2D}
                z_m\vartheta_m(a)^3\right|^2
 \ll_{K,S,\eta}(A_3D)^\eta(A_3+D)\sum_m|z_m|^2.
\]
This is the required column-indexed orientation and allows arbitrary
complex coefficients and squarefree ideals of nonsquarefree rational norm.
The canonical family has primitive trivial-infinity-type characters, a
fixed finite reciprocity group, and the same-class primitive-product
property required by Definition~1. At each $p\mid m$ outside $S$, the
local tame sextic character has order six, so its cube has order two and
conductor exponent one. Thus its outside-$S$ primitive conductor is
exactly $m$; the $S$-part has only finitely many fixed conductor types.
Accordingly both sides of the comparison are zero at $(m,a)>1$.
No conductor-drop value replaces the raw zero, and no global formula
for the full conductor is required. The fixed $\nu$ stays in coefficients.

Enlarge the positive row sum separately on each finite part and apply
this sieve. Ideal counting gives $O_K(A_j)$ possibilities for each frozen
$r_j$, including arbitrary $r_6$, hence $O_K(R_0/A_3)$ frozen tuples.
Summing gives $(A_3D)^\eta(R_0+DR_0/A_3)\sum|b_m|^2$.
Since $A_3\le H$, adjust $\eta$ to obtain
Proposition~\ref{qp:bound}; literal unit multiplicity is at most six.
Together with the prior block bound this proves
\eqref{qp:min}. It does not change the actual coefficient
normalization or principal residues.

On the face in the core, the old bound has exponents
$a+1/4$, $5/4$, $11/12+2a/3$, whereas the quadratic bound has
$a+1/4$ and $1+2a$. Its improvement for $a<1/8$ leaves the balanced
point and $a\ge1/8$ with the old $5/4$ leading exponent. Even the
improved endpoint $a=0$ costs $D$, above the target $D^{3/4}$.
The pinned GL PDF has SHA256
\texttt{5cc5f03388d3996b4333f07be51f989f}\allowbreak\texttt{60779fe172f8bd24c7aa6ddad64dd304}.
This proof invokes that classical result; finite exponent checks do not
reprove it, and no new formalization is claimed.

For the last condition, use $R_0/M\le H^{1/3}$ from the exact
capacity inequalities. The hypothesis implies
$R_0/A_3\ge H^{1/3}D^{-\eta}$, so
$A_3/M=(R_0/M)/(R_0/A_3)\le D^\eta$.
It also implies $R_0/M\ge H^{1/3}D^{-\eta}$, which is exactly
the previous slack hypothesis. Finally
$\min(a+b,c)\le\min(a,c)+b$ for nonnegative $a,b,c$
proves the separated envelope, with no replacement of a signed moment
by a positive lower bound.

% Compact selected proposal; exact algebra, no new analytic exponent.
\section{Joint completion before conductor absolute values}\label{jdq:core}
In the original C--C expansion write $m=hr,n=ht$ with $h,r,t$
pairwise coprime squarefree, $q=rt\ne1$, $L=hq$, $Q=Nq$, $\ell=NL$.
Set $\theta=x_r\overline{x_t}$ and $\xi_q=x_q^3$.
The original internal masks give $\theta^3=\xi_q$ and primitive
$\theta^j$ on $q$ for $1\le j\le5$, while
$\mu(hr)\mu(ht)=\mu(q)$. Fix generators $\alpha$ of $h$, $\beta$ of
$q$, choosing $\alpha=1$ for $h=1$. With $\delta=1+2\omega$, put
$e_\beta(z)=\exp(2\pi i\operatorname{Tr}(z/(\delta\beta)))$,
$\tau_\beta(\eta)=\sum_{z\bmod q}\eta(z)e_\beta(z)$, and
$c_h(k)=\prod_{p\mid h}(Np\mathbf1_{k\in p}-1)$, $c_1=1$.
For $f_j(v)=\theta(v)^j\mathbf1_{(v,h)=1}$, $1\le j\le5$,
$f_6(v)=\mathbf1_{(v,L)=1}$ and
$G_j(k)=\sum_{v\bmod L}f_j(v)e_{\alpha\beta}(kv)$, joint completion gives
\begin{equation}\label{jdq:gauss}
 \prod_{j=1}^6G_j(k_j)=Q^2\xi_q(-\alpha)\tau_\beta(\xi_q)
 \overline{\theta(\prod_{j=1}^5k_j^j)}
 \prod_{j=1}^5c_h(k_j)c_L(k_6).
\end{equation}
This holds when the first five frequencies are $q$-units; otherwise
the left side is zero. The mixed Gauss phase is orientation-independent,
with every common-prime zero retained. The exact row Fourier weight
still contains all squarefree, ray, inter-factor, capacity and profile
restrictions, including the full $A_3$ face. Its weighted joint
conductor/frequency estimate remains open. This identity and its finite
Parseval budget give no improved global exponent or GRH conclusion.

% Compact selected derivation; longer research notes remain preserved.
\section{The surviving joint identity and its missing estimate}\label{jdq:details}
CRT gives $G_j(k)=\theta(\alpha)^j\tau_\beta(\theta^j)
\overline{\theta(k)^j}c_h(k)$ for $q$-unit $k$, and $G_6=c_L$.
Pairing Gauss powers $1,5$ and $2,4$ yields
$\prod_{j=1}^5\tau_\beta(\theta^j)=Q^2\xi_q(-1)\tau_\beta(\xi_q)$;
the CRT factors multiply to $\theta(\alpha)^{15}=\xi_q(\alpha)$.
This proves~\eqref{jdq:gauss}, with all additive-generator phases explicit.

On a fixed literal class write $u=u_0\prod_{j=1}^6v_j^j$ and let $F(v)$
be its exact finite row weight. Only $v_1,\ldots,v_5$ are squarefree and
pairwise coprime; $v_6$ is arbitrary with its permitted overlaps.
Retain every original mask and the sharp capacity boundary in $F$.
Define $\widehat F_L(k)=\sum_vF(v)e_{\alpha\beta}(-\sum_jk_jv_j)$ and,
with $t=q/r$, put
\begin{align*}
 a(r)&=r_U(hr)r_U(ht)\nu(r)\overline{\nu(t)}\\
 &\qquad{}\times W(Nh\,Nr/D)\overline{W(Nh\,Nt/D)}\theta_{r,t}(u_0),\\
 A(K)&=\sum_{r\mid q}a(r)\overline{\theta_{r,q/r}(K)},\\
 B(K)&=\sum_{\substack{k\bmod L,\ (k_1\cdots k_5,q)=1\\
                   \prod_{j=1}^5k_j^j=K\bmod q}}
       \widehat F_L(k)\prod_{j=1}^5c_h(k_j)c_L(k_6).
\end{align*}
Set $a(r)=0$ at excluded orientations; the actual factor counts and
common possibly complex $W$ remain. Fourier inversion gives exactly
\begin{equation}\label{jdq:joint}
 E_{\mathcal T}=\frac1D\sum_{(h,q)\in\mathcal T}
 \mu(q)\frac{Q^2\xi_q(-\alpha)\tau_\beta(\xi_q)}{\ell^6}
 \sum_{K\in(\mathcal O/q)^\times}B(K)A(K).
\end{equation}
Sum the original $u_0$ classes and blocks. This retains low gcd,
including $h=1$, before conductor absolute values. For C--$\mu$ replace
the second $r_U(ht)$ by $1$. A column-orientation-dependent restriction
cannot be hidden in a common $F$ when applying the next identity.

Distinct orientations have nontrivial cubic quotient, so
\[
 \sum_K|A(K)|^2=\varphi(q)\sum_{r\mid q}|a(r)|^2,\qquad
 |E_{h,q}|\le\frac{Q^{5/2}}{D\ell^6}
  \left(\varphi(q)\sum_{r\mid q}|a(r)|^2\right)^{1/2}\|B\|_2.
\]
The product-frequency fibers have size $\varphi(q)^4$ under uniform
averaging. The actual $B$ is not uniform: all orientations agree on
the coherent set $x_p(K)^2=1$ for every $p\mid q$, of size
$\varphi(q)/3^{\omega(q)}$. Centered representatives have
$|k_j|\le\sqrt\ell$. Restricting to any frequency subset $\Omega$ has
error bounded by the same displayed budget with $\|B-B_\Omega\|_2$;
no small lower-height tail is proved.

The full face $A_1=A_2=A$, $AA_3=H^{1/3}$, $A_4=A_5=A_6=1$ remains.
At its pure $A_3$ endpoint, inversion in the fixed factors leaves
$(\sum_ra(r))\sum_vF_3(v)\xi_q(v)\mathbf1_{(v,h)=1}$:
the actual row character there is orientation-independent. Literal
length-one Fourier factors can undo complete phase balance.
A quantitative signed estimate for the actual $h,q,K$ sum is still
required. Exact split/inert Gauss, full-mask CRT and orientation tests
pass $991$ assertions; receipt:
\path{outputs/checks/2026-11-05-joint-dual/research/joint-gauss-check-v1.json}.
These diagnostics do not prove an incomplete-moment saving or GRH.

% Fresh proposal; ordinary classical sieve application, no new GRH/Lean theorem.
\section{A radical-conductor refinement of the raw envelope}\label{mcs:core}
Retain the original literal symbols $\chi_m(u)=(u/m)_6$, extended by zero,
the fixed good set $S$, and the squarefree ideal columns $D\le Nm\le2D$.
On an actual dyadic row block write
\[
 u=u_0r_1r_2^2r_3^3r_4^4r_5^5r_6^6,\qquad
 A_j\le Nr_j<2A_j,\qquad \prod_{j=1}^6A_j^j\le H.
\]
Only $r_1,\ldots,r_5$ are pairwise-coprime squarefree ideals;
$r_6$ is arbitrary and may overlap them. Fix primary generators,
a literal unit and an $S$-supported factor in $u_0$. Put
$R=\prod_jA_j$, $M=\max(A_1,A_2)$ and $Q=R/A_6$.
For arbitrary complex coefficients $b_m$, independent of every variable row,
define
\[
 \mathcal E_A(b)=\sum_{u\ \mathrm{in\ the\ actual\ block}}
       \left|\sum_m b_m\chi_m(u)\right|^2,
 \qquad \Lambda_A=\sup_{b\ne0}\frac{\mathcal E_A(b)}{\sum_m|b_m|^2}.
\]
\begin{proposition}[An ordinary conductor comparison]\label{mcs:block}
For every $\varepsilon>0$,
\begin{equation}\label{mcs:block-bound}
 \Lambda_A\ll_{S,\varepsilon}(HD)^\varepsilon
             A_6\{D+Q^2\}.
\end{equation}
Consequently the original residue-subtracted remainder satisfies,
uniformly for $H,D,U\ge1$ and $U^2\le D$,
\begin{equation}\label{mcs:uniform}
 \sum_{0<Nu\le H}|E_u|^2
 \ll_{\nu,S,\varepsilon}\|W\|_\infty^2(HD)^\varepsilon
       \{H+DH^{1/6}+(DH)^{2/3}\}.
\end{equation}
Here $\nu,S$ are fixed and $W\in C_c^\infty((0,\infty))$ is one common,
possibly complex profile supported in $[1,2]$; there is no profile supremum.
\end{proposition}
The input to~\eqref{mcs:block-bound} is the classical ordinary sieve for
\emph{all} primitive characters of conductor norm at most $Q$, with constant
$D+Q^2$, not a higher-order-family theorem for arbitrary conductor labels.
The following section gives its standard Gauss-sum and two-dimensional
lattice proof, including prime-power conductors at $S$.
For fixed $r_6$, the first five factors have primitive outside-$S$
conductor $q=r_1\cdots r_5$. The literal primitive character agrees with
the raw symbol even at collisions with $q$; the sixth-power factor instead
remains the coefficient mask $\mathbf1_{(m,r_6)=1}$.
At most $5^{\omega(q)}\ll_\eta(Nq)^\eta$ tuples give the same primitive
character. Applying the ordinary sieve and then counting $r_6$ proves
\eqref{mcs:block-bound}. No rational-norm-squarefree restriction is imposed.

To obtain~\eqref{mcs:uniform}, combine this with the accepted powered-BGL
block bound from Section~\ref{wc:core}. With
$S_A=R/M$ and $T_A=D^{2/3}R/M^{1/3}$, capacity gives
$R\le H$, $A_6\le H^{1/6}$, $RS_A\le H$ and
$T_A\le(DH)^{2/3}$. Hence
\begin{align*}
 \min(R+DS_A+T_A,DA_6+R^2/A_6)
 &\le R+T_A+DA_6+\min(DS_A,H^2/S_A^2)\\
 &\le H+DH^{1/6}+2(DH)^{2/3}.
\end{align*}
Dyadic and fixed-$S$ summation, the uniform actual coefficient norm,
and the principal-residue estimate give the assertion for $E_u=I_u-P_u$.

At $H=D^{3/4}$ the new envelope has exponent $7/6$, improving the old
$5/4$ by $1/12$; its excess above the target $H$ is still $D^{5/12}$.
On the \emph{entire} old face $A_1=A_2=A$, $A_3=B$,
$A_4=A_5=A_6=1$, $AB=H^{1/3}$, one has
$Q=A^2B=AH^{1/3}\le H^{2/3}$, so~\eqref{mcs:block-bound} costs
$D^{1+\varepsilon}$ there. This includes $B\le A$.
The new combined upper bounds reach $7/6$ at
$(\log_D A_1,\log_D A_2)=(5/12,1/6)$ with other coordinates zero;
the pure point $\log_D A_1=3/4$ also retains the old BGL cross term.
These are upper-envelope witnesses, not arithmetic energy lower bounds
or counterexamples to signed cancellation.

The simpler proposed replacement of the radicand $ab^2$ by its radical
$ab$ inside BGL remains unjustified: after the existing fixed reciprocity
phases, the mixed reciprocity quotient is
$\vartheta_b(a')/\vartheta_a(b')$, which varies outside $S$.
The present application does not require that replacement.
No target signed moment, type-II covariance estimate, positive arithmetic
determinant comparison, zero-location result, new Lean theorem or GRH
conclusion follows from~\eqref{mcs:uniform}.

% Proof proposal; full raw zeros and arbitrary fixed coefficients retained.
\section{Proof of the ordinary-conductor refinement}\label{mcs:details}
\subsection{Primitive Gauss sums at every conductor}
Let $\mathcal O=\mathbb Z[\omega]$, $\delta=1+2\omega$, and choose one
generator $f$ of each nonzero ideal $\mathfrak f$. Write
$e_f(x)=\exp(2\pi i\operatorname{Tr}(x/(\delta f)))$,
$U(\mathfrak f)=(\mathcal O/\mathfrak f)^\times$, and
$\varphi(\mathfrak f)=|U(\mathfrak f)|$.
This additive character is primitive, since the different is $(\delta)$.
For a primitive multiplicative character $\rho$ modulo $\mathfrak f$,
extended by zero at nonunits, set
$\tau_f(\rho)=\sum_{x\bmod\mathfrak f}\rho(x)e_f(x)$.
We verify the full-modulus identity
\begin{equation}\label{mcs:gauss}
 |\tau_f(\rho)|^2=N\mathfrak f,\qquad
 \sum_{x\bmod\mathfrak f}\rho(x)e_f(zx)
       =\overline{\rho(z)}\tau_f(\rho)\quad(z\in\mathcal O).
\end{equation}
No squarefree-conductor hypothesis is used, so this includes wild
prime powers in the fixed $S$-part.

The Ramanujan sum is
\[
 c_{\mathfrak f}(y)=\sum_{x\in U(\mathfrak f)}e_f(xy)
   =\sum_{\mathfrak d\mid(\mathfrak f,(y))}
                   N\mathfrak d\,\mu(\mathfrak f/\mathfrak d).
\]
Indeed insert $\mathbf1_{(x,\mathfrak f)=1}
=\sum_{\mathfrak b\mid(x,\mathfrak f)}\mu(\mathfrak b)$;
the complete additive sum with $\mathfrak b\mid x$ is
$N(\mathfrak f/\mathfrak b)$ when
$\mathfrak f/\mathfrak b\mid(y)$ and zero otherwise.
Changing variables $x=ay$ in the Gauss square now gives
\begin{align*}
 |\tau_f(\rho)|^2
 &=\sum_{a\in U(\mathfrak f)}\rho(a)c_{\mathfrak f}(a-1)\\
 &=\sum_{\mathfrak d\mid\mathfrak f}N\mathfrak d\,
       \mu(\mathfrak f/\mathfrak d)
       \sum_{\substack{a\in U(\mathfrak f)\\a\equiv1\pmod{\mathfrak d}}}
                  \rho(a)=N\mathfrak f.
\end{align*}
For every proper $\mathfrak d\mid\mathfrak f$, the reduction
$U(\mathfrak f)\to U(\mathfrak d)$ is surjective by local lifting and CRT.
Primitivity says $\rho$ is nontrivial on its kernel, so that inner sum
vanishes. For $\mathfrak d=\mathfrak f$ it equals one.
For unit $z$, substitution proves the second identity in~\eqref{mcs:gauss}.
For all frequencies together, finite additive Parseval gives
\[
 \sum_{z\bmod\mathfrak f}
   \left|\sum_x\rho(x)e_f(zx)\right|^2
       =N\mathfrak f\,\varphi(\mathfrak f).
\]
The unit frequencies already exhaust this quantity by the first identity;
every nonunit-frequency sum is therefore zero. This proves the asserted
zero extension. For $\mathfrak f=\mathcal O$ there is one residue class
and one character; take $\tau_1=1$ and the single frequency zero.

\subsection{The classical ordinary sieve in these lattice conventions}
For $D,Q\ge1$, coefficients $z_\alpha$ supported on any subset of
$\{\alpha\in\mathcal O:|\alpha|^2\le C_0D\}$ satisfy
\begin{equation}\label{mcs:ordinary}
 \sum_{N\mathfrak f\le Q}\ \sum_{\rho\ \mathrm{primitive}\bmod\mathfrak f}
 \left|\sum_\alpha z_\alpha\rho(\alpha)\right|^2
 \ll_{\mathcal O,C_0}(D+Q^2)\sum_\alpha|z_\alpha|^2.
\end{equation}
Here the primitive characters may be restricted to the unit-trivial
subfamily. We give the standard proof to make the conductor comparison
explicit.

First the reduced fractions $x/f$, with $x\in U(\mathfrak f)$ and
$N\mathfrak f\le Q$, are distinct modulo $\mathcal O$.
Their denominator ideal is intrinsic:
$\{a\in\mathcal O:a(x/f)\in\mathcal O\}=\mathfrak f$.
Thus equality modulo $\mathcal O$ forces the same denominator ideal
and residue numerator. For different fractions and every $k\in\mathcal O$,
\[
 \left|\frac{x}{f}-\frac{y}{g}-k\right|
   =\frac{|xg-yf-kfg|}{|fg|}\ge\frac1Q,
\]
because the numerator is a nonzero Eisenstein integer and
$|fg|=(N\mathfrak f N\mathfrak g)^{1/2}\le Q$.
The fixed codifferent factor $1/\delta$ changes only the lattice constants.
The modulus-one frequency zero obeys the same separation from all others.

For completeness, a $\rho_0$-separated set $t_r$ in a fixed
two-dimensional torus obeys the additive bound
\begin{equation}\label{mcs:additive}
 \sum_r\left|\sum_{|\alpha|\ll\sqrt D}
        z_\alpha\exp(2\pi i\operatorname{Tr}(\alpha t_r))\right|^2
       \ll(D+\rho_0^{-2})\sum_\alpha|z_\alpha|^2.
\end{equation}
Use fixed linear coordinates to write the lattice as $\mathbb Z^2$ and
its additive pairing as $e(n\cdot t)$. The support lies in
$|n_i|\le T$ with $T\asymp\max(1,\sqrt D)$.
To prove the dual inequality, majorize this box indicator by
$e^{2\pi}\exp(-\pi\|n\|^2/T^2)$, expand the square, and apply Gaussian
Poisson summation. Up to fixed constants the resulting kernel is
\[
 T^2\sum_{\ell\in\mathbb Z^2}
       \exp(-cT^2\|t_r-t_s+\ell\|^2).
\]
Take $\rho_0$ at most a fixed torus constant. The lifted points are
$c'\rho_0$-separated, so disjoint disks bound their count in a disk
of radius $L$ by $O(1+(L/\rho_0)^2)$.
Gaussian annular summation bounds every absolute kernel row sum by
$O(T^2+\rho_0^{-2})$. Schur's inequality and finite-dimensional duality
prove~\eqref{mcs:additive}, with constants independent of the number
of frequencies.

Put $\mathcal S(x/f)=\sum_\alpha z_\alpha e_f(\alpha x)$.
By~\eqref{mcs:gauss},
\[
 \rho(\alpha)=\tau_f(\overline\rho)^{-1}
       \sum_{x\in U(\mathfrak f)}\overline{\rho(x)}e_f(\alpha x).
\]
After taking squared absolute values, enlarging the primitive-character
sum to all characters of $U(\mathfrak f)$ and using orthogonality yields
\[
 \sum_{\rho\ \mathrm{primitive}\bmod\mathfrak f}
       \left|\sum_\alpha z_\alpha\rho(\alpha)\right|^2
 \le\frac{\varphi(\mathfrak f)}{N\mathfrak f}
       \sum_{x\in U(\mathfrak f)}|\mathcal S(x/f)|^2
 \le\sum_{x\in U(\mathfrak f)}|\mathcal S(x/f)|^2.
\]
Apply~\eqref{mcs:additive} with $\rho_0\gg Q^{-1}$ to prove
\eqref{mcs:ordinary}.

Finally $\mathcal O$ has class number one. Choose one generator
$\alpha_m$ of each ideal column, with $|\alpha_m|^2=Nm$; these are an
injective subset of the lattice disk. A primitive trivial-infinity-type
Hecke character is a unit-trivial primitive residue character on
$\mathcal O/\mathfrak f$, and its ideal value equals its value on
$\alpha_m$. Descent of the residue character to a proper divisor would
preserve unit-triviality and contradict primitive Hecke conductor.
Thus~\eqref{mcs:ordinary} applies to arbitrary ideal coefficients of
norm at most $2D$, including squarefree ideals of nonsquarefree rational norm.

\subsection{Raw zeros, tuple multiplicity and block transfer}
Freeze $r_6$ and the literal unit/$S$-factor, putting
\[
 c_m=b_m\chi_m(u_0)\mathbf1_{(m,r_6)=1},\qquad
       \sum_m|c_m|^2\le\sum_m|b_m|^2.
\]
This is independent of every first-five variable. In particular the
arbitrary overlapping sixth-power factor is not replaced by one.
Let $v=\prod_{j=1}^5r_j^j$ denote the chosen literal product.
Every outside-$S$ valuation is between one and five.
BGL, \href{https://arxiv.org/pdf/1112.1650v1}{arXiv:1112.1650v1},
Lemma~2.1, together with the established finite $S$-unit comparison
in~\eqref{wc:literal-reciprocity}, gives its primitive literal Kummer
character $\kappa_v$ of trivial infinity type, with
\[
 q=r_1\cdots r_5\mid\operatorname{cond}(\kappa_v)\mid\mathfrak c_Sq,
 \qquad \kappa_v(m)=\chi_m(v)\quad((m,S)=1).
\]
The equality includes collisions with $q$: both sides are zero.
The fixed $\mathfrak c_S$ includes all wild conductor types and literal
comparison twists. No row-dependent phase enters the coefficients.

A fixed primitive character determines its outside-$S$ conductor $q$.
Each prime of $q$ has at most five assignments to the pairwise-coprime
$r_1,\ldots,r_5$, so tuple multiplicity is at most
$5^{\omega(q)}\ll_\eta(Nq)^\eta$.
All conductors in the dyadic block have norm at most $C_SQ$.
Enlarge the positive row sum, apply~\eqref{mcs:ordinary}, and count
$O(A_6)$ possibilities for the frozen arbitrary $r_6$.
This proves~\eqref{mcs:block-bound}; the actual capacity, squarefree,
ray and inter-factor restrictions only shrink the positive sum.
Equivalently $\mathcal E_A(b)\le\Lambda_A\sum_m|b_m|^2$;
there is no coefficient norm factor in the definition of $\Lambda_A$.

\subsection{Uniform coefficient and principal-residue assembly}
The other block estimate is exactly the accepted powered-BGL input
\eqref{wc:block}, whose squarefree powers are justified in BGL
Section~3.2. The capacity comparisons and elementary minimum in
Section~\ref{mcs:core} give
\[
 \Lambda_A\ll(HD)^\varepsilon
          \{H+DH^{1/6}+(DH)^{2/3}\}.
\]
Sum logarithmically many blocks with smaller internal exponents.
An $S$-supported literal factor $\mathfrak b$ reduces capacity to
$H/N\mathfrak b$; the required sums over $\mathfrak b$ converge
at powers $1$, $1/6$, $2/3$. Literal units have multiplicity at most six.
This proves the raw bound with arbitrary fixed $b_m$.

For the actual coefficients use~\eqref{wc:l2}:
$|A_U(m)|\le3^{\omega(m)}\ll_\eta(Nm)^\eta$ gives
\[
 b_m=D^{-1/2}A_U(m)\nu(m)W(Nm/D),\qquad
       \sum_m|b_m|^2\ll_\eta\|W\|_\infty^2D^\eta,
\]
uniformly for $U\ge1$, $U^2\le D$. The common $W$ is not optimized
separately for different rows. If the primitive $\psi_u$ is principal,
$\nu$ is unramified outside $S$, so tame Kummer inertia forces every
outside-$S$ valuation of $u$ to be divisible by six. Ideal counting and
the convergent $S$-supported sum at power $1/6$ give
$O_{\nu,S}(H^{1/6})$ such literal rows.
In the exact formula~\eqref{wc:principal}, $|\rho_T|\ll_K1$,
$|\widehat W(1)|\le\|W\|_\infty$, and
\[
 \sum_d\frac{|\beta_U(d)|}{Nd}
 \le\left(\sum_{Na\le U}\frac1{Na}\right)^2
 \ll\log^2(2U).
\]
Thus $|P_u|\ll\sqrt D\|W\|_\infty\log^2(2U)$ on principal rows,
and
\[
 \sum_{0<Nu\le H}|P_u|^2
       \ll DH^{1/6}\|W\|_\infty^2\log^4(2U).
\]
This fits the new envelope. The exact identity $E_u=I_u-P_u$ and
$|E_u|^2\le2|I_u|^2+2|P_u|^2$, with smaller internal exponents,
prove~\eqref{mcs:uniform}. No cutoff limit or profile supremum occurs.

The ordinary comparison is an application refinement. A direct BGL
radical relabelling would additionally require reciprocity of its mixed
family, not just its conductor count: for pairwise-coprime $a,b,a',b'$,
the quotient after the original fixed ray phases is
$\vartheta_b(a')/\vartheta_a(b')$, which varies outside $S$.
GL, \href{https://arxiv.org/pdf/1112.1642v2}{arXiv:1112.1642v2},
Definition~1(2), is therefore not supplied by that shortcut.
This does not rule out a different mixed-family theorem; it is simply
unneeded for the proved weaker ordinary comparison.
The remaining exponent gap and every independent signed, spectral,
formal and GRH obligation stated in Section~\ref{mcs:core} remain open.

% Original finite arithmetic learning construction; no asymptotic or GRH claim.
\section{Learning the actual signed arithmetic correlations}\label{arl:core}
The unsigned refinement in Section~\ref{mcs:core} leaves a signed problem.
Here is a separate, executed finite learning objective on the original
correlations, with a certificate for every coordinate and observable.
It is not a uniform estimate for an expanding family.
For ordered distinct squarefree columns $m=hr,n=ht$, put $q=rt$ and
\[
 \Gamma_i=\sum_{u\in\mathcal R}\chi_r(u)\overline{\chi_t(u)}
                         \mathbf1_{(u,h)=1},\qquad i=(h,q,r).
\]
The symbols vanish at every collision with $q$; the common-prime mask is
retained separately. With principal $\nu$, actual complementary divisor counts
$r_U(m)$, and one common complex profile $W_m=W(Nm/D)$, define
\[
 d_i^{CC}=r_U(m)r_U(n)W_m\overline{W_n},\qquad
 d_i^{C\mu}=r_U(m)W_m\overline{W_n},\qquad
 \lambda_i=1+|d_i^{CC}|^2+|d_i^{C\mu}|^2.
\]
The exact diagonal $\Delta_X$ and these offdiagonal coordinates give
\begin{equation}\label{arl:observable}
 E_X=\Delta_X+D^{-1}\sum_i\mu(q_i)d_i^X\Gamma_i,
       \qquad X\in\{CC,C\mu\}.
\end{equation}
Thus $E_{CC}=\sum_{u\in\mathcal R}|C_u|^2$ and
$E_{C\mu}=\sum_{u\in\mathcal R}C_u\overline{M_u}$ for the actual normalized
row sums. Neither the M\"obius sign nor a complex conjugate is discarded.

\begin{proposition}[A finite supervised certificate]\label{arl:certificate}
Let $g_i$ be complex readout coordinates and
$L(g)=\sum_i\lambda_i|g_i-\Gamma_i|^2$.
For $K_X=\sum_i|d_i^X|^2/\lambda_i$ and
$\lambda_{\min}=\min_i\lambda_i$, the substituted readouts satisfy
\begin{align*}
 |\widehat E_X-E_X|^2&\le D^{-2}K_X L(g),\\
 \max_i|g_i-\Gamma_i|^2&\le L(g)/\lambda_{\min}.
\end{align*}
The executed metric-gradient update $g^{(t+1)}=(g^{(t)}+\Gamma)/2$
gives $L(g^{(t)})=4^{-t}L(g^{(0)})$. Therefore every positive rational
requested tolerance on these three squared bounds is attained after finitely
many updates on this fixed exact instance, with unrestricted resources.
\end{proposition}
The implementation's resource guard refuses requests beyond its update budget;
termination in this proposition is not a practical runtime guarantee.
The labels $\Gamma_i$ are computed exactly before training, so no speedup over
exact label evaluation or generalization to unseen conductors is claimed.

The supplied instance uses $D=6000,U=60,H=625$, $S$ above $2,3$, principal
$\nu$, and all $15$ good squarefree ideal columns at activated norms
$7507,7519,8281,10507$ (respectively $2,4,1,8$ ideals).
One explicitly specified common smooth four-bump profile is zero at every
other integer norm. Both the $2262$ literal rows and the separately labelled
$2148$ guaranteed nonprincipal subset are evaluated; the latter is not
asserted to be the whole nonprincipal family.
There are $95$ offdiagonal $(h,q)$ blocks, $65$ with $h=1$, and $210$
ordered learned coordinates. Every row case retains $4530$ raw column zeros.
Real and complex profiles give four requested-$10^{-12}$ certificates,
using $53,55,53,55$ metric updates respectively: $216$ executed updates,
with every full rational state saved. A separate $200$ executed ordinary
coordinate updates are budget-limited and carry no requested-error claim.
The standalone replay is linked from the book and project page.
No analytic saving, positive determinant, zero-location assertion,
upstream quantum training, new Lean declaration or GRH conclusion follows.

\section{Finite phase features, gradients and exact verification}\label{arl:details}
For each block let $G_q=(\mathbb Z/6\mathbb Z)^{\omega(q)}$.
A literal admissible row determines local phases $e_p$ at every prime of $q$;
rows with a $q$-collision or a common-prime collision with $h$ are deleted
before forming the exact sparse histogram $n(e)$. Orientation $r$ has
\[
 \theta_r(e)=\zeta_6^{\sum_{p\mid q}s_p e_p},\qquad
 s_p=1\ (p\mid r),\quad s_p=-1\ (p\mid t),\qquad
 \Gamma_r=\sum_e n(e)\theta_r(e).
\]
Different active orientations have distinct exponent vectors modulo $6$,
so their characters are orthonormal for the uniform measure on $G_q$.
The model and target projection are
\[
 F_g(e)=\sum_{r\in J_q}g_r\overline{\theta_r(e)},\qquad
 F_\Gamma(e)=\sum_{r\in J_q}\Gamma_r\overline{\theta_r(e)}.
\]
Indeed $F_\Gamma$ is the orthogonal projection of $|G_q|n$ onto this
character subspace, and $\mathbb E_{G_q}F_g\theta_r=g_r$.
The positive spectral metric assigns weight $\lambda_r$ to this orthonormal
basis. Its squared error is exactly $\sum_r\lambda_r|g_r-\Gamma_r|^2$.
Thus the run evaluates the full finite feature-projection loss by Parseval;
it does not claim to have executed a sampled or stochastic training loss.
This projection need not be a nonnegative density. A Hermitian readout matrix
need not be positive semidefinite.

\paragraph{Gradient and contraction.}
Use $\zeta_6$ coordinates $g=a+b\zeta_6$, so
$|g|^2=a^2+ab+b^2$ and
$M=\left(\begin{smallmatrix}1&1/2\\1/2&1\end{smallmatrix}\right)$.
For the objective $L/2$, the ordinary real-coordinate gradient in block $i$
is $\lambda_i M(x_i-x_i^*)$.
In the block metric $\lambda_iM$, the gradient is $x_i-x_i^*$.
A metric step of size $1/2$ therefore halves every complex error and quarters
$L$. Conjugate initial coordinates and this step preserve the Hermitian
reverse orientations. The ordinary trajectory treats all ordered coordinates
independently; asymmetric weights need not preserve that symmetry.
For ordinary steps $\eta=2/(3\lambda_{\max})$, the Hessian eigenvalues
lie between $\lambda_{\min}/2$ and $3\lambda_{\max}/2$, whence
\[
 L_{t+1}\le\left(1-\frac{\lambda_{\min}}{3\lambda_{\max}}\right)^2L_t.
\]
This last bound is conservative; the saved $50$ ordinary steps per case are
not asserted to attain the requested tolerance.

\begin{proof}[Proof of Proposition~\ref{arl:certificate}]
Subtract~\eqref{arl:observable} with the same exact diagonal. Weighted
Cauchy--Schwarz, applied only to the error, gives
\[
 \left|\sum_i\mu(q_i)d_i^X(g_i-\Gamma_i)\right|^2
 \le\left(\sum_i\frac{|d_i^X|^2}{\lambda_i}\right)
     \left(\sum_i\lambda_i|g_i-\Gamma_i|^2\right).
\]
Here $|\mu(q_i)|=1$ because $q_i$ is squarefree; the predicted sum itself
keeps its signs and phases. Each summand in the second factor is nonnegative,
which proves the coordinate bound. The exact affine update gives
$g^{(t)}-\Gamma=2^{-t}(g^{(0)}-\Gamma)$, proving contraction.
For finite rational positive $\lambda_i,K_X,L_0$ and $\tau>0$, some integer
$t$ makes all three bounds at most $\tau^2$. The zero-error case stops
immediately. No varying-family constants enter this finite assertion.
\end{proof}

\paragraph{The activated original family.}
Let
\[
 \phi(t)=\begin{cases}\exp(1-1/(1-t^2)),&|t|<1,\\0,&|t|\ge1,\end{cases}
 \qquad W(y)=\sum_{N\in\{7507,7519,8281,10507\}}
                     w_N\phi(4(6000y-N)).
\]
These disjoint bumps are supported in $(1,2)$ and $\phi(0)=1$.
For the real case $w_N=1$; for the complex case
$w_N=1+\zeta_6(4N-30000)/6000$.
At integer norms $W(n/6000)=w_n$ at the four activated norms and is zero
elsewhere. Hence all active original squarefree columns are present, rather
than an undeclared column submatrix. No exponential evaluation is needed
at these exact nodes. The profile is a fixed finite demonstration and does
not establish a profile-uniform analytic estimate.
The counts $r_U$ are $0,2,4,2$ at the four norms.
The guaranteed nonprincipal subset removes precisely those literal rows whose
outside-$S$ norm is $1$; no classification of the omitted $114$ rows as a
whole is asserted. Every nonunit outside-$S$ factor has a prime exponent
from $1$ to $3$, since its prime norm is at least $7$ and $7^4>625$;
its nonzero residue modulo $6$ gives a nontrivial local character.
All six units and all original $S$-power rows belong to
the separately saved full literal case.

The seed $20261105$ selects deterministic rational conjugate initial
coordinates by SHA256 of the ordered column identifiers. There is no random
sampling. Exact $\mathbb Q(\zeta_6)$ arithmetic stores every metric and
ordinary state. Training accepts only when both observable certificates and
the uniform-coordinate certificate are at most $\tau^2$, preventing a
scalar cancellation from masquerading as accurate correlation learning.
The real full-row targets are $E_{CC}=2393/125$ and $E_{C\mu}=431/50$;
the guaranteed-subset targets are $4293/250$ and $771/100$.
Complex-profile targets and the complete phase histograms are saved alongside.
Each case has $64$ nonzero negative-$\mu(q)$ C--C target orientations.
Replacing all signs by plus is an explicit failing control, not an estimator.

\paragraph{Reproduction and scope.} The complete run is reproducible.

The source is \texttt{examples/}\allowbreak\texttt{arithmetic\_learning/}\allowbreak\texttt{learn.py}; its README
states the exact smooth profile, CLI, resource guard and loss convention.
The accepted immutable output is
\texttt{outputs/experiments/}\allowbreak\texttt{2026-11-05-arithmetic-learning-v2/}.
All source/helper hashes, exact input bytes and artifact hashes are bound
before and after execution and checked on replay. Existing outputs,
nonpositive or binary-float tolerances, and stale source/input/case bytes
are refused. The preserved v1 logging failure is not the accepted run.
The focused verifier independently reconstructs $31668$ local symbols and
$450$ Gram targets, then checks the full saved metric trajectory.
Independent mathematical review and exact code checks are recorded separately;
finite arithmetic checks are not a proof of an asymptotic signed estimate.
Full exact training labels are supervision, so this is a certified finite
reference architecture, without an efficiency, generalization, PSD,
zero-location or GRH claim. Analytic and categorical Lean remain unfinished.

% New scoped consequence for moving common profiles; no fixed-profile saving.
\section{Why the common profile must stay fixed}\label{prf:core}
The unsigned envelope of Section~\ref{mcs:core} and the unresolved signed
estimate have different profile quantifiers. We give an obstruction to an
$H$-sized bound uniform through the sup norm alone. It uses the \emph{actual}
coefficients and nonprincipal rows, rather than replacing the columns by
arbitrary weights. It does not contradict an estimate for one fixed smooth
profile or one with adequate derivative dependence.

Fix $K=\mathbb Q(\omega)$, $\omega^2+\omega+1=0$, principal $\nu=1$, and a
finite good set $S$ containing the primes above $2,3$. Write
$\zeta=1+\omega$, $\delta=1+2\omega$, so $\zeta$ has order six and
$\delta^2=-3$. Retain the literal raw symbols and normalization
\[
 I_u(W)=D^{-1/2}\sum_{\substack{m\ \mathrm{squarefree}\ (m,S)=1}}
 A_U(m)\chi_m(u)W(Nm/D),\qquad E_u=I_u-P_u.
\]
Here and below squarefree means \emph{ideal} squarefree. For $U^2=D$ the
accepted high-window identity gives
$A_U(m)=\mu(m)(r_U(m)-1)$ on $D<Nm<2D$, with $r_U(m)$ even and
$|A_U(m)|\ge1$. The multiset of prime-ideal norms of a squarefree ideal
is determined by its rational norm. Thus $A_U(m)=A_N$ depends only on
$N=Nm$, including divisor-norm multiplicities. Moreover
\begin{equation}\label{prf:unit}
 \eta(m):=\chi_m(\zeta)=\zeta^{(Nm-1)/6},\qquad (m,S)=1.
\end{equation}
The primitive character $\eta$ has conductor $(2)^2(\delta)^3$, of norm
$432$; its powers $1,5$ have conductor norm $432$, powers $2,4$ norm $27$,
and power $3$ norm $16$. All five powers are nonprincipal. These row
conductors do not remove the raw zero at any prime dividing a sixth factor.

Put $\phi(t)=\exp(1-1/(1-t^2))$ for $|t|<1$ and zero otherwise.
For integer $D$ set $w_N=\operatorname{sgn}(A_N)\overline{\zeta^{(N-1)/6}}$
when a good squarefree ideal of norm $N$ exists and $D<N<2D$, and zero
otherwise. Define one common, complex profile
\begin{equation}\label{prf:profile}
 W_D(y)=\sum_{D<N<2D}w_N\phi(4(Dy-N)).
\end{equation}
Its disjoint supports lie strictly inside $(1,2)$, its sup norm is one for
large $D$, and $\|W_D^{(j)}\|_\infty\ll_jD^j$. In particular the endpoints
have value zero. This is a profile varying with scale, independent of the
row, not a separate choice for each ideal or each row.

\begin{proposition}[A moving-profile lower bound]\label{prf:lower}
For every fixed $0<\theta<1$, take $H=D^\theta$, $U=\sqrt D$, and
$X=H^{1/6}$. As integer $D\to\infty$,
\begin{equation}\label{prf:lower-bound}
 \sum_{0<Nu\le H}|E_u(W_D)|^2\ge(c_S+o(1))DH^{1/6},\qquad c_S>0.
\end{equation}
This lower bound already comes from literal rows
$u_v=\zeta\beta_v^6$, where $v$ is an arbitrary good ideal of norm at most
$X$ and $\beta_v$ is a generator. These rows have nonprincipal primitive
character $\eta$, so $P_{u_v}=0$. Their exact values are
\begin{equation}\label{prf:row}
 E_{u_v}(W_D)=D^{-1/2}
 \sum_{\substack{m\ \mathrm{squarefree},\ (m,Sv)=1\\D<Nm<2D}}
                      |A_U(m)|.
\end{equation}
Every collision with $v$ remains deleted, even when the primitive conductor
drops to $432$. No squarefree condition is imposed on $v$.
\end{proposition}
At $H=D^{3/4}$ the lower exponent is $9/8$, whereas the current upper
exponent is $7/6$. Thus this does not establish sharpness of $7/6$ or a new
upper saving. For $0<\theta<2/3$, the term $DH^{1/6}$ dominates the current
unsigned envelope; it is necessary up to subpower losses in this
\emph{varying-profile} class, even on nonprincipal rows.

For $k\ge0$ put
$M_k(W)=\sum_{j=0}^k\|(y\partial_y)^jW\|_\infty$.
If, for every $\varepsilon>0$, a bound uniform over all common profiles
with constants independent of $W,D$ were to assert
\[
 \sum_{0<Nu\le H}|E_u(W)|^2
 \ll_\varepsilon M_k(W)^\alpha(HD)^\varepsilon H,
 \qquad \alpha\ge0,
\]
then necessarily $k\alpha\ge1-5\theta/6$. At $\theta=3/4$ this is
$k\alpha\ge3/8$. For a single fixed $\varepsilon$ the necessary condition
is only $k\alpha+(1+\theta)\varepsilon\ge1-5\theta/6$.
These are necessary cost comparisons, not sufficient estimates. A
fixed-$W$ result, a class with bounded seminorms, or the existing
$M_4(W)^2$ dependence is not refuted.

The same rows expose a concrete centering error in the joint completion:
removing the orientation mean of the actual correlation can delete a
nonzero signed contribution on \emph{nonprincipal} rows. The extended proof
retains the literal unit phase, gives an exact four-orientation example,
and explains why orientation projection alone supplies no saving. The
standalone \href{../profile-uniformity.html}{profile and coherent-mean inspector}
shows the complete finite norm groups, arbitrary sixth-factor masks and
exact residual quantities. Its finite tables are diagnostics of this
construction, not asymptotic estimates, training or GRH evidence.

% Exact proof with fixed principal nu, moving profile, and all raw masks.
\section{Proof and joint-mean boundary of the moving-profile construction}\label{prf:details}
\subsection{Norm groups and primitive conductors}
For a split rational prime $p$, a squarefree ideal contains zero, one or
both conjugate primes. Its rational norm exponent is respectively $0,1,2$;
at exponent one there are two choices, and at exponent two its divisor
norms include two copies of $p$. For an inert prime, the exponent is zero
or two, and the nontrivial choice is the single prime ideal of norm $p^2$.
Thus the rational norm determines $\mu$, the divisor-norm multiset and
$r_U$, even though it need not determine the ideal. For example norm $49$
has squarefree ideal divisor norms $1,7,7,49$, whereas the inert prime of
norm $25$ has divisor norms $1,25$. Additional fixed $S$ deletions remove
some ideals without changing this constancy among those remaining.
The high-window formula and its parity are the already proved facts of
Section~\ref{sb:core}, not a new coefficient estimate.

At a good prime ideal $\mathfrak p$ of norm $q$, Euler's definition gives
$\chi_{\mathfrak p}(\zeta)=\zeta^{(q-1)/6}$. All such norms are $1$ modulo
$6$; for $a=1+6r,b=1+6s$ one has
$(ab-1)/6=r+s+6rs$. Multiplicativity proves~\eqref{prf:unit}.
Write $\chi_4(N)=(-1)^{(N-1)/2}$ for odd $N$. On good ideals,
\begin{align*}
 \eta&=\chi_4(N)\omega^{(N-1)/3},\\
 \eta^2&=\omega^{2(N-1)/3}=\omega^{-(N-1)/3},\qquad
 \eta^3=\chi_4(N).
\end{align*}
The inverse orientation in the middle expression is essential.

For completeness the local conductors are exact, not merely upper bounds.
The quadratic norm character factors modulo $(2)^2=(4)$.
Since $\operatorname{Tr}(\delta\gamma)\in3\mathbb Z$ for every
$\gamma\in\mathcal O$, expanding the norm gives
$N(\alpha+\delta^3\beta)\equiv N\alpha\pmod9$.
The cubic norm character therefore factors modulo $\delta^3$.
The coprime local witnesses
\begin{align*}
 \alpha_2&=13+6\omega,\quad N\alpha_2=127,\quad
 \alpha_2-1=2\omega\delta^3,\\
 \alpha_3&=-11,\quad N\alpha_3=121,\quad
 \alpha_3-1=4\delta^2.
\end{align*}
show minimality. The first is $1$ modulo $2$, is $1+2\omega$ modulo $4$,
and is $1$ modulo $\delta^3$; its $\eta$ value is $-1$.
The second is $1$ modulo $4$ and $-2=1+\delta^2$ modulo $\delta^3$,
but not $1$ modulo $\delta^3$; its $\eta$ value is $\omega$ and its square
is $\omega^2$. Thus the quadratic local exponent is two and the cubic
one is three. The two moduli are coprime, giving the stated norms
$432,27,16$ for the nontrivial powers. Global units have norm one, so
these residue characters are unit trivial and descend to ideal Hecke
characters in this class-number-one field. A norm-seven prime has value
$\zeta$, proving order six. A larger finite $S$ does not principalize the
character: CRT gives representatives of the same nontrivial ray classes
avoiding the extra excluded primes. Primitive conductor reduction never
restores raw primes excluded by $S$ or the sixth factor.

\subsection{Exact alignment, with arbitrary overlapping sixth factors}
The bumps in~\eqref{prf:profile} have radius $1/(4D)$ and center $N/D$.
Distinct supports are disjoint. At an active integer node their value is
exactly $w_N$, and all endpoints are zero. Differentiation gives
$\|W_D^{(j)}\|_\infty\le(4D)^j\|\phi^{(j)}\|_\infty$.
Expanding $(y\partial_y)^j$ into ordinary derivatives and using $1<y<2$
gives $M_k(W_D)\ll_kD^k$.

Every nonzero ideal $v$ has a generator $\beta_v$. The literal element
$u_v=\zeta\beta_v^6$ is independent of that choice, since all units have
sixth power one; distinct ideals yield distinct rows. Its norm is $(Nv)^6$.
For squarefree good $m$, exact raw multiplicativity says
\[
 \chi_m(u_v)=\eta(m)\mathbf1_{(m,v)=1}.
\]
Repeated primes of $v$ are allowed and its collisions with $m$ remain
zero. The primitive row character is always $\eta$, so the actual
principal indicator in~\eqref{wc:principal} is zero when $\nu=1$.
Substituting the nodal profile proves~\eqref{prf:row} with no factor six,
unit averaging or residue replacement.

\subsection{The coprime-triple density and the selected $h=1$ block}
Put $\kappa=\pi/(3\sqrt3)$. Lattice counting, with six unit generators
per ideal, gives
\[
 A_K(Y):=\#\{a:Na\le Y\}=\kappa Y+O(\sqrt Y+1),
 \qquad A_K(Y)\le C Y\quad(Y\ge1).
\]
This is the existing elementary ideal count. Let $T(D,X)$ count ordered
triples $(m,n,v)$ with $m,n$ squarefree good, $D<Nm,Nn<2D$, $Nv\le X$,
$v$ good and arbitrary, and the three ideals pairwise coprime. For
$D,X\to\infty$, $X\le D$, we claim
\begin{equation}\label{prf:triple}
 \begin{aligned}
 T(D,X)&=(\kappa^3\mathcal C_S+o(1))D^2X,\\
 \mathcal C_S&=\prod_{\mathfrak p\in S}(1-t_{\mathfrak p})^3
 \prod_{\mathfrak p\notin S}
       (1-5t_{\mathfrak p}^2+6t_{\mathfrak p}^3-2t_{\mathfrak p}^4)>0,\\
 t_{\mathfrak p}&=(N\mathfrak p)^{-1}.
 \end{aligned}
\end{equation}
Here a product over $S$ means its prime ideals. To prove this, first impose
the local restrictions at a fixed finite set of primes containing $S$.
They are finite Boolean combinations of divisibility by $\mathfrak p$
and $\mathfrak p^2$ in $m,n$, and by $\mathfrak p$ in $v$.
Finite inclusion-exclusion and the count for fixed ideal multiples give
the product of the local densities, with relative error tending to zero.
At a good prime the allowed cases are no positive valuation, a single
valuation one in $m$ or $n$, or an arbitrary positive valuation in $v$.
The factor is therefore
\[
 (1-t)^3\left(1+2t+\frac{t}{1-t}\right)
   =(1-t)^2(1+2t-2t^2)=1-5t^2+6t^3-2t^4.
\]
At $S$ all three valuations must vanish, giving $(1-t)^3$.
The geometric series for $v$ is essential; $v$ is not squarefree.

An omitted failure belongs to one of five events:
$\mathfrak p^2\mid m$, $\mathfrak p^2\mid n$,
$\mathfrak p\mid(m,n)$, $\mathfrak p\mid(m,v)$ or
$\mathfrak p\mid(n,v)$. Divisibility counts and $A_K(Y)\le CY$ bound
each by $O(D^2X/(N\mathfrak p)^2)$, uniformly; if a quotient scale is
below one its count is zero. Thus the tail after norm $P$ is
\[
 O\left(D^2X\sum_{N\mathfrak p>P}(N\mathfrak p)^{-2}\right)=o_P(D^2X).
\]
The series converges by the ideal count (or $\zeta_K(2)<\infty$).
Take the large-scale limit first at fixed $P$, then $P\to\infty$.
Each local factor is positive and differs from one by $O(t^2)$, proving
\eqref{prf:triple}. Strict annulus endpoints have no effect on the main
term by the same lattice count. The secondary range $X/2\le Nv<X$ has
half this limiting volume, though it may be empty at small $X$.

On the selected literal unit class every expanded pair term in
\eqref{prf:row} is nonnegative. Keeping only $(m,n)=1$ gives its exact
$h=1$ A--A contribution
\[
 \mathcal E^{AA}_{\mathrm{selected},h=1}
 =D^{-1}\sum_{\substack{v\ \mathrm{good}\ Nv\le X}}
  \sum_{\substack{m,n\ \mathrm{squarefree,good}\ D<Nm,Nn<2D\\
                   (m,n)=(m,v)=(n,v)=1}}|A_U(m)A_U(n)|
 \ge D^{-1}T(D,X).
\]
Consequently, with $X=H^{1/6}$,
\[
 \sum_{0<Nu\le H}|E_u|^2
 \ge\sum_{Nv\le X,\ v\ \mathrm{good}}|E_{u_v}|^2
 \ge\mathcal E^{AA}_{\mathrm{selected},h=1}
 \ge(\kappa^3\mathcal C_S+o(1))DX.
\]
The first inequality is positivity of the \emph{outer squared row sum}.
The selected $h=1$ block does not lower-bound the signed $h=1$ sum over
other literal rows or a unit-complete family. No positive-semidefinite
claim about a gcd-truncated Gram matrix is needed. This is A--A, not a
positive-density assertion for C--C long-long weights, which may vanish.
This proves Proposition~\ref{prf:lower} with $c_S=\kappa^3\mathcal C_S$.
The derivative tradeoff follows by comparing powers in
$DH^{1/6}\ll D^{k\alpha+(1+\theta)\varepsilon}H$ and then letting
$\varepsilon\downarrow0$ only when the proposed bound holds for every
positive $\varepsilon$.

\subsection{Projection preserves the coupled problem, not a saving}
Use the exact joint notation of Section~\ref{jdq:details}. For fixed
$h,q$ with $q\ne1$ (including $h=1$), let $n=\varphi(q)$, $\ell=N(hq)$, $Q=Nq$, and
\begin{align*}
 c&=\frac{Q^2\xi_q(-\alpha)\tau_\beta(\xi_q)}{\ell^6},\\
 \Gamma_r^F&=c\sum_K B_F(K)\overline{\theta_r(K)},\qquad
 A(K)=\sum_{r\in J}a_r\overline{\theta_r(K)}.
\end{align*}
The active orientations $J$ include all original restrictions, and the
literal unit phase belongs to $a_r$. The distinct $\theta_r$ are
orthonormal for $n^{-1}\sum_K$. If $P$ projects onto their span, then
\[
 PB_F=(cn)^{-1}\sum_{r\in J}\Gamma_r^F\theta_r,\qquad
 \sum_K|PB_F(K)|^2=(|c|^2n)^{-1}\sum_r|\Gamma_r^F|^2.
\]
Replacing $B_F$ by $PB_F$ in the full bilinear pairing is exact.
Cauchy--Schwarz is precisely
$|c\sum_KB_FA|\le\|a\|_2\|\Gamma^F\|_2$.
The conductor and group-size factors cancel. Projection rank alone gives
no arithmetic saving for this particular coupled vector.

For $s=\omega(q)$ let $C=\{K:x_{\mathfrak p}(K)^2=1\text{ for all }
\mathfrak p\mid q\}$. All orientations agree on $C$. If $M_C$ denotes
multiplication by its indicator, the matrix in the orthonormal orientation
basis is
\[
 PM_CP=3^{-s}\mathbf1\mathbf1^*,\qquad
 \lambda=|J|/3^s,\qquad
 \|(1-P)M_CP\|^2=\lambda(1-\lambda).
\]
Thus coherent restriction does not commute with projection. For example
\[
 \left|\sum_C(B_F-PB_F)A\right|
 \le\sqrt{\lambda(1-\lambda)}\|(1-P)B_F\|_2\|A\|_2,
\]
where all norms in this inequality are counting $\ell^2$ norms. Normalized
norms would require an additional factor $n$ on the right. Counting a
coherent subset alone is not a lower bound for its actual signed pairing.

\subsection{A nonzero orientation mean on actual nonprincipal rows}
For the selected rows let $V$ be any finite set of good ideals, put the
first five row factors equal to one, and retain the sixth-factor mask.
For each $h,q$ define $Z_{h,q}=\#\{v\in V:(v,hq)=1\}$.
After removing the separately displayed literal unit phase, the exact
correlation is
\[
 \Gamma_r^F=Z_{h,q}
    =c\sum_KB_V(K)\overline{\theta_r(K)}
 \quad\text{for every active orientation }r.
\]
To pass from the published C--C identity to A--A replace each $r_U$ by
$r_U-1$. At $h=1$, $m=r,n=q/r$, the coefficient is
\[
 a_r^{AA}=(r_U(m)-1)(r_U(n)-1)W_m\overline{W_n}
                                   \theta_{m,n}(\zeta),
 \qquad \mu(q)a_r^{AA}=|A_U(m)A_U(n)|.
\]
This retains $\mu(q)$, the common complex profile and its conjugation,
all deleted common primes, and the literal $\zeta$ phase.
If $Z_{1,q}>0$ and at least one active interior A--A coefficient is
nonzero, the actual pairing is positive, whereas
$\sum_r|\Gamma_r^F-\operatorname{mean}\Gamma^F|^2=0$.
An orientation-variance-only replacement with no mean term is therefore
false. The mean is not a principal residue: these rows have character
$\eta$, with nonprincipal powers. This does not refute a properly
formulated fixed-profile signed estimate.

A fresh exact example uses $D=4000$, $U=63$ and chosen prime ideals of
norms $7,13,79,67,97$. Their product has four active orientations:
the pairs $6097/7663$, $7189/6499$ and both reversals.
Here $r_U=2$, and $A_U$ is $-1$ on the three-prime columns and $+1$ on
the two-prime columns. The four good ideals with $10\le Nv\le20$ have
norms $13,13,19,19$. Exactly one collides with the chosen $13$-prime,
so $Z=3$ for all four orientations. The exact selected $h=1$ contribution
is $12/4000=3/1000$ despite zero orientation variance.
Dropping the literal $\zeta$ phase incorrectly changes the coefficient
sum from $4$ to $-1$, giving $-3/4000$.
The coherent fraction is $4/243$ and the squared leakage is $956/59049$.
The residue group and infinite Fourier tail are not estimated by this
finite check. The fixed row conductor $432$ must be distinguished from
the pair-column conductor $q=mn$, whose norm here and in the asymptotic
annulus is comparable to $D^2$.

The exact implementation at \texttt{examples/\allowbreak profile\_uniformity/}
keeps complete norm groups and arbitrary sixth factors, including distinct
ideals of norm $49$ with different masks. The saved six-scale examples,
exact joint fixture and their independent checks certify finite identities
only. Full-union masks, coherent means, fixed-profile uniformity and the
remaining mixed/pure upper-envelope neighborhoods still require a new
signed estimate. No determinant comparison, zero location, analytic or
categorical Lean theorem, quantum result or GRH proof is obtained.

% Actual signed gcd-one budget and character identification; no power saving.
\section{A bound for the signed coprime slice and an exact character twist}\label{sw:core}
The coprime-column contribution can be negative, so it needs its own bound.
Fix the original $\nu,S$, one common possibly complex smooth $W$ supported
in $[1,2]$, $D>1$, $U\ge1$ and $U^2\le D$. Consider an actual block
\[
 u=u_0v_1v_2^2,\quad A_j\le Nv_j<2A_j\ (j=1,2),\quad Nu\le H,
 \qquad R=A_1A_2,\quad M=\max(A_1,A_2).
\]
The good ideals $v_1,v_2$ are squarefree and coprime, with their literal
primary/ray representatives; all other good factors are the unit ideal.
Let $F$ be the actual indicator, including capacity and every row mask.
In the remaining neighborhoods the first factor is nonunit, so the
primitive row character is nonprincipal and $E_u=I_u$ there.
For $X,Y\in\{A_U,C_U,\mu\}$ define the actual signed slice
\begin{equation}\label{sw:slice}
 S_1^{XY}=\frac1D\sum_uF(u)
 \sum_{\substack{m,n\ \mathrm{squarefree,good}\\(m,n)=1}}
 X(m)Y(n)\nu(m)\overline{\nu(n)}\chi_m(u)\overline{\chi_n(u)}
 W(Nm/D)\overline{W(Nn/D)}.
\end{equation}
All columns in this formula are the original annular columns, including
raw zero sentinels. No positive-semidefinite claim about the coprime
truncation is made.

\begin{proposition}[A bound for the actual signed $h=1$ slice]\label{sw:budget}
Uniformly in the original cutoff and this actual block,
\begin{equation}\label{sw:bound}
 |S_1^{XY}|\ll_{\nu,S,\epsilon}\|W\|_\infty^2(HD)^\epsilon
 \min\left\{R+\frac{DR}{M}+\frac{D^{2/3}R}{M^{1/3}},\ D+R^2\right\}.
\end{equation}
The exact identity underlying the estimate is
\begin{equation}\label{sw:mobius}
 \begin{aligned}
 DS_1^{XY}&=\sum_{\substack{d\ \mathrm{squarefree}\\(d,S)=1}}\mu(d)
 \sum_{\substack{u\ \mathrm{in\ the\ actual\ block}\\(d,u)=1}}
 T_d^X(u)\overline{T_d^Y(u)},\\
 T_d^X(u)&=\sum_{\substack{a\ \mathrm{squarefree}\\(a,dS)=1}}
 X(da)\nu(a)\chi_a(u)W(Nd\,Na/D).
 \end{aligned}
\end{equation}
In particular $X(da)$ is not replaced by $X(d)X(a)$, and the mask
$(d,u)=1$ remains.
\end{proposition}
The proof applies the accepted block sieves inside each positive row norm
in~\eqref{sw:mobius}, then explicitly loses the $d$ signs in an absolute
triangle. At $H=D^{3/4}$ the two remaining power-scale neighborhoods
$A_1\asymp D^{5/12},A_2\asymp D^{1/6}$ and
$A_1\asymp D^{3/4},A_2\asymp1$, with feasible capacity constants, both
retain power $7/6$. The desired $3/4$ still needs a signed saving of
$5/12$. This proves a separate signed-slice bound; it does not improve the
published global envelope or supply the missing cancellation.

There is also an exact arithmetic description of the two-pivot phase.
Put $\delta=1+2\omega$ and, at each good prime ideal $p$, set
$j(p)=J_p(\chi_p,\chi_p^2)$. Extend $j$ multiplicatively on squarefree
columns. Let $\Psi(a)$ be the unique generator of an ideal prime to
$2\delta$ satisfying $\Psi(a)\equiv1\pmod{2\delta}$.
\begin{proposition}[The actual Jacobi column twist]\label{sw:twist}
The map $\Psi$ is an algebraic Hecke character with exact finite conductor
$(2\delta)$, of norm $12$, and principal convention
$\Psi((\alpha))=\alpha$ for $\alpha\equiv1\pmod{2\delta}$
(infinity type $(1,0)$). For good squarefree columns,
\begin{equation}\label{sw:character}
 (j(p))=p,\quad -j(p)=\Psi(p),\qquad
 \mu(m)j(m)=\Psi(m).
\end{equation}
Define $\vartheta(m)=\Psi(m)/\sqrt{Nm}$ and
$\rho(m)=j(m)/\sqrt{Nm}=\mu(m)\vartheta(m)$ on those columns.
In the high window the actual gauged coefficients are
$b_A=(r_U-1)\vartheta\nu W$, $b_C=r_U\vartheta\nu W$ and
$b_\mu=\vartheta\nu W$.
\end{proposition}
This does not replace the original row kernel. Its compensating inverse
gauge remains, so all ordinary coefficient/operator norms are unchanged.
The new character has nonzero infinity type and is not silently covered
by a theorem stated only for finite-order twists. The identities in
\eqref{sw:character} involving $\mu$ must not be extended to arbitrary
nonsquarefree sixth factors.

A fresh fixed-profile example takes $\nu=1$, $S$ the primes above $2,3$,
literal unit one and primary row generators $1$ modulo $3$. It lies in an
actual $A_6=1$ block at $D=60000$, $U=244$, $H=3833$, $A_1=37$, $A_2=7$.
All $20$ rows of the stated literal class are nonprincipal. A selected
$q$ has exactly two legal column orientations, of norms $97051/96091$
and its reversal. Their correlations in the $\omega$ basis are
$(-4,-2)$ and $(-2,2)$, with mean $-3$. Their selected contribution is
$S_{1,q}^{AA}=-1/10000$; deleting the mean gives zero.
The centered squared norm is $6$, not zero, and the proposed
variance-only Cauchy bound would incorrectly require $36\le12$.
This refutes that centering rule in a fixed-profile row block, not the
asymptotic target. Other $q$ contributions remain.
The \href{../profile-uniformity.html#fixed-h1}{coprime-kernel inspector}
shows the actual rows, raw zeros, two-column Gram matrix and compensated
gauge; no training or arithmetic spectrum is inferred from that finite view.

% Full scoped proof; original signed estimate and determinant bridge remain open.
\section{Proof of the signed-slice budget and its surviving analytic burden}\label{sw:details}
\subsection{Gcd inversion before an absolute bound}
Insert $\mathbf1_{(m,n)=1}=\sum_{d\mid m,\ d\mid n}\mu(d)$ in
\eqref{sw:slice}, and write $m=da,n=db$. Squarefreeness forces
$(a,d)=(b,d)=1$. The common factors give
$\nu(d)\overline{\nu(d)}=1$ and
$\chi_d(u)\overline{\chi_d(u)}=\mathbf1_{(d,u)=1}$.
This proves~\eqref{sw:mobius}, retaining the original coefficient
$X(da)$ and the full row-zero mask. For $X=Y$ each inner squared row sum
is nonnegative; its $\mu(d)$-weighted sum need not be.

Put $L=D/Nd$. Whenever a column survives, $L\ge1/2$.
The actual divisor-assignment bound and ideal counting give
\[
 \sum_{\substack{a\ \mathrm{squarefree}\\(a,dS)=1}}
 |X(da)\nu(a)W(Nd\,Na/D)|^2
       \ll_\eta D^\eta\|W\|_\infty^2 L.
\]
Indeed $N(da)\le2D$, $|X(da)|\ll_\eta D^\eta$ uniformly in $d,U$,
and there are $O(L)$ ideals of norm at most $2L$ for $L\ge1/2$.
This includes $C_U$ by $C_U=A_U+\mu$ on the active annulus.
There is no multiplicativity assertion for $X(da)$.
For $L\ge1$, the powered block bound and the ordinary-conductor comparison
of Section~\ref{mcs:core}, followed by Cauchy for the cross term, imply
\begin{equation}\label{sw:divisor-budget}
 |G_d^{XY}|\ll (HD)^\eta\|W\|_\infty^2 L
 \min\left\{R+\frac{LR}{M}+\frac{L^{2/3}R}{M^{1/3}},\ L+R^2\right\},
\end{equation}
where $G_d^{XY}$ denotes the inner row sum in~\eqref{sw:mobius}.
The raw $(d,u)$ restriction is only enlarged inside these positive norms;
$u_0$ is a fixed column multiplier of modulus at most one.
For $1/2\le L<1$, only the unit column can survive. The direct row count
is $O(R)$, and the right side of~\eqref{sw:divisor-budget} is at least a
constant times $R$ since $R\ge1$ and $L\ge1/2$.
Thus this endpoint does not require a sieve with an unstated subunit length.

Now take the absolute triangle in $d$ explicitly. The two estimates may
be summed separately and their minimum taken. The elementary ideal sums
\[
 \sum_{Nd\le2D}\frac1{Nd}\ll\log(2D),\qquad
 \sum_d(Nd)^{-2}\ll1,\qquad \sum_d(Nd)^{-5/3}\ll1
\]
turn the first estimate into
$R\log(2D)+DR/M+D^{2/3}R/M^{1/3}$ after division by $D$;
the second gives $D+R^2\log(2D)$.
Absorb the logarithms with smaller internal exponents to obtain
\eqref{sw:bound}. All raw zeros, literal/ray/squarefree/capacity restrictions
were retained before the stated positive enlargements. No bound for an
indefinite submatrix was inferred from positivity of a larger Gram matrix.

At the mixed power-scale point the powered exponents are
$7/12,7/6,10/9$, and the ordinary exponents are $1,7/6$.
At the pure point they are $3/4,1,7/6$ versus $1,3/2$.
These are upper-bound exponents with feasible dyadic constants, not actual
energy lower bounds. A sufficient $h=1$ target for the A--A combination is
\begin{equation}\label{sw:target}
 \left|\sum_d\mu(d)G_d^{AA}\right|
       \ll_{\nu,S,W,\epsilon}DH(HD)^\epsilon.
\end{equation}
At $H=D^{3/4}$ this asks for power $7/4$ before division by $D$, while
\eqref{sw:bound} gives $13/6$, leaving $5/12$.
This is sufficient for this slice, not a necessary bound for every separate
C--C/C--$\mu$ term if their signed assembly cancels further.
The quantities $G_d$ depend on prime supports, $X(da)$ and row deletions;
they are not an independent smooth function of $Nd$ to which a scalar
M\"obius partial-sum estimate can simply be applied.

\subsection{The correctly oriented Jacobi generator}
At a good prime ideal $p$, let $n=Np$ and $\chi=\chi_p$ be the literal
positive sextic residue character. In $\mathcal O/p$, its value reduces
to $x^{(n-1)/6}$. For $j(p)=J(\chi,\chi^2)$, the elementary Gauss--Jacobi
identity and the three nontrivial Gauss norms give $|j(p)|^2=n$.
Reduction in that \emph{same} residue field gives
\[
 j(p)\equiv\sum_x x^{(n-1)/6}(1-x)^{(n-1)/3}=0\pmod p.
\]
Every exponent after expansion lies strictly between $0$ and $n-1$,
so the finite-field power sums vanish. Hence $(j(p))=p$.
This includes an inert prime of norm $p_0^2$, and inverting the literal
character conjugates the ideal at a split prime.

This raw sum equals the old $J(\chi^2,\chi^3)$: the quotient of the two
Gauss--Jacobi expressions is
$G(\chi)G(\chi^5)/G(\chi^3)^2=1$.
It differs from the theta-curve Frobenius normalization by its explicitly
stated factor $-\chi^3(-1)$; that sign is not imported silently.
No new Gauss or Jacobi norm theorem is claimed.

Let $f=(2\delta)$. The six units map bijectively onto
$(\mathcal O/f)^\times$: the components at $2$ and $\delta$ have orders
three and two, with complementary kernels. Since $\mathcal O$ is a PID,
every ideal prime to $f$ has a unique generator $\Psi(a)\equiv1\pmod f$.
Uniqueness gives multiplicativity. If $u_f(\alpha)$ is the unit having
the residue of $\alpha$ modulo $f$, then
\[
 \Psi((\alpha))=\alpha/u_f(\alpha),\qquad
 |\Psi(a)|^2=Na.
\]
This explicitly defines the finite part and the principal infinity-type
convention. The finite conductor divides $f$. A proper divisor omits a
component: the unit $-1$ is $1$ modulo $2$, and $\omega$ is $1$ modulo
$\delta$. A smaller defining modulus would require the unit ideal to
have value $-1$ or $\omega$. Thus the finite conductor is exactly $f$.

To identify the actual $j(p)$, now reduce in the \emph{character-value}
ring, a different operation from reduction modulo $p$ above.
Modulo $\delta$, a sixth root $v$ satisfies $v=v^3$ and $v^2=1$.
The zero branches at $x=0,1$ are kept, so
$j(p)\equiv\sum_{x\ne0,1}\chi^3(x)=-1\pmod\delta$.
Modulo $2$, $v=v^4$, and
$J(\chi,\chi^2)\equiv J(\chi^4,\chi^2)=-\chi^4(-1)=-1$.
The inverse-character Jacobi identity here follows by the bijection
$x/(1-x)$ onto the nonzero field elements other than $-1$.
Consequently $-j(p)\equiv1\pmod f$, proving $-j(p)=\Psi(p)$.
Multiplication proves~\eqref{sw:character} on squarefree columns.
For nonsquarefree ideals $\mu$ vanishes, so that last identity must not
be extended to arbitrary sixth factors. This construction establishes the
fixed character directly; it assumes no elliptic modularity, spectral
comparison, or distribution of prime arguments.

\subsection{The actual two-factor completion, not a new row kernel}
Freeze the literal $u_0,v_3,v_4,v_5,v_6$ and put
$Z=u_0v_3^3v_4^4v_5^5v_6^6$. The last factor is arbitrary, with its allowed
overlaps. For $h=1$, write $q=mn\ne1$, $Q=Nq$,
$\theta=\chi_m\overline{\chi_n}$ and $\xi_q=\theta^3$.
For $q=(\beta)$ use the original inverse-different pairing
$e_\beta(z)=\exp(2\pi i\operatorname{Tr}(z/(\delta\beta)))$.
Let $F(v_1,v_2)$ be the exact common finite weight, retaining the first-five
coprimality, primary/ray and capacity restrictions. It must be independent
of the column orientation. Define
\begin{align*}
 \widehat F_q(k_1,k_2)&=\sum_{v_1,v_2}F(v_1,v_2)
                       e_\beta(-k_1v_1-k_2v_2),\\
 T_{q,F}(\theta)&=\sum_{k_1,k_2\in(\mathcal O/q)^\times}
           \widehat F_q(k_1,k_2)\overline{\theta(k_1k_2^2)}.
\end{align*}
Periodization sums every actual representative; no injectivity assumption
is needed. The primitive powers $\theta,\theta^2$ have zero Gauss sums
at nonunit frequencies, including at inert primes. Exact inversion gives
\[
 \sum F\theta(v_1)\theta(v_2)^2
  =Q^{-2}\tau_\beta(\theta)\tau_\beta(\theta^2)T_{q,F}(\theta)
  =Q^{-2}J_q(\theta,\theta^2)\tau_\beta(\xi_q)T_{q,F}(\theta).
\]
Jacobi sums CRT without additive-generator multipliers, so
$J_q=j(m)\overline{j(n)}=\sqrt Q\,\rho(m)\overline{\rho(n)}$.
The quadratic central phase is the already recorded period-four law
\[
 \gamma(\beta)=\tau_\beta(\xi_q)/\sqrt Q
     =\tfrac12\{1+i^{-b}+i^a+i^{b-a}\},\qquad \beta=a+b\omega.
\]
For composite squarefree $q$, identifying $\tau_\beta(\xi_q)$ with
$\sum_xe_\beta(x^2)$ requires CRT expansion of the local root counts;
$1+\xi_q$ is not the root-count function on a composite ring.
The known quadratic Gaussian evaluation is reused, not retested as a new
phase theorem. Its pairing is exactly the displayed $e_\beta$.

With $b_X=\rho X\nu W$, the actual two-pivot identity is
\begin{equation}\label{sw:dual}
 DS_1^{XY}=\sum_q\frac{\gamma(\beta_q)}{Q}
   \sum_{mn=q} b_X(m)\overline{b_Y(n)}\theta_{m,n}(Z)
                        T_{q,F}(\theta_{m,n}).
\end{equation}
The two Gauss magnitudes, equivalently the central Gauss and Jacobi
magnitudes, give exactly $1/Q$.
There is no extra $\xi_q(-1)$ in this two-factor identity; the full
six-factor formula of Section~\ref{jdq:core} retains that sign.
If $\beta$ changes to $u\beta$, both $\gamma$ and $T$ acquire the same
quadratic sign $\xi_q(u)$, which cancels in their product.
If $Z$ shares any prime with $q$, $\theta(Z)=0$, including at a
sixth-power collision. No dropped primitive-conductor prime is restored.

Multiplying columns by $\rho$ changes the kernel by its compensating
inverse unitary gauge. It preserves operator norms and eigenvalues even
for an indefinite coprime-truncated matrix. A fixed-character coefficient
estimate cannot replace a bound for the actual $q$-dependent $T_{q,F}$.
At the pure endpoint $v_2=c$ with $(c,q)=1$, exact summation gives
\[
 T_{q,F}(\theta)=\theta(c)^2\tau_\beta(\overline{\theta^2})
       \sum_{k_1\in(\mathcal O/q)^\times}\widehat F_1(k_1)
                                      \overline{\theta(k_1)}.
\]
Since $\theta^2(-1)=1$, the extra Gauss pair multiplies to $Q$,
returning the original one-pivot scalar $\tau_\beta(\theta)/Q$.
A nonunit $c$ gives the original zero. Thus a putative central-phase gain
cannot be retained after discarding a literal length-one Fourier factor.

\subsection{A central reciprocal remains after the gauge}
On squarefree columns put $B_U=\mu A_U$ and
$\vartheta=\Psi/\sqrt N$. The original ternary assignment gives exactly
\[
 B_U(m)=\sum_{abc=m\atop Na,Nb\le U}\mu(c),\qquad
 \rho(m)A_U(m)=\vartheta(m)B_U(m).
\]
All factors are squarefree and pairwise coprime. Fix a nonzero numerator
$z$ and keep the full deleted set $T=S\cup\operatorname{supp}(z)$.
Write $\Xi$ for the unitary primitive character of
$\vartheta\nu\kappa_z$, with these raw deletions retained.
Its nonzero infinity type prevents a silent reuse of a finite-order-only
moment theorem. For $\Re s>1$,
\begin{align*}
 J_U^{\sharp,T}(s)&=\sum_{m\ \mathrm{squarefree},\ (m,T)=1}
                   B_U(m)\Xi(m)(Nm)^{-s}=N_U^T(s)/L^T(s,\Xi),\\
 N_U^T(s)&=\sum_{\substack{a,b\ \mathrm{squarefree}\\(a,b)=1,\ (ab,T)=1\\Na,Nb\le U}}
 \frac{\Xi(ab)}{(Nab)^s}\prod_{p\mid ab}(1-\Xi(p)(Np)^{-s})^{-1}.
\end{align*}
Fixing $a,b$ leaves the M\"obius Euler product in $c$; this proves the
identity by absolute convergence. The numerator is finite and analytic
for $\Re s>0$. At $U=1$ it is one, so the series is exactly $1/L^T(s,\Xi)$.
That boundary identity is not a high-window counterexample. It shows why
cancellation of central reciprocal poles cannot simply be inferred from
the column-character identification.

For $x_p=\Xi(p)(Np)^{-s}$, $\alpha_p=(Np)^{-v}$ and
$\beta_p=(Np)^{-w}$, the gauged three-placement Euler factor is
$1+x_p(\alpha_p+\beta_p-1)$, with the opposite signs from the ungauged
factor in Section~\ref{sb:core}. In the initial absolute domain it gives
\[
 \prod_{p\notin T}\{1+x_p(\alpha_p+\beta_p-1)\}
 =\frac{L^T(s+v,\Xi)L^T(s+w,\Xi)}{L^T(s,\Xi)}\prod_{p\notin T}H_p,
\]
where
\[
 H_p=\frac{\{1+x_p(\alpha_p+\beta_p-1)\}
                 (1-x_p\alpha_p)(1-x_p\beta_p)}{1-x_p}.
\]
Expansion gives
\begin{align*}
 H_p-1&=\frac{x_p^2(\alpha_p+\beta_p-\alpha_p^2-\beta_p^2-\alpha_p\beta_p)}{1-x_p}\\
 &\quad+\frac{x_p^3\alpha_p\beta_p(\alpha_p+\beta_p-1)}{1-x_p}.
\end{align*}
Thus the correction converges normally in the open region $\Re s>0$
where the real parts of all of
\[
 \begin{gathered}
 2s+v,\ 2s+w,\ 2s+2v,\ 2s+2w,\ 2s+v+w,\\
 3s+2v+w,\ 3s+v+2w,\ 3s+v+w
 \end{gathered}
\]
exceed one. In particular, for $\Re s\ge1/2+\eta$ and
$\Re v,\Re w\ge0$, its terms differ from one by
$O_\eta((Np)^{-1-2\eta})$. Deleted primes only remove factors.
Finite correction factors may vanish; no zero-free product is asserted.
The displayed factorization is first an absolute-domain identity and is
used meromorphically only where the functions have been continued.
At $v=w=0$ it simplifies to $L^T(s,\Xi)/L^T(2s,\Xi^2)$, a special
slice cancellation which does not remove the general central reciprocal.
No finite-height inversion is declared exact without truncation errors;
no new contour, reciprocal family moment or analytic saving is proved.
Most importantly, this fixed-$z$ coefficient series is not the actual
$q$-dependent joint weight in~\eqref{sw:dual}.

\subsection{The fixed-profile remaining-sector diagnostic}
Use the fixed smoothstep $s(t)=e^{-1/t}/(e^{-1/t}+e^{-1/(1-t)})$ for
$0<t<1$, extended by zero and one at the two sides, and set
$W(y)=s(8(y-1))s(8(2-y))$. It is smooth, supported in $[1,2]$, and
identically one on $[9/8,15/8]$, independently of $D$.
Take $\nu=1$ and $S$ the primes above $2,3$.
At $D=60000,U=244,H=3833$, select the actual block $A_1=37,A_2=7$,
literal unit one and primary row generators $1$ modulo $3$. Complete enumeration gives $44$ candidate pairs,
with two coprimality and $22$ capacity rejections, leaving $20$ distinct
nonprincipal rows. The first outside-$S$ valuation is one in each row.

Choose $\mathfrak p_p=(p,\omega-\lambda_p)$ for
\[
 (p,\lambda_p)\in\{(37,10),(43,6),(61,13),(307,17),(313,98)\}.
\]
Their product $q$ has exactly two legal annular column orientations:
$97051/96091$ and its reversal. Direct small- and long-divisor assignments
give $r_U=(0,2)$, $\mu=(-1,1)$, $C_U=(0,2)$ and $A_U=(1,1)$.
Both profile values are exactly one. Twelve orientation/row evaluations
are raw zeros. Exact residue-field exponentiation yields
\[
 \Gamma_1=-4-2\omega,\qquad \Gamma_2=-2+2\omega,
 \qquad \Gamma_{\rm av}=-3.
\]
Here $\Gamma_{\rm av}$ is the orientation average.
With $a^{AA}=(-1,-1)$ and $\mu(q)=-1$,
\[
 S_{1,q}^{AA}=-6/60000=-1/10000,\qquad
 \sum_r a_r(\Gamma_r-\Gamma_{\rm av})=0.
\]
The centered norm is
$\sum_r|\Gamma_r-\Gamma_{\rm av}|^2=6$ and $\sum_r|a_r|^2=2$,
whereas $|\sum_ra_r\Gamma_r|^2=36>12$.
A variance-only Cauchy bound therefore fails even with nonzero variance.
It is not a negative full row energy or an asymptotic counterexample.

The corresponding \emph{selected two-column} Gram matrix has diagonal
$14,20$ and the displayed conjugate off-diagonal entries. Its eigenvalues
are $17\pm\sqrt{21}$ and its $A$ quadratic form is $28$ before division
by $D$. Removing the diagonal for coprimality gives eigenvalues
$\pm2\sqrt3$ and the $A$ form $-6$. Other original columns and $q$ values
have not been suppressed in a full-family claim: this display is a
selected-$q$ contribution with its complete stated row block.
All finite data and scalar identities have independent exact checks;
no finite spectrum or learned label implies a determinant or GRH theorem.
The next estimate must retain the signed mean, reciprocal burden and
actual coupled $q$ weights, closing the still-open $5/12$ gap.

% Actual signed divisor variable; no new asymptotic saving or Lean declaration.
\section{The signed divisor mean and a complete-column laboratory}\label{sd:core}
The preceding signed-slice estimate loses the sign of the common ideal
 divisor. We now retain that variable in an exact analytic reduction and
in a complete finite computation. Neither supplies the missing
$5/12$ saving at $H=D^{3/4}$. The computations below use a new fixed
profile; they do not fill in the uncomputed columns of the earlier
broad-plateau profile.

Let $\mathcal C$ be original squarefree good ideal columns in $[D,2D]$,
and let $\mathcal R$ be an actual finite literal row block, with
nonnegative multiplicities $w_u$. All shared-prime zeros of the original
$\chi_m(u)$ remain. For arbitrary complex coefficient vectors define the
linear-first form
\begin{equation}\label{sd:form}
 DQ_1(x,y)=\sum_{u\in\mathcal R}w_u
 \sum_{\substack{m,n\in\mathcal C\\(m,n)=1}}
 x_m\overline{y_n}\chi_m(u)\overline{\chi_n(u)}.
\end{equation}
The original component is obtained with
$x_m=X_U(m)\nu(m)W(Nm/D)$ and the corresponding $y_m$.
On the nonprincipal row blocks used below the original residue term
vanishes. Equation~\eqref{sd:form} is a raw coprime form on a generic
block; principal rows would still require their separate residue assembly.

For every ideal divisor of an active column put
\[
 F_d(x;u)=\sum_{m\in\mathcal C,\ d\mid m}x_m\chi_m(u),\qquad
 Z_d(x,y)=\sum_u w_uF_d(x;u)\overline{F_d(y;u)}.
\]
The exact identity is
\begin{equation}\label{sd:incidence}
 DQ_1(x,y)=\sum_d\mu(d)Z_d(x,y).
\end{equation}
There is no extra $\mu(q)$: the vectors already contain the original
coefficients. Full divisors $d=m$ are essential to cancel self-pairs.
Divisors and squarefreeness refer to ideals, including distinct ideals
with the same rational norm.

\begin{proposition}[A weighted divisor-incidence lift]\label{sd:lift}
Suppose $\lambda_d\ge0$ is an arbitrary-coefficient bound
\[
 \sum_u w_u|F_d(x;u)|^2\le\lambda_d
                 \sum_{m\in\mathcal C,\ d\mid m}|x_m|^2.
\]
With $B(m)=\sum_{d\mid m}\lambda_d$ and
$\|x\|_B^2=\sum_mB(m)|x_m|^2$, one has
\begin{equation}\label{sd:weighted}
 |DQ_1(x,y)|\le\|x\|_B\|y\|_B.
\end{equation}
For a finite literal phase table one may take the exactly computed
\[
 \lambda_d=E_d:=\sum_u w_u\#\{m\in\mathcal C:d\mid m,\ \chi_m(u)\ne0\}.
\]
This last choice is a Frobenius certificate, not a new analytic estimate.
For $x=t+e$ it gives the finite perturbation bound
\begin{equation}\label{sd:certificate}
 |Q_1(x,x)-Q_1(t,t)|\le
 D^{-1}\bigl(2\|t\|_B\|e\|_B+\|e\|_B^2\bigr).
\end{equation}
\end{proposition}
\begin{proof}
Insert ideal-divisor inversion in~\eqref{sd:form}. Apply row Cauchy to
each $Z_d$ and weighted Cauchy to the resulting sum over $d$.
Then $\sum_d\lambda_d\|x_d\|_2^2=\|x\|_B^2$.
The squared Frobenius norm of the actual $F_d$ matrix is exactly $E_d$.
Expand $Q_1(t+e,t+e)-Q_1(t,t)$ to obtain the last assertion.
Positivity of the coprime-truncated matrix was never assumed.
\end{proof}
Using an accepted positive-row sieve bound for $\lambda_d$ only recovers
the old exponent up to ideal-divisor factors. It does not estimate the
signed sum before the triangle in~\eqref{sd:incidence}.

There is also a genuinely signed scalar expansion. Fix squarefree residual
ideals $a,b$, which may overlap, and an actual row $u$.
Keep the full deleted set $\mathcal T=S\cup\operatorname{supp}(abu)$ and set
\[
 \mathcal Z_{U;a,b,\mathcal T}(s)=
 \sum_{\substack{d\ \mathrm{squarefree}\\(d,\mathcal T)=1}}
 \frac{\mu(d)A_U(da)A_U(db)}{(Nd)^s},\qquad\Re s>1.
\]
For four positive cutoff variables put
$P_p=1-(Np)^{-v_1}-(Np)^{-v_2}$ and
$Q_p=1-(Np)^{-w_1}-(Np)^{-w_2}$.
The common-divisor Euler factor is $1-(Np)^{-s}P_pQ_p$.
Its first-order factorization has four numerator and five denominator
zeta factors:
\begin{equation}\label{sd:nine}
 \prod_{p\notin\mathcal T}(1-(Np)^{-s}P_pQ_p)
 =\frac{\prod_{i=1}^2\zeta_K^{\mathcal T}(s+v_i)
                         \zeta_K^{\mathcal T}(s+w_i)}
 {\zeta_K^{\mathcal T}(s)
       \prod_{i,j=1}^2\zeta_K^{\mathcal T}(s+v_i+w_j)}\,
       \mathcal H_{\mathcal T}(s;v,w).
\end{equation}
The correction is normally convergent at
$\Re s\ge\tfrac12+\eta$, $\Re v_i,\Re w_j\ge0$, for each fixed
$\eta>0$. This says nothing about zeros or residues of the zeta ratio.
Finite-height Perron kernels, with their explicit endpoint/error costs,
recover the cutoffs by successive limits.

A direct finite cutoff expansion gives more explicit information:
\begin{equation}\label{sd:finite-numerator}
 \mathcal Z_{U;a,b,\mathcal T}(s)
       =\frac{\mathcal N_{U;a,b,\mathcal T}(s)}
                    {\zeta_K^{\mathcal T}(s)}.
\end{equation}
The numerator is a finite sum over eight nonzero joint short-divisor
states, holomorphic on $\Re s>0$. Uniformly for $\sigma\ge\sigma_0>0$,
\[
 |\mathcal N_{U;a,b,\mathcal T}(\sigma+i\tau)|
 \ll_{K,\sigma_0,\epsilon}(Na\,Nb)^\epsilon
             U^{4\max(1-\sigma,0)+\epsilon}.
\]
The costs near $\sigma=1/2$ persist. The remaining long state has the
M\"obius sign, but the eight finite short states cannot be omitted.
A fixed-residual reciprocal zeta factor also does not prove an
uncancelled pole in the complete coherent sum.

\paragraph{Exact complete-column reference.}
Choose once $\delta=1/240000$ and
$W_\delta(y)=s((y-1)/\delta)s((2-y)/\delta)$, where $s$ is the smooth
step from zero to one. For the finite grid $D\le60000$ every strict
integer norm node has $W_\delta=1$; endpoints vanish. Its derivatives
cost $O_k(\delta^{-k})$. The three cases have different declared row
blocks, not a common asymptotic scaling experiment. They use principal
$\nu$, $S$ above $2,3$, literal class one and nonprincipal actual rows.
\begin{center}
\begin{tabular}{r|r|r|r|r}
 $D$ & columns & divisors & rows & $S_1^{AA}$\\\hline
 $144$ & $41$ & $53$ & $5$ & $-19/36$\\
 $576$ & $164$ & $212$ & $11$ & $-221/288$\\
 $60000$ & $16953$ & $21750$ & $20$ & $-5119/5000$
\end{tabular}
\end{center}
In the last case the exact even/odd M\"obius masses are
$543416$ and $604844$, giving numerator $-61428$ before division by $D$.
All $339060$ raw column--row values, including $47908$ zeros, and all nine
cross components are retained. In particular $1361$ columns have
nonsquarefree rational norm while remaining squarefree as ideals.
The \href{https://amelie-iska.github.io/GRH-research/signed-divisor.html}{signed-divisor inspector}
exposes the actual states and dimensions. These finite negative slices
are neither whole-family asymptotic counterexamples nor a positive
arithmetic determinant or GRH result.

% Full proofs for actual divisor incidence, signed common-d series and finite laboratory.
\section{Proofs and limits of the signed-divisor development}\label{sd:details}
\subsection{Actual divisor incidence and a finite error certificate}
The identity
$\mathbf1_{(m,n)=1}=\sum_{d\mid m,\ d\mid n}\mu(d)$ is an identity of
ideal divisors. Expanding the finite sums proves~\eqref{sd:incidence}.
When $m=da$ is squarefree, $(d,a)=1$. If $(d,u)=1$, multiplication by
$\chi_d(u)$ is a rowwise unit phase; if $(d,u)>1$ it is the original zero.
For original $x_m=X_U(m)\nu(m)W(Nm/D)$, precisely
$F_d(x;u)=\nu(d)\chi_d(u)T_d^X(u)$ on $(d,u)=1$, and it is zero
otherwise. Thus their row norms agree on the actual restricted support.
For general vectors the coefficient $x_{da}$ remains unfactored.
No identity $X_U(da)=X_U(d)X_U(a)$ is used. Residual $a,b$ may overlap.

Write $\|x_d\|_2^2=\sum_{d\mid m}|x_m|^2$.
The hypotheses of Proposition~\ref{sd:lift} give
\[
 \begin{split}
 |DQ_1(x,y)|&\le\sum_d |Z_d(x,y)|
 \le\sum_d\lambda_d\|x_d\|_2\|y_d\|_2\\
 &\le\left(\sum_d\lambda_d\|x_d\|_2^2\right)^{1/2}
       \left(\sum_d\lambda_d\|y_d\|_2^2\right)^{1/2}
 =\|x\|_B\|y\|_B.
 \end{split}
\]
This is a divisor-incidence lift of positive-row bounds to an indefinite
coprime form. It is not an inference from positivity of a larger Gram
matrix. In particular $B(m)\le\lambda_{\max}\tau_{\rm ideal}(m)$,
with $\tau_{\rm ideal}(m)=2^{\omega(m)}$ on squarefree ideals.
The subpolynomial divisor bound recovers the existing exponent; this
triangle does not supply the desired signed cancellation.

For the original unweighted two-pivot block ($w_u=1$), the accepted
arbitrary-coefficient inputs
at $L=D/Nd\ge1$ allow
\[
 \lambda_d\ll_\epsilon(HD)^\epsilon
 \min\left\{R+\frac{LR}{M}+\frac{L^{2/3}R}{M^{1/3}},\ L+R^2\right\},
 \qquad R=A_1A_2,\quad M=\max(A_1,A_2).
\]
Their constants are classical-input constants, not computed numbers.
Any removal of a $d$-row mask occurs only inside the positive norm.
For $1/2\le L<1$, $Na\in[L,2L]$ permits at most the unit ideal; its
operator norm is bounded directly by the actual row measure.
For an empty subsequence take $\lambda_d=0$.

For the fully computable choice $\lambda_d=E_d$, the matrix entries are
$\sqrt{w_u}\chi_m(u)$ on $d$-divisible columns. The squared Frobenius
norm is precisely the stated raw nonzero-entry count. It is valid for
arbitrary coefficient vectors, not just the saved trajectory. No unknown
sieve constant enters that finite certificate.
For approximations $t_x,t_y$ with errors $e_x=x-t_x,e_y=y-t_y$,
\[
 |Q_1(x,y)-Q_1(t_x,t_y)|
 \le D^{-1}\bigl(\|e_x\|_B\|y\|_B+
                    \|t_x\|_B\|e_y\|_B\bigr).
\]
The AA specialization follows by expanding three terms. To verify it
without floating square roots, put $T=\|t\|_B^2,E=\|e\|_B^2$ and
$L=D|Q_1(x,x)-Q_1(t,t)|-E$.
If $L>0$, the certificate is the equivalent real inequality
$L^2\le4TE$; if $L\le0$ it is immediate. In the saved exact
$\mathbf Z[\omega]$ computations with rational weights, these are
integer/rational checks rather than floating-root tests.
These are finite coefficient-error statements. They do not prove a
varying-family remainder estimate, spectral positivity or a determinant
comparison.

\subsection{A signed common-divisor Euler product}
From here restore ALL original annular squarefree good columns, not an
arbitrary subset from the vector lemma, and use the actual unweighted
raw row block of Section~\ref{sw:core}.
On the stated nonprincipal blocks it is the corresponding component of
the residue-subtracted moment. On a block containing principal rows the
raw identity remains valid, but the original principal correction is
additional and is not discarded.
Expand the exact $d$ identity while keeping its sign. For fixed $a,b,u$
the inner sum is
\begin{equation}\label{sd:inner}
 \sum_{\substack{d\ \mathrm{squarefree}\\(d,Sabu)=1}}
 \mu(d)A_U(da)A_U(db)
       W(Nd\,Na/D)\overline{W(Nd\,Nb/D)}.
\end{equation}
The exact outer coefficient is
$\nu(a)\overline{\nu(b)}\chi_a(u)\overline{\chi_b(u)}$.
The residual $a,b$ are individually squarefree but need not be coprime.
If their quotient character loses a common prime, that prime remains in
$\mathcal T=S\cup\operatorname{supp}(abu)$ and in every original zero.
The common $d$ character cancels only to $\mathbf1_{(d,u)=1}$.

For $\Re s>1$ the series $\mathcal Z$ is absolutely convergent, with
coefficient majorant
$3^{\omega(a)+\omega(b)}9^{\omega(d)}$, uniformly in $U$.
Use four positive Perron variables. Write $\mathcal E_{\mathcal T}(s;v,w)
=\prod_{p\notin\mathcal T}(1-(Np)^{-s}P_pQ_p)$.
The fixed residual prefactor is
$P_a(v)Q_b(w)=\prod_{p\mid a}P_p(v)\prod_{p\mid b}Q_p(w)$;
a common prime of $a,b$ contributes both factors.
The common $d$ prime has nine joint cutoff positions and the sign
$\mu(p)=-1$, giving local factor $1-zPQ$ with
$z=(Np)^{-s}$, $\alpha_i=(Np)^{-v_i}$ and $\beta_j=(Np)^{-w_j}$.
Thus the exact correction in~\eqref{sd:nine} is
\begin{equation}\label{sd:local-correction}
 H_p=\frac{(1-zPQ)\prod_{i=1}^2(1-z\alpha_i)
                           \prod_{j=1}^2(1-z\beta_j)}
                {(1-z)\prod_{i,j=1}^2(1-z\alpha_i\beta_j)}.
\end{equation}
The linear coefficient of numerator and denominator is
$-1-\sum_{i,j}\alpha_i\beta_j$, fixing all nine zeta signs.
At all four auxiliary variables zero the product is
$1/\zeta_K^{\mathcal T}(s)$ and $H_p=1$.
At $v_2=w_2=0$ it is $1/\zeta_K^{\mathcal T}(s+v_1+w_1)$, again
with $H_p=1$. These are generating-function checks, not finite-cutoff
identifications.

For $\Re s\ge1/2+\eta$, $\Re v_i,\Re w_j\ge0$, the degree-five
numerator and denominator of~\eqref{sd:local-correction} have equal
constant and linear coefficients. The sums of absolute coefficient
majorants in degrees two through five are at most $146$ and $26$.
Since $|PQ|\le9$ and $|\alpha_i|,|\beta_j|\le1$,
\begin{equation}\label{sd:normal}
 |H_p-1|\le C_\eta(Np)^{-1-2\eta},\qquad
 C_\eta=\frac{172}{(1-7^{-1/2-\eta})^5}.
\end{equation}
Indeed the five denominator factors have modulus at least $1-|z|$;
good prime ideal norms are at least seven. Ideal counting proves normal
convergence, uniform in deleted primes and vertical phases, and
\[
 |\mathcal H_{\mathcal T}|
 \le\exp\left(C_\eta\sum_p(Np)^{-1-2\eta}\right).
\]
Local correction factors can vanish; no inverse-correction bound is
asserted. Normal convergence controls neither the four numerator nor
five reciprocal zeta factors in the critical region. Their poles, zeros
and crossed residues cannot be suppressed.

\subsection{Finite Perron heights and the exact eight-state numerator}
Set $V=\lfloor U\rfloor+1/2$ and
$\Delta_V=\min_{n\ge1}|\log(V/n)|\ge1/(2V+1)$.
For four finite vertical heights $T_j$ on positive real lines $c_j$,
let $I_T(s)$ be the fourfold integral of
$P_aQ_b\mathcal E_{\mathcal T}(s;v,w)$ against
$\prod_j V^{z_j}\,dz_j/(2\pi i z_j)$.
In the initial half-plane $\sigma=\Re s>1$, finite-height Fubini is
justified by the absolute placement majorant. Put
\[
 \begin{split}
 A_{\rm abs}(\sigma,c)={}&
 \prod_{p\mid a}(1+q^{-c_1}+q^{-c_2})
 \prod_{p\mid b}(1+q^{-c_3}+q^{-c_4})\\
 &{}\cdot\prod_{p\notin\mathcal T}
 \{1+q^{-\sigma}(1+q^{-c_1}+q^{-c_2})
                     (1+q^{-c_3}+q^{-c_4})\},\quad q=Np.
 \end{split}
\]
Then
\begin{equation}\label{sd:perron-error}
 |I_T(s)-\mathcal Z_{U;a,b,\mathcal T}(s)|
 \le V^{\sum c_j}A_{\rm abs}(\sigma,c)
 \left\{\prod_{j=1}^4
    \left(1+\frac{2}{\pi T_j\Delta_V}\right)-1\right\},
\end{equation}
and $A_{\rm abs}\le3^{\omega(a)+\omega(b)}\zeta_K(\sigma)^9$.
For one short integer norm $n$, the indicator is at most $(V/n)^c$,
and the truncated Perron error is at most
$2(V/n)^c/(\pi T\Delta_V)$. Expand the difference of four products and
sum every nonempty error subset to prove~\eqref{sd:perron-error}.
Consequently successive Perron limits recover the exact cutoffs.
The half-integer gap also gives a height-uniform $C_U(V/n)^c$ bound,
allowing dominated convergence in the initial absolutely convergent
series. Four infinite Perron integrals are not declared absolutely
convergent or interchanged silently with Mellin inversion.

There is an entirely finite alternative. Partition
$a=a_0a_1a_2$ and $b=b_0b_1b_2$, internally coprime in each column.
Cross-column overlaps remain allowed. A common $d$ prime occupies one
state $(i,j)\in\{0,1,2\}^2$. Separate its $(0,0)$ part $e$ from eight
pairwise coprime squarefree atomic ideals $h_{ij}$ avoiding $\mathcal T$,
and let $h$ be their product. Define
\[
 h_i^L=\prod_{j=0}^2h_{ij}\ (i=1,2),\qquad
 h_j^R=\prod_{i=0}^2h_{ij}\ (j=1,2).
\]
An atomic $h_{ij}$ with both indices nonzero belongs to one short
factor in each column. These four short factors must not be forced
globally disjoint. Let $\mathcal A_U$ comprise precisely these tuples
satisfying
$N(a_i h_i^L)\le U$ and $N(b_jh_j^R)\le U$ for $i,j=1,2$.
The numerator in~\eqref{sd:finite-numerator} is exactly
\begin{equation}\label{sd:eight}
 \mathcal N_{U;a,b,\mathcal T}(s)=
 \sum_{\mathcal A_U}
 \frac{\mu(a_1a_2)\mu(b_1b_2)
       \mu(h_{11}h_{12}h_{21}h_{22})}{(Nh)^s}
       \prod_{p\mid h}(1-(Np)^{-s})^{-1}.
\end{equation}
At a common prime the sign is
$(-1)^{1+\mathbf1_{i>0}+\mathbf1_{j>0}}$.
It is negative for $(0,0)$, positive for the four single-short states,
and negative for the four double-short states. Summing the remaining
squarefree $e$, coprime to $\mathcal T h$, gives
$1/\zeta_K^{\mathcal T}(s)$ times the local product in~\eqref{sd:eight}.
This proves the exact identity, including its sign.

Multiplying the four cutoff inequalities yields $Nh\le U^4$; double
states are counted twice in that product. Hence the numerator is finite
and holomorphic on $\Re s>0$. At $U=1$ it equals one: this boundary is
not a high-window counterexample or a global pole obstruction.
For $\sigma\ge\sigma_0>0$ and every $\epsilon>0$,
\begin{equation}\label{sd:numerator-bound}
 |\mathcal N_{U;a,b,\mathcal T}(\sigma+i\tau)|
 \ll_{K,\sigma_0,\epsilon}(Na\,Nb)^\epsilon
               U^{4\max(1-\sigma,0)+\epsilon}.
\end{equation}
There are at most $3^{\omega(a)+\omega(b)}$ residual partitions and
$8^{\omega(h)}$ atomic placements at fixed $h$. The finite local product
is $\ll_{\sigma_0,\epsilon}(Nh)^\epsilon$. Apply the ideal-divisor bound
and ideal counting up to $U^4$, using smaller internal exponents; a
logarithm at the transition is absorbed in $U^\epsilon$.
Uniformity as $\sigma_0\to0$ is not claimed. For
$\sigma=1/2+\eta$ with $0<\eta<1/2$, the old-type cutoff cost is
$U^{2-4\eta+\epsilon}$, not a new saving.

\subsection{A uniform common-profile Mellin embedding of the actual mean}
For residual $a,b$ put $M=\max(Na,Nb)$ and
$\varrho=\min(Na,Nb)/M$. At $\varrho\le1/2$ the product of the two
profile values vanishes, including equality because of endpoint
vanishing. Otherwise put
\[
 \Phi_\varrho(t)=W(t)\overline{W(\varrho t)}
\]
when $Na\ge Nb$, and use its conjugate for the reversed orientation.
It is supported in $[1/\varrho,2]\subset[1,2]$.
With $M_k(W)=\max_{j\le k}\|(t\,d/dt)^jW\|_\infty$,
\begin{equation}\label{sd:profile-mellin}
 |\widehat\Phi_\varrho(\sigma+i\tau)|
 \le C_{I,k}M_k(W)^2(1+|\tau|)^{-k}
\end{equation}
uniformly for $\sigma$ in any fixed compact strip $I$ and all
$\varrho\in[1/2,1]$. Leibniz and log-coordinate integration by parts
prove this bound; no ratio-dependent change of $W$ is made.
The transform is entire. For $c>1$, ordinary Mellin inversion after the
finite numerator identity gives the exact scalar~\eqref{sd:inner} as
\begin{equation}\label{sd:mellin}
 \frac1{2\pi i}\int_{\Re s=c}
 \widehat\Phi_\varrho(s)(D/M)^s
 \frac{\mathcal N_{U;a,b,\mathcal T}(s)}
                      {\zeta_K^{\mathcal T}(s)}\,ds.
\end{equation}
This integral is vertically absolutely convergent in its initial
half-plane. The original profile forces $M\le2D$, so the residual $a,b$
and actual row sums are finite.

Keep the full local deletion factor
\[
 \frac1{\zeta_K^{\mathcal T}(s)}=
 \frac1{\zeta_K(s)}\prod_{p\in\mathcal T}(1-(Np)^{-s})^{-1}.
\]
At fixed $\sigma_0>0$ it costs
$\ll_{S,\sigma_0,\epsilon}(N\operatorname{rad}(abu))^\epsilon$,
hence at most $(HD^2)^\epsilon$ on actual support.
This estimates the finite mask dependence, not the central reciprocal.
Multiplying~\eqref{sd:mellin} by every original outer coefficient and
summing the actual rows and residuals gives
\begin{equation}\label{sd:coherent}
 DS_1^{AA}=\frac1{2\pi i}\int_{\Re s=c}
                   \frac{D^s\mathcal A_{D,H,U,W}(s)}{\zeta_K(s)}\,ds.
\end{equation}
Explicitly, with each profile orientation retained,
\[
 \begin{split}
 \mathcal A_{D,H,U,W}(s)
 ={}&\sum_{u\in\mathcal R}
 \sum_{\substack{a,b\ \mathrm{squarefree,good}\\Na,Nb\le2D,\ \varrho>1/2}}
 \nu(a)\overline{\nu(b)}\chi_a(u)\overline{\chi_b(u)}\,M^{-s}\,\widehat\Phi_\varrho(s)\\
 &{}\cdot\mathcal N_{U;a,b,\mathcal T}(s)
       \prod_{p\in\mathcal T}(1-(Np)^{-s})^{-1}.
 \end{split}
\]
This is a \emph{finite coherent sum}, not an independent marginal
moment or a positive sum. All raw zeros remain in its outer coefficients.
For fixed finite $D,H,U$ it is holomorphic on $\Re s>0$ and rapidly
decreasing on fixed vertical strips there, with constants depending on
those finite parameters. No uniform target estimate for it is supplied.
This is an exact signed embedding, not a zero-location theorem.

\subsection{A nonzero short-state correction and the unproved step}
An actual annular control uses good prime ideals of norms
$7,13,19,31,43$ with primary generators
$(1,3),(4,3),(-2,3),(1,6),(7,6)$ in the $1,\omega$ basis.
Take $D=1500,U=25$, $d=p_7$, $a=p_{13}p_{19}$,
$b=p_{13}p_{31}$ and the nonprincipal row $u=p_{43}$ in the pure block
$A_1=37,A_2=1,H=241$, literal class one.
The full raw supports are disjoint from $u$, and
$N(da)=1729,N(db)=2821$ lie in the annulus. A single fixed smooth
profile with plateau $[1.1,1.9]$ equals one at both nodes.
Direct assignments give
$A_U(a)=A_U(b)=-1$, $A_U(da)=1$, $A_U(db)=-1$.
The actual signed $d$ coefficient is $+1$. The $(0,0)$ state alone
would give $-1$; states $(1,0)$ and $(2,0)$ each contribute $+1$.
Dropping all short states therefore misses a nonzero $e=1$ correction.
This illustrates the new eight-state identity, not a new discovery of
coefficient nonmultiplicativity.

The original pair has gcd $p_7p_{13}$. Its contributions across
$d=1,p_7,p_{13},p_7p_{13}$ cancel exactly. Thus a fixed-residual
reciprocal does not prove an uncancelled pole or an obstruction in the
complete coherent mean. The missing step is to estimate that full mean
by $DH(HD)^\epsilon$, retaining fixed-profile/seminorm dependence,
raw masks and conductor/ray/capacity restrictions. Scalar Mertens
cancellation in the long state, a normal correction product, or an
absolute cutoff numerator bound does not prove it. Principal and short
assembly, independent positive determinant/zero-location bridges and
separate Dirichlet, Hecke, Dedekind, elliptic and Artin obligations remain
open. No holomorphy or desired zero location is assumed to close them.

\subsection{Exact arithmetic and source-bound computational scope}
The evaluator streams each actual row and accumulates original
$x_m\chi_m(u)$ into every ideal-divisor bucket $d\mid m$. It then
forms all nine cross products for $A_U,C_U,\mu$ and sums with $\mu(d)$.
Write $n_{\rm col}=|\mathcal C|$,
$I=\sum_m\tau_{\rm ideal}(m)$, $K$ for the distinct divisors and $r$
for the listed rows. This bucket stage costs $O(r(I+K))$ arithmetic operations with a
fixed three channels. Its streaming storage is $O(K)$ plus incidences
and either a streamed row or separately retained phase table.
Character construction, factorization and integer bit lengths are
additional costs, not declared free.
For coefficients bounded by $C$, bucket coordinates are $O(n_{\rm col}C)$,
products $O(n_{\rm col}^2C^2)$ and accumulated values
$O(r n_{\rm col}^2C^2)$ for unit row
weights. For compressed integer multiplicities replace $r$ in this
bit-size bound by $\sum_u w_u$ and account for weight multiplication.

For $D=60000$, the actual phase table has shape $20\times16953$;
the incidence has $73420$ nonzeros over $21750$ divisor buckets, and
$F_d$ has $20\times21750\times3$ complex states.
There are $271688516$ ordered coprime column pairs.
The naive all-pair row loop would visit $5748084180$ pairs and was
\emph{not executed} at this size. This is an operation count, not an
unexecuted timing benchmark. Independent direct pair evaluation covers
the two complete small families; the large family is checked linearly
from literal field arithmetic and divisor-poset inversion.

The exact large results also include
\[
 S_1^{CC}=\frac{661}{7500},\qquad
 S_1^{C\mu}=\frac{13931}{30000},\qquad
 S_1^{\mu\mu}=-\frac{229}{1250}.
\]
For the ordered vectors $(A_U,C_U,\mu)$ the full pairing table is
\[
 \frac1{30000}
 \begin{pmatrix}
 -30714&-11287&19427\\
 -11287&2644&13931\\
 19427&13931&-5496
 \end{pmatrix}.
\]
This is the table of three original coefficient vectors, not a spectral
carrier for GRH. The conjugate $\mu C$ entry and the AA identity retain
their actual linear-first convention. All three saved computations use source v1.
A separate v2 helper fixes its unused $N=U=1$ high-window assertion;
every one of the $17158$ executed columns has $N>U$, so no accepted
row/table or source snapshot was silently replaced.
The fixed-profile flat-node implementation refuses requests above its
stated $D\le60000$ certificate. Neither these three different finite
blocks nor their measured runtimes imply an asymptotic saving.

\paragraph{Classical reference}NIST Digital Library of Mathematical Functions, Section25.15, \href{https://dlmf.nist.gov/25.15#E5}{Dirichlet L-functions, equations 25.15.5--6}; primitive functional equation and finite Gauss sums (snapshot pinned 4 November 2026). The normalization and explicit finite Gauss sums used here are stated above.\end{document}
